The first three targets are the three-axis Weyl estimate, the flat Weyl and
Bach symbols, and the chiral curvature projection. The later sections record
the conditional neck argument, finite energy accounting and the open problem of
supplying the geometric hypotheses. Section references point to the working
\emph{Sections and Connections} manuscript. Section~\ref{sec:master-variational}
follows the master manuscript from a single positive-class competitor to
the marked-cut test and the physical energy scales.

\section{Three-axis Weyl estimate}

\begin{definition}[Null evaluations]
\label{def:null-evaluations}
\lean{SectionsAndConnections.axisNullVector, SectionsAndConnections.axisNullEvaluation, SectionsAndConnections.axisMagnitude}
\leanok
Let $V$ be a real Euclidean three-space with orthonormal basis
$n_1,n_2,n_3$. For cyclic $(i,j,k)$ set
$\zeta_i=(n_j+\mathrm{i}n_k)/\sqrt{2}$. For a real symmetric
trace-free endomorphism $W$ of $V$, let
$q_i=|\langle W\zeta_i,\zeta_i\rangle_{\mathbb C}|$.
The pairing is the complex bilinear extension of the real inner product.
The Lean declarations use coordinates in the specified orthonormal basis.
\end{definition}

\begin{lemma}[Three-axis Weyl estimate]
\label{lem:three-axis-weyl}
\uses{def:null-evaluations}
\lean{SectionsAndConnections.threeAxisWeylEstimate}
\leanok
Write $d_i$ for the diagonal entries of $W$ and $w_{ij}$ for its
off-diagonal entries in this basis. Then
\[
  \sum_i q_i^2
    =\tfrac34\sum_i d_i^2+\sum_{i<j}w_{ij}^2
    \ge \tfrac12\lVert W\rVert_{\mathrm{HS}}^2 .
\]
\end{lemma}

\begin{proof}
\leanok
For cyclic $(i,j,k)$ the complex pairing equals
$(d_j-d_k)/2+\mathrm{i}w_{jk}$. Summing the squared magnitudes gives
the equality because $\sum_i d_i=0$ implies
$\sum_{\mathrm{cyc}}(d_j-d_k)^2=3\sum_i d_i^2$.
The squared Hilbert--Schmidt norm is
$\sum_i d_i^2+2\sum_{i<j}w_{ij}^2$. This is the algebraic step used
for the section-energy lower bound in \S2.1.
\end{proof}

\section{Flat Weyl and Bach symbols}

\begin{definition}[Flat curvature symbols]
\label{def:flat-curvature-symbols}
In Euclidean four-space let $\delta$ denote the metric, $\xi$ a real
covector, and $H$ a real symmetric two-tensor. Identify covectors and
vectors using $\delta$. For symmetric two-tensors the Kulkarni--Nomizu
product has slots
\[
  (h\odot k)_{ij\ell m}
    =h_{i\ell}k_{jm}+h_{jm}k_{i\ell}-h_{im}k_{j\ell}-h_{j\ell}k_{im}.
\]
With the Fourier rule $\partial_j\mapsto\mathrm{i}\xi_j$, define
\[
\begin{aligned}
  R_\xi(H) &= \tfrac12(\xi\otimes\xi)\odot H,
  & (r_\xi(H))_{ij} &= \sum_k(R_\xi(H))_{kikj},\\
  s_\xi(H) &= \operatorname{tr} r_\xi(H),
  & A_\xi(H) &= \tfrac12\left(r_\xi(H)-\tfrac16 s_\xi(H)\delta\right),\\
  W_\xi(H) &= R_\xi(H)-A_\xi(H)\odot\delta,
  & (B_\xi(H))_{ij} &= -\sum_{k,l}\xi_k\xi_l (W_\xi(H))_{ikjl}.
\end{aligned}
\]
These are the coordinate symbols of the corresponding flat linearizations
in \S14.2, with the curvature convention giving positive sectional
curvature on the unit round sphere.
\end{definition}

\begin{lemma}[Weyl and Bach symbols on transverse traceless amplitudes]
\label{lem:flat-weyl-bach}
\uses{def:flat-curvature-symbols}
Suppose $\xi\ne 0$, $H\xi=0$, and $\operatorname{tr}H=0$. Then
\[
  B_\xi(H)=-\tfrac14|\xi|^4H,\qquad
  \lVert W_\xi(H)\rVert_{\mathrm{tens}}^2
    =\tfrac12|\xi|^4\lVert H\rVert_{\mathrm{HS}}^2 .
\]
The tensor norm on $W_\xi(H)$ sums over all four slots.
\end{lemma}

\begin{proof}
In an orthonormal basis with $\xi=|\xi|e^1$, all $H_{1i}$ vanish.
Thus $r_\xi(H)=|\xi|^2H/2$, $s_\xi(H)=0$, and
$A_\xi(H)=|\xi|^2H/4$. For transverse indices $a,b$, the curvature
component $(R_\xi(H))_{1a1b}=|\xi|^2H_{ab}/2$ gives
$(W_\xi(H))_{1a1b}=|\xi|^2H_{ab}/4$. The double contraction defining
$B_\xi$ gives its claimed symbol.

Counting the four ordered placements of the longitudinal index pair
gives $\lVert R_\xi(H)\rVert_{\mathrm{tens}}^2
=|\xi|^4\lVert H\rVert_{\mathrm{HS}}^2$.
For trace-free $A$ in dimension four,
$\lVert A\odot\delta\rVert_{\mathrm{tens}}^2=8\lVert A\rVert_{\mathrm{HS}}^2$.
The Weyl part pairs to zero with $A_\xi(H)\odot\delta$ by trace-freeness.
Subtracting its squared norm gives the Weyl symbol identity in \S14.2.
\end{proof}

\section{Chiral curvature projection}

\begin{definition}[Wedge projection of chiral curvature]
\label{def:chiral-wedge-projection}
Let $g$ be a smooth Riemannian metric on an oriented four-manifold and
$\Sigma_i$ a local frame of self-dual two-forms satisfying
$\Sigma_i\wedge\Sigma_j=2\delta_{ij}\,dV_g$.
Use the orientation in \S3.1 for this frame and the internal generators
$T_i$ with $[T_i,T_j]=\epsilon_{ijk}T_k$.
For a metric connection $A$ on $\Lambda^+$, express its curvature as
$F_A^i$ in those generators and define
\[
  C_{ij}\,dV_g=-\tfrac12 F_A^i\wedge\Sigma_j .
\]
\end{definition}

\begin{proposition}[Chiral curvature of the Levi--Civita connection]
\label{prop:chiral-curvature}
\uses{def:chiral-wedge-projection}
Take $A=\Gamma^+(g)$, the connection on $\Lambda^+$ induced by the
Levi--Civita connection, and write $s_g$ for scalar curvature. In the
orthonormal frame $\Sigma_i/\sqrt{2}$ the matrix $C$ satisfies
\[
  C=W^++\tfrac{s_g}{12}I .
\]
\end{proposition}

\begin{proof}
Identify each two-form $\alpha$ with the skew operator $K_\alpha$ by
$g(K_\alpha X,Y)=\alpha(X,Y)$ and take $T_i=K_{\Sigma_i}/2$.
The curvature operator $\mathcal R$ on two-forms uses the pairing
$g(R(X,Y)V,U)$ for the slots $(X,Y,U,V)$. In these conventions the
induced chiral curvature has $F_A^i=-\mathcal R\Sigma_i$.
This step uses the action of the self-dual summand on $\Lambda^+$ and
the trivial action of the anti-self-dual summand. The
$\Lambda^+\to\Lambda^+$ block of $\mathcal R$ is $W^++s_g I/12$.
Since $\Sigma_j$ is self-dual, wedging with $\Sigma_j$ retains this
block and the factor $2$ from normalization cancels the $1/2$ in the
definition of $C$. See \S3.1 for the curvature conventions and block
calculation.
\end{proof}

For the unit round $S^4$ the formula gives $C=I$. On the conformally
flat product cylinder $dt^2+g_{S^3}$ with unit round cross-section the
chiral curvature two-forms are tangential, so $F_A^i\wedge F_A^j=0$.
These models test the choice of sign and the behavior of the metric
formula on a collar with degenerate wedge pairing.

\section{Energy in a supplied annular chart}

\begin{definition}[Annular data]
\label{def:annular-data}
For $r>0$ and $R\ge2r$ set
$A=\{r\le |x|\le R\}$ and $A^*=\{r/2<|x|<2R\}$ in $\mathbb R^4$.
For a smooth Riemannian metric $g=\delta+v$ on $A^*$ define
\[
\begin{aligned}
  \varepsilon &= \sup_{A^*}|v|,
  & K &= \lVert D^2v\rVert_{L^2(A^*,dx)},\\
  \mathcal W_A(g) &= \int_A |W_g|_{\mathrm{tens}}^2\,dV_g,
  & \mathfrak B_{A^*}(g) &= \int_{A^*}|B_g|_g\,dV_g .
\end{aligned}
\]
The metric deviation and its derivatives use Euclidean norms. Both
integrals use $g$ for the norms and volume. The Weyl norm sums over
all four covariant slots.
\end{definition}

\begin{theorem}[Length-independent annular energy bound]
\label{thm:one-chart-energy}
\uses{def:annular-data}
There are universal constants $C$ and $\varepsilon_0\in(0,1/4]$ such
that in a supplied chart with $\varepsilon\le\varepsilon_0$ and finite
$K$ and $\mathfrak B_{A^*}(g)$,
\[
  \mathcal W_A(g)\le C\varepsilon
    \bigl[\mathfrak B_{A^*}(g)+(K+\varepsilon)(1+\varepsilon K)\bigr].
\]
The constants are uniform in $R/r$.
\end{theorem}

\begin{proof}
\uses{lem:buffered-euclidean-interpolation}
The proof in \S14.3 tests the Weyl first variation against $\eta^4v$,
where $\eta=1$ on $A$ and decays to zero on end shells in $A^*$.
The leading second-derivative terms of the Weyl variation in direction
$v$ agree with those of $W_g$. The remaining terms are estimated using
$\lVert\eta Dv\rVert_{L^4}^2\le C\varepsilon(K+\varepsilon)$.
Small $\varepsilon$ absorbs the term from varying the contractions of
the Weyl tensor. Young's inequality bounds the products of the square
root of the tested energy with the cutoff and interpolation errors.
Fixed cutoff profiles on the two shells give scale-invariant
$L^4$ bounds on $D\eta$ and $L^2$ bounds on $D^2\eta$ in dimension
four. The estimate therefore remains uniform as the annulus gets longer.
\end{proof}

\begin{definition}[Supplied control along neck scales]
\label{def:supplied-neck-control}
\uses{def:annular-data}
Suppose $r_i/R_i\to0$ at an interface with an annular identification.
For each fixed large $L$, assume that for all sufficiently large $i$
the middle region is covered by the chart of Definition~\ref{def:annular-data}
with radii $Lr_i$ and $R_i/L$, including the enlarged annulus.
Denote the corresponding data by $\varepsilon_i(L)$, $K_i(L)$ and
$\mathfrak B_i(L)$. Supplied control means bounds on $K_i(L)$ and
$\mathfrak B_i(L)$ that are independent of large $L$ and hold for all
sufficiently large $i$ at each such $L$, together with
\[
  \lim_{L\to\infty}\limsup_{i\to\infty}\varepsilon_i(L)=0 .
\]
\end{definition}

\begin{corollary}[No neck energy under supplied control]
\label{cor:supplied-no-neck}
\uses{def:supplied-neck-control}
Under the hypotheses of Definition~\ref{def:supplied-neck-control},
\[
  \lim_{L\to\infty}\limsup_{i\to\infty}
  \mathcal W_{\{Lr_i<|x|<R_i/L\}}(g_i)=0 .
\]
\end{corollary}

\begin{proof}
\uses{thm:one-chart-energy}
For all sufficiently large $L$, the chart deviations are eventually below
$\varepsilon_0$. The uniform bounds on $K_i(L)$ and $\mathfrak B_i(L)$
bound the bracket in Theorem~\ref{thm:one-chart-energy} by one constant.
Take the upper limit in $i$ for fixed $L$ and then let $L$ tend to infinity.
\end{proof}

\begin{openproblem}[Control of almost Bach-flat necks]
\label{problem:almost-bach-flat-control}
\uses{def:supplied-neck-control}
\notready
For smooth Riemannian metrics on a fixed oriented smooth homotopy
four-sphere, consider the positive-Yamabe band
$4\pi^2\le E_+\le8\pi^2-\gamma$
for fixed $\gamma>0$, bounded $L^2$ curvature and
$\int_M|B_{g_i}|_{g_i}\,dV_{g_i}\to0$. At a concentration interface
with separated scales, can one obtain the compatible charts and
bounds of Definition~\ref{def:supplied-neck-control}, or prove the
no-neck estimate by intrinsic control? Here
$E_+(g)=\int_M |W_g^+|_{\mathrm{HS}}^2\,dV_g$.
\end{openproblem}

The energy identity discussed in \S15 also needs the smooth base and ALE
comparison data and control of the full complement of the truncated
components. The local estimate in Corollary~\ref{cor:supplied-no-neck}
addresses the energy between a pair of concentration scales.

\section{Finite accounting of chiral Weyl energy}

We work on one common source sequence. The base and all bubbles contribute
in one chirality. All measures in the accounting theorems have finite
mass. For the geometric application the source measure is
$\mu_i=|W^+_{g_i}|_{\mathrm{HS}}^2\,dV_{g_i}$. A smooth compactified
base can carry positive energy after all bubble energies vanish.

\begin{lemma}[Finite accounting with a strict residual]
\label{lem:finite-bubble-accounting}
\lean{SectionsAndConnections.noPositiveBubbleOfFiniteAccounting}
\leanok
Put $p=4\pi^2$. Let $J$ be finite, $\delta>0$ and $p\le e\le2p-\delta$.
Suppose $e=b+\sum_{j\in J}x_j$, where $b,x_j\ge0$. Each of $b$ and
the $x_j$ is zero or strictly larger than $p$. If $e-x_j>0$ whenever
$x_j>0$, then every $x_j=0$ and $e=b$.
\end{lemma}
\begin{proof}
\leanok
For a positive $x_j$, the residual $e-x_j$ is less than $p$. It is
the sum of $b$ and all the other bubble energies. Each is nonnegative
and strictly below $p$, so each is zero. This contradicts $e-x_j>0$.
\end{proof}

\begin{lemma}[Source and component budgets in scalar form]
\label{lem:scalar-component-budgets}
\lean{FirstOrderTwistor.sourceBandScalarBudget, FirstOrderTwistor.sourceCurvatureBudget, FirstOrderTwistor.strictComponentResidual, FirstOrderTwistor.uniformComponentResidual}
\leanok
Write $Y_{\rm rd}=8\sqrt6\,\pi$. Given the source Euler identity
$Y^2=48(8\pi^2-E)+12D$, $0\le Y\le Y_{\rm rd}$, $D\ge0$,
and $E\le8\pi^2-\delta$ with $\delta>0$, one has $Y\ge\sqrt{48\delta}$
and $D\le4E$. If the full curvature integral is $8E+2D+Y^2/6$, it is
at most $192\pi^2-16\delta$. For a component with energy $x$, suppose
the limiting Euler identity holds at $(e,y,D)$ and $4x=d\le D$. Then
\[
 e-x\ge (Y_{\rm rd}^2-y^2)/48.
\]
If $0\le y\le\widehat Y<Y_{\rm rd}$ the right hand side is positive.
If $\widehat Y\le Y_{\rm rd}-d_\delta$ for a supplied
$0<d_\delta<Y_{\rm rd}$, the lower bound is at least
$[Y_{\rm rd}^2-(Y_{\rm rd}-d_\delta)^2]/48$.
\end{lemma}
\begin{proof}
\leanok
Substitute the Euler identity and $Y_{\rm rd}^2=384\pi^2$. On the
nonnegative real axis, squaring preserves the Yamabe comparisons.
\end{proof}

\begin{lemma}[Counting ends from the numerical end budget]
\label{lem:scalar-ale-end-count}
\lean{FirstOrderTwistor.oneALEEndOfSubcriticalBounds, FirstOrderTwistor.q8EtaContradictionOfSubcriticalBand}
\leanok
In the notation of \S11.3 let $k\ge1$ count the end orders $q_j\ge1$,
all bounded by $Q$. Suppose $e\le8\pi^2-\delta$, $\delta>0$,
$e_+,e_-\ge0$, and
\[
 (8\pi^2-e)+(e_++e_-)/2+4\pi^2(k-2+\textstyle\sum_j1/q_j)
 \le 8\pi^2/Q.
\]
Then $k=1$. If both chiral energies belong to $[0,8\pi^2-\delta]$,
the $Q_8$ signature defect $|e_+-e_-|=9\pi^2$ is impossible.
\end{lemma}
\begin{proof}
\leanok
If $k\ge2$, insert $\sum_j1/q_j\ge k/Q$ and subtract $8\pi^2/Q$.
The difference is at least $\delta+4\pi^2(k-2)(1+1/Q)>0$.
The two nonnegative chiral energies differ by less than $9\pi^2$.
\end{proof}

\begin{lemma}[Energy on a fixed comparison domain]
\label{lem:fixed-domain-density}
\lean{FirstOrderTwistor.energyMeasure_real_range_of_embedding, FirstOrderTwistor.tendstoImageEnergyOfUniformPullback, FirstOrderTwistor.chiralHSDensity_integrableOn_and_tendsto}
\leanok
For a measurable embedding $\phi_i$ of a fixed finite-volume domain,
suppose the pullback volume has nonnegative density $J_i$ relative to
$d\nu$. For an integrable nonnegative source density $f_i$ on its image,
the image mass is $\int J_i(f_i\circ\phi_i)\,d\nu$.
Uniform convergence of these densities to an integrable nonnegative
density gives convergence of image masses. For given chiral matrices
$W_i$ in oriented orthonormal frames, almost-everywhere convergence of
the nine entries and $J_i$ also gives convergence of
$\int J_i\sum_{a,b}(W_i)_{ab}^2\,d\nu$ under measurability and a common
integrable majorant.
\end{lemma}
\begin{proof}
\leanok
Transport measure by the embedding and use the formula for integration
against a measure with density. Apply the finite-volume uniform bound
or dominated convergence to the finite sum of squares.
\end{proof}

\begin{lemma}[The mass of the complement]
\label{lem:source-remainder-split}
\lean{FirstOrderTwistor.measureReal_univ_eq_sum_add_compl, FirstOrderTwistor.measureReal_le_sum_add_uncovered}
\leanok
For finitely many measurable $C_j$ with pairwise null intersections,
put $R=(\bigcup C_j)^c$. Then $\mu(X)=\sum_j\mu(C_j)+\mu(R)$.
If $\mu(S_k)\le d_k$ for finitely many strips, then
\[
 \mu(R)\le\sum_k d_k+\mu(R\setminus\textstyle\bigcup_k S_k).
\]
\end{lemma}
\begin{proof}
\leanok
Use finite additivity up to null sets for the first equation and
subadditivity on the covered portion for the second.
\end{proof}

\begin{theorem}[The energy defect and equality under a strip estimate]
\label{thm:covered-measure-accounting}
\uses{lem:source-remainder-split}
\lean{FirstOrderTwistor.TruncationFamily.sum_componentEnergy_le, FirstOrderTwistor.TruncationFamily.tendsto_limsup_remainder, FirstOrderTwistor.energyIdentityOfCoveredRemainder}
\leanok
Suppose $\mu_i(X)\to e$. Each of finitely many component measures
$\nu_j$ has increasing truncations $K_L^j$ exhausting its space up to
a null set. For fixed $L$, the source images $C_{Li}^j$ are eventually
measurable and pairwise disjoint up to null sets, with
$\mu_i(C_{Li}^j)\to\nu_j(K_L^j)$. With $R_{Li}=(\bigcup_j C_{Li}^j)^c$
and $s=\sum_j\nu_j(X_j)$, one has
\[
 s\le e,\qquad \lim_{L\to\infty}\limsup_{i\to\infty}\mu_i(R_{Li})=e-s .
\]
Suppose also that each strip mass is at most its estimated cost, and
that both the sum of costs and the uncovered mass tend to zero in the
following sense. For every $\varepsilon>0$ and all sufficiently large
$L$, each quantity is eventually at most $\varepsilon$ as $i\to\infty$.
Then $e=s$. The finite strip family may depend on $L$ and $i$.
\end{theorem}
\begin{proof}
\uses{lem:source-remainder-split}
\leanok
At fixed $L$, take the limit of the exact mass split. The finite sum of
truncated masses then converges to $s$ by continuity of measure from
below. Bounding each of the two errors by $\varepsilon/2$ gives $e=s$.
\end{proof}

\begin{corollary}[Bubble energies with supplied geometric budgets]
\label{cor:covered-bubble-energies}
\uses{thm:covered-measure-accounting,lem:scalar-component-budgets,lem:finite-bubble-accounting}
\lean{FirstOrderTwistor.noPositiveBubbleOfCoveredRemainder, FirstOrderTwistor.bubbleEnergyZeroOfRicciYamabeBounds}
\leanok
With one distinguished base and the band and gap hypotheses of
Lemma~\ref{lem:finite-bubble-accounting}, the covered-remainder hypotheses
give zero mass for every bubble and $e=e_0$ if each positive bubble has
strict residual. For the same conclusion it suffices to supply the
limiting Euler identity, $4e_j\le D$, and
$0\le y\le\widehat Y_j<Y_{\rm rd}$ for each positive bubble.
\end{corollary}
\begin{proof}
\uses{thm:covered-measure-accounting,lem:scalar-component-budgets,lem:finite-bubble-accounting}
\leanok
Derive equality by exhaustion. The component budget supplies the strict
residual in the second formulation, so the finite accounting lemma applies.
\end{proof}

\begin{lemma}[Sequence limits of the supplied annular bounds]
\label{lem:annular-bound-limits}
\lean{FirstOrderTwistor.noNeckOfSuppliedCharts, FirstOrderTwistor.noNeckOfSuppliedField, FirstOrderTwistor.annularSmallness_iff_tendsto_limsup, FirstOrderTwistor.smallSumOfFiniteStripErrors}
\leanok
Given the numerical estimate of Theorem~\ref{thm:one-chart-energy}, a
uniform bound on $K,\mathfrak B$ and iterated smallness of $\varepsilon$
give iterated smallness of the annular energy. The same holds for
the intrinsic bound $8\Lambda\mathfrak B+4\|W\|_2\mathfrak d$
when $\Lambda,\|W\|_2$ are uniformly bounded and both error inputs tend
to zero in the iterated sense. For nonnegative energies eventually
bounded above at each fixed large $L$, this is equivalent to the
iterated limsup being zero. Individual estimates sum over a fixed
finite strip family.
\end{lemma}
\begin{proof}
\leanok
Bound the coefficients independently of large $L$ and sufficiently
large $i$ at that $L$. Apply the defining eventual bounds for the
inner limsup and sum over finitely many strips.
\end{proof}

\begin{lemma}[Covered portion of a single end]
\label{lem:diagonal-covered-end}
\lean{FirstOrderTwistor.exists_slow_diagonal, FirstOrderTwistor.exists_slow_diagonal_with_vanishing_annularMeasure}
\leanok
Suppose each fixed truncation eventually passes its comparison test.
Suppose its annular mass between $A$ and $m$ is at most the mass of
a model annulus in the tail beyond $A$ plus an error tending to zero
with $m$. If the model tail mass tends to zero with $A$, there is
a choice $m(i)\to\infty$ passing every test through $m(i)$ eventually
and giving iterated smallness of this covered annular mass.
\end{lemma}
\begin{proof}
\leanok
Choose the largest $m\le i$ for which the tests through $m$ pass.
For each fixed $m$ this prefix eventually passes. Compare at the
selected endpoint and then exhaust the inner tail. Energy up to a
different prescribed outer scale requires a separate estimate.
\end{proof}

\section{Descent from one marked spherical cut}

Let $T=4\pi^2$. The source is an oriented smooth closed homotopy
four-sphere. We use the Hilbert--Schmidt norm of the Weyl endomorphism
on the self-dual two-forms. A chosen smooth ALE compactification
$\widehat X$ has a marked normal ball at its added point. The comparison
must embed the whole protected core and the two collar buffers into a
single source. Several of the lemmas below apply to general measure
spaces or manifolds; their geometric applications use these choices.

\begin{lemma}[Strict exterior margin]
\label{lem:strict-exterior-margin}
\lean{FirstOrderTwistor.exteriorEnergyLtThreshold, FirstOrderTwistor.exteriorEnergyLtThresholdOfStrictError}
\leanok
Let $T,\delta>0$, $e\le2T-\delta$, $v\le e-x+\eta$ and
$0\le\eta\le\delta$. If $x>T$, then $v<T$. If $x\ge T$, the
same conclusion holds when $\eta<\delta$.
\end{lemma}
\begin{proof}
\leanok
At $x=T$, $e=2T-\delta$ and $\eta=\delta$ the exterior bound is $T$.
Either strict inequality in the statement lowers this upper bound.
\end{proof}

\begin{lemma}[Protected core and physical exterior]
\label{lem:cut-core-complement}
\lean{FirstOrderTwistor.FixedTruncationComparison.tendsto_complMass, FirstOrderTwistor.eventually_cappedExteriorEnergy_le_of_coreComparisons}
\leanok
Assume fixed-core image-mass comparison, a finite model of energy $x$
exhausted almost everywhere by increasing protected cores, and total
source energy tending to $e$. Suppose the exterior costs at most the
physical complement of the protected core plus a splice error tending
to zero in the iterated order. For every $\varepsilon>0$ and all
sufficiently large fixed truncations $L$, the exterior energy is at most
$e-x+\varepsilon$ for all sufficiently late source indices at $L$.
\end{lemma}
\begin{proof}
\leanok
Subtract the mass of the measurable embedded image from the total
source mass. At fixed $L$ this difference tends to $e$ minus the model
core mass. Use exhaustion and split the allowance between the remaining
model tail and the splice.
\end{proof}

\begin{lemma}[Cap form and splice estimates]
\label{lem:cut-cap-form}
\lean{FirstOrderTwistor.capPartitionMass, FirstOrderTwistor.capCutoffLoss, FirstOrderTwistor.capYamabeInfimumLowerBound, FirstOrderTwistor.smoothCapSpliceWeylBound, FirstOrderTwistor.smoothCapSpliceBachBound, FirstOrderTwistor.firstOrderCapSpliceWeylBound}
\leanok
For nonzero $L^4$ tests, let $\chi^2+\xi^2=1$ and suppose the two
localized quadratic forms have nonnegative lower coefficients $a,b$.
Supply the IMS identity with error $6\int Fu^2$ where $F\ge0$ and
$6\|F\|_2\le L$. Under the stated measurability and integrability
conditions the infimum of the tested Rayleigh quotients is at least
$\min(a,b)-L$. On a splice of volume at most $V$, a pointwise Weyl
bound $C\eta_2$ costs at most $V C^2\eta_2^2$ in squared norm.
An independent pointwise Bach bound $C'\eta_4$ costs at most
$VC'\eta_4$ in integrated norm. In the first-order model, put
$q,A$ for the old and new squared norms, with the quotient normalization.
If the normalized weighted $L^2$ splice error is at most $\Delta$ and
the cutoff has modulus at most one, the net added energy is at most
$2\sqrt q\,\Delta+\Delta^2$.
\end{lemma}
\begin{proof}
\leanok
Apply the triangle inequality to $u^2=(\chi u)^2+(\xi u)^2$ in
$L^2$ and use Cauchy--Schwarz on the cutoff loss. Divide the resulting
form bound by $\|u\|_4^2$. Integrate the two pointwise bounds on the
fixed splice. For the first-order estimate apply the $L^2$ triangle
inequality on the pair of splices and square.
\end{proof}

\begin{lemma}[One buffer, radius and index for both tests]
\label{lem:cut-ordered-selection}
\uses{lem:strict-exterior-margin,lem:cut-cap-form}
\lean{FirstOrderTwistor.selectFixedSphericalCut}
\leanok
Suppose $e\le2T-\delta$, $x\ge T$, and $0<4\varepsilon<\delta$.
Supply a cutoff loss tending to zero with width, compact core exhaustion
and collar closeness as the radius tends to zero for each width, and
source comparisons at each chosen fixed radius. Assume the same exterior
metric satisfies the cap energy and Yamabe estimates in the normal chart.
One can choose a width $\ell\ge1$, then an arbitrarily small fixed
$s>0$, so that eventually in the source index
\[
 E_{\rm ext}<e-x+4\varepsilon,\qquad E_{\rm ext}<T,\qquad
 Y_{\rm ext}>\tfrac12\min(y_0,Y_{\rm rd}/2)>0.
\]
\end{lemma}
\begin{proof}
\leanok
Choose $\ell$ to control the cutoff loss. In that width make the model
tail and collar error small at a fixed radius. Intersect the eventual
conditions for total energy, protected core comparison and splice error
at that radius. Their four allowances give the displayed energy bound.
\end{proof}

\begin{lemma}[The normal ball and the signed slice]
\label{lem:cut-normal-mark}
\lean{FirstOrderTwistor.normalCutPolarDiffeomorph, FirstOrderTwistor.chartedNormalCap_boundary, FirstOrderTwistor.chartedNormalCut_sides, FirstOrderTwistor.exists_normalCutRadius_for_buffer, FirstOrderTwistor.signedCollar_inside_capNeighborhood}
\leanok
In normal coordinates put $P_s(t,u)=se^{-t}Ju$ for $s>0$, $u\in S^3$
and an orthogonal $J$. The signed cylinder maps diffeomorphically onto punctured four-space,
with inverse given by the logarithmic radius and the normalized direction. The boundary of the marked ball $v\mapsto sJv$
is precisely the middle slice. The positive side lies inside the ball
and the negative side lies outside the closed ball. A radius as small
as desired places the ball and any fixed finite buffer in a chart
about zero. An open neighborhood of an entire middle sphere contains
a two-sided strip of positive uniform width.
\end{lemma}
\begin{proof}
\leanok
The norm is $se^{-t}$, giving signed time $-\log(\|y\|/s)$ and the
three radius tests. The largest radius on a buffer of half-length
$L$ is $se^L$. For the final assertion apply the tube lemma to
$\{0\}\times S^3$ in the inverse image of the open neighborhood.
\end{proof}

\begin{proposition}[Smooth reassembly on a marked collar]
\label{prop:cut-marked-reassembly}
\lean{FirstOrderTwistor.exists_markedExteriorRecognition, FirstOrderTwistor.exists_recognition_eqOn_capNeighborhood, FirstOrderTwistor.originalTypeOfRecognizedExterior, FirstOrderTwistor.SmoothCutGluing.toDiffeomorph_preservesOrientation}
\leanok
Given an oriented recognition $f:\widehat V\to S^4$ and an oriented
ambient extension of its recorded cap germ, a corrected recognition
agrees with the prescribed chart near the closed cap and maps its closed
exterior onto the opposite hemisphere. Supply the retained-side and
normal-ball maps, the core inverse, and the two open covers with their
common signed collar charts. When these charts satisfy the recorded
patch formulas and preserve tangent orientations on their domains, they
induce an oriented smooth diffeomorphism $\Phi:M\to\widehat X$
with those core and exterior restrictions.
\end{proposition}
\begin{proof}
\leanok
Compose $f$ with the inverse ambient extension. Compose the corrected
map with the retained-side and normal-ball charts. The resulting
exterior map and the core inverse coincide on the source overlap, and
the common collar maps onto the target overlap. Glue both maps and
their inverses on these opens. Their differentials on the patch domains
carry the given orientations.
\end{proof}

\begin{proposition}[Marked competitor and the variational residual]
\label{thm:one-cut-same-type-descent}
\uses{prop:cut-marked-reassembly,lem:scalar-component-budgets}
\lean{FirstOrderTwistor.classYamabeFour_pullback, FirstOrderTwistor.chiralFrameEnergy_pullback, FirstOrderTwistor.MarkedMetricTransfer.competitor, FirstOrderTwistor.lowerEnergyOnOriginalType, FirstOrderTwistor.excludesExactMinimum, FirstOrderTwistor.uniformDescentOnOriginalType}
\leanok
Use the specified positive-block coefficients for the energy functional
on all source metrics. On the compactified metric $g$ of energy $x$,
suppose $Y([g])>0$ and the oriented patch maps identify the two blocks
up to almost-everywhere orthogonal conjugacy. The actual pullback
$\Phi^*g$ has the same positive Yamabe constant and energy. The infimum
on $M$ for this coefficient family, denoted by $I_W(M)$, therefore
satisfies $I_W(M)\le x$. The strict scalar residual gives $I_W(M)\le x<e$ under its Ricci and
subround Yamabe hypotheses. With a supplied band-uniform deficit the
stronger inequality is $I_W(M)\le x\le e-\kappa_\delta$.
\end{proposition}
\begin{proof}
\leanok
Pullback identifies the two sets of normalized total scalar-curvature
values as the smooth logarithmic conformal factor varies. Conjugacy
preserves the squared Hilbert--Schmidt norm, and Riemannian volume
changes variables by the same diffeomorphism. Insert the resulting
metric into the source infimum and apply the strict or uniform residual.
\end{proof}

\begin{theorem}[Curvature and descent on the original smooth type]
\label{thm:cut-curvature-descent}
\uses{prop:cut-marked-reassembly,thm:one-cut-same-type-descent}
\lean{FirstOrderTwistor.geometricChiralBlock, FirstOrderTwistor.geometricChiralBlock_pullback, FirstOrderTwistor.geometricPullbackCompetitor, FirstOrderTwistor.geometricPositiveEnergies_eq_of_diffeomorph, FirstOrderTwistor.geometricDescentOnOriginalType, FirstOrderTwistor.oneCutDescent}
\leanok
Compute the positive block of the actual lowered Levi--Civita curvature
in an oriented orthonormal basis by restricting to the normalized
self-dual wedge basis and removing the scalar trace. Let $\mathcal E(M)$
be the energies so obtained from positive-Yamabe metrics and pointwise
oriented orthonormal frames with integrable squared block norm.
An oriented diffeomorphism identifies the two variational sets.
For the single-cut implication, assume $e\le 2T-\delta$, $x>T$ and
$0\le\eta\le\delta$ with $\delta>0$. On the specified exterior supply a
positive Yamabe metric of integrable chiral energy at most $e-x+\eta$.
Assume its strict recognition principle, the hemisphere germ-extension
theorem and the marked local patches for this cut. The map returned by reassembly agrees with the fixed core inverse
and pulls the compactified metric back to the given source. On the
marked cut, a positive-Yamabe compactification with integrable curvature
energy $x$ yields $x\in\mathcal E(M)$. The Route 02 strict residual
implies $\inf\mathcal E(M)\le x<e$. The uniform residual gives
$\inf\mathcal E(M)\le x\le e-\kappa_\delta$ with its extra Yamabe deficit.
\end{theorem}
\begin{proof}
\leanok
Transport the chosen tangent basis by the inverse differential. It
remains oriented and orthonormal for the pulled-back metric. Riemann
naturality identifies all four curvature slots, so the two positive
blocks have identical entries in these paired frames. Pullback preserves
integrability, the energy integral and the Yamabe value. Apply the
construction to the inverse map for the reverse inclusion of the
variational sets. Use the source infimum and the two scalar residuals.
\end{proof}

\begin{openproblem}[Realizing the marked cut from the source hypotheses]
\label{problem:cut-geometric-realization}
\notready
Produce the required smooth compactification, marked full-core embedding,
cap metric estimates and recognized exterior under the hypotheses of the
single-cut theorem. Identifying the formal infimum $\inf\mathcal E(M)$
with the intrinsic $I^+(M)$ in the manuscript requires invariance under
arbitrary changes of oriented orthonormal frame and integrability of the
intrinsic density for each smooth metric. Formal recognition still needs the topology of the
actual capped exterior and the strict sphere theorem. The expected strict
residual depends on the geometric Euler, Ricci and Yamabe inputs. General
existence of a suitable cut and further descent from a positive closed
critical class are separate questions.
\end{openproblem}

\section{Spinor transmission and the residual chiral energy}

In an oriented Riemannian spin four-manifold let $P$ be the twistor
operator on positive half-spinors. We use the curvature operator with
positive sectional curvature on the round sphere, the Hilbert--Schmidt
norm on $W^+$, and the Hilbert norm on spinor-valued two-forms.

\begin{lemma}[Weighted finite cover]
\label{lem:finite-cover-defect-bound}
\lean{FirstOrderTwistor.finiteCoverDefectBound, FirstOrderTwistor.finiteCoverDefectBoundOfIntegrals}
\leanok
Let $r\le u+\sum_{j\in J} w_j$ for a finite $J$. If $c_j>0$ and
$c_j^2 w_j\le 2H_j$ for $j\in J$, then
$r\le u+\sum_{j\in J}2H_j/c_j^2$.
For a nonnegative, integrable density on a measurable residual $S$,
take $w_j$ to be its integral on $S\cap U_j$ and $u$ its integral on
$S\setminus\bigcup_{j\in J}U_j$, for measurable strips $U_j$.
\end{lemma}
\begin{proof}
\leanok
The energy measure restricted to $S$ is finite. Subadditivity on the
union gives the cover inequality even with overlaps. Divide each strip
estimate by its positive squared floor and sum. An empty family leaves
the whole residual in the uncovered term.
\end{proof}

\begin{proposition}[Projected twistor curl in a supplied normal frame]
\label{prop:supplied-normal-spin-defect}
\lean{FirstOrderTwistor.selfDualCliffordProjection_orthogonal, FirstOrderTwistor.spinCovariantDerivative_commutator, FirstOrderTwistor.spinConnectionLift_commutator, FirstOrderTwistor.selfDualCurvatureBlock_metric_isSymm, FirstOrderTwistor.metricProjectedFrameTwistorCurl_normSq, FirstOrderTwistor.metricCoordinateTwistorDefect, FirstOrderTwistor.spinNormalGauge_leviCivita_at_center, FirstOrderTwistor.metricNormalChartTwistorDefect}
\leanok
In the positive quaternionic half-spin fiber the unit self-dual forms
act by $C_a=\sqrt{2}\,I_a$ for $I_a=i,j,k$. Write $m$ for contraction
and $\Pi=1-m^*m/6$. In a coordinate model supply a $C^2$ field, a
differentiable frame connection vanishing at the point, upper and lower
Clifford symbols with their normal-frame values and zero first derivatives,
and an identification of the frame curvature with the metric curvature
in the chosen oriented orthonormal frame. For the coordinate twistor
operator and its exterior covariant derivative,
\[
 \Pi(d^\nabla P\psi)^+_a=-\tfrac12 c(W^+\omega_a)\psi,
 \qquad
 \sum_a|\Pi(d^\nabla P\psi)^+_a|^2
   =\tfrac12|W^+|_{\rm HS}^2|\psi|^2 .
\]
In an aligned boundaryless normal chart the hypotheses on the symbols and
curvature follow from the metric Christoffel matrices and a supplied
twice differentiable moving frame with identity value and zero first
derivative at the center and a differentiable local inverse. For spinor
interpretation choose an oriented orthonormal frame field and its positive
spin lift. The covector connection term vanishes at the normal center.
The squared formula also holds for every ordered two-jet realizing
the metric spin curvature commutator in that frame.
\end{proposition}
\begin{proof}
\leanok
The commutator of two coordinate covariant derivatives equals the
curvature of the supplied connection by symmetry of the ordinary
Hessian. The positive spin lift preserves brackets of skew matrices.
Self-dual alternation of the Clifford correction in $P$ lies in the
image of $m^*$, as does the scalar part of the curvature block.
The trace-free symmetric block lies in $\ker m$, and $mm^*=6$.
The norm of left Clifford multiplication by a real self-dual form
squares to twice the squared form norm times the squared spinor norm.
\end{proof}

\begin{proposition}[Measured strip costs on the physical remainder]
\label{prop:spin-remainder-costs}
\uses{lem:finite-cover-defect-bound,prop:supplied-normal-spin-defect,thm:covered-measure-accounting}
\lean{FirstOrderTwistor.spinDefectFloorBound, FirstOrderTwistor.spinDefectEnergyDecomposition, FirstOrderTwistor.SpinDefectStrip.validOn_ofMetricJets, FirstOrderTwistor.RemainderEstimates.ofSpinDefectStrips, FirstOrderTwistor.RemainderEstimates.ofNormalizedSpinDefectStrips, FirstOrderTwistor.SpinDefectStrip.normalizedCost_le_of_singleton_lpNorm, FirstOrderTwistor.energyIdentityOfSpinorRatioBounds}
\leanok
On each strip intersected with the physical remainder, assume equality
of the restricted physical and strip chiral energy measures. In the
strip metric write $S_j=\sum_a|\psi_{ja}|^2$ and
$H_j=\sum_a|\Pi(d^\nabla P\psi_{ja})^+|^2$. Given the a.e. per-field
identities and the required integrability, a joint floor $S_j\ge c_j^2>0$
gives cost $2\int H_j/c_j^2$. Alternatively use the joint nonzero set
and cost $2\int H_j/S_j$ there. In the latter case the uncovered set
consists of points outside every contributing joint nonzero set. Each
strip uses its own fields and volume. If the summed costs and the uncovered
energy both vanish, first in the eventual sequence bounds at fixed
truncation and then on exhaustion, these estimates give the remainder
hypothesis of Theorem~\ref{thm:covered-measure-accounting}.
\end{proposition}
\begin{proof}
\leanok
Sum the identities over fields and integrate. On the positive locus the
ratio is exactly half the curvature density. On a floor set use its
lower bound in the unnormalized identity. Transfer only the restricted
energy measure. At each pair of indices take the finite union and retain
its complement once. The number of strips may vary with both indices.
\end{proof}

\begin{lemma}[First-order perturbations and a differentiated interior bound]
\label{lem:spin-strip-derivative}
\uses{lem:buffered-euclidean-interpolation}
\lean{FirstOrderTwistor.spinStrip_operatorPerturbation, FirstOrderTwistor.spinStrip_bufferedCoefficientBound, FirstOrderTwistor.spinStrip_sumSlabBounds, FirstOrderTwistor.spinStrip_absorbTwistedDirac, FirstOrderTwistor.spinStrip_projectedDerivative, FirstOrderTwistor.harmonicMetricCoordinateTwistedDiracDerivative_normSq}
\leanok
For a supplied operator difference $a\nabla+b$ on a reference strip,
gradient norm bound one and Sobolev constant $S$, the operator norm
from the first-order space to $L^2$ is at most
$\|a\|_\infty+S\|b\|_4$. For the interior estimate on $K\Subset K'$,
supply the flat cutoff estimate, small twisted-operator perturbation,
the differentiated harmonic equation with its cutoff and Ricci forcing
bounds, the return to the actual norms, an incoming error $d$, a
first-order comparison error $\epsilon_1$, and a floor $|\psi|\ge c>0$
on $K$. With constants from those estimates,
\[
 \left\|\frac{|\Pi(d^\nabla P\psi)^+|}{|\psi|}\right\|_{2,K}
 \le \frac{C}{c}\left(\epsilon_1+d+
       \|\,|\mathrm{Ric}|\,|\psi|\,\|_{2,K'}\right).
\]
\end{lemma}
\begin{proof}
\leanok
Apply H\"older with exponents $(\infty,2)$ and $(4,4)$ in the supplied
measure. Buffered interpolation bounds the coefficient budget under its
chart comparison hypotheses. For the derivative estimate, absorb the
relative error in the flat estimate when $C_0\epsilon\le1/2$. Bound the
cutoff term on the outer buffer, then divide the projected norm by the
floor on the inner set.
\end{proof}

\begin{proposition}[Relative Bochner criterion without a positive minimum]
\label{prop:spin-relative-bochner}
\lean{FirstOrderTwistor.normalizedSpinCutoff_fderiv_bound_frame, FirstOrderTwistor.spinRelativeBochner_bound, FirstOrderTwistor.spinMixedRelative_sum_le, FirstOrderTwistor.spinCutoffIntegral_in_gauge, FirstOrderTwistor.RemainderEstimates.ofSpinCutoffCover, FirstOrderTwistor.energyIdentityOfSpinBochnerBudgets}
\leanok
For measured joint spinor mass $S$, derivative density $Q$, and cutoff
$\chi$, set $v^2=Q/S$ on the joint nonzero locus. Assume $S>0$ almost
everywhere on the cutoff support. Given the integrated regularized
twisted Bochner inequality and finite relative costs in the strip metric,
\[
 \int\chi^2|W^+|^2 \le C\int
 \big(\chi^2|\mathrm{Ric}|^2+\chi^2|\mathrm{Rm}|v^2+
          \chi^2v^4+|d\chi|^2v^2\big).
\]
For a finite family of measured patches the sum of the mixed terms
is at most $\sqrt{\mathcal K\mathcal V}$, where $\mathcal K$ and
$\mathcal V$ are the sums of $\int\chi^2|\mathrm{Rm}|^2$ and
$\int\chi^2 v^4$ in the respective gauges. Where the cutoff squares
sum to at least one in the physical energy measure, the weighted costs
bound the covered energy. The mass of the remaining set enters once.
For a uniform constant, bounded $\mathcal K$, vanishing Ricci-square,
relative fourth-power and cutoff costs, and vanishing uncovered energy
in the iterated order, Route~02 gives the component energy identity
under its truncation and total convergence hypotheses.
\end{proposition}
\begin{proof}
\leanok
Writing $D=S+\epsilon$, differentiating $\chi/\sqrt D$ gives
$|d(\chi/\sqrt D)|^2Q\le 2|d\chi|^2 Q/D+2\chi^2(Q/D)^2$
in the Hilbert frame model. Use the factor $S/(S+\epsilon)$ on the left and dominate the regularized
right side by the finite unregularized costs. Dominated convergence
takes $\epsilon$ to zero. Apply Cauchy--Schwarz first to each integral,
then to the finite sum over patches.
\end{proof}

\begin{proposition}[Incoming modes and attenuation]
\label{prop:spin-mode-transmission}
\lean{FirstOrderTwistor.forcedCylinderMode_equation, FirstOrderTwistor.cylinderMode_energy_balance, FirstOrderTwistor.fixedDomainDiracPerturbation_inverse_norm_le, FirstOrderTwistor.fixedDomainDiracScattering_sameEnd_norm_le, FirstOrderTwistor.diracWeighted_remoteTrace_norm_le, FirstOrderTwistor.roundCylinderMidpoint_jointFloor_Lp}
\leanok
For $\partial_t+A$ on $[0,L]$, a positive signed mode uses its left
datum, a negative mode its right datum. On a supplied complete
orthonormal modal model with $|\lambda|\ge3/2$, a finite family with
joint floor $c^2$ on the middle slice satisfies
\[
 c^2\operatorname{vol}(Y)
 \le e^{-3L/2}\sum_a(\|x_a\|_2^2+\|y_a\|_2^2)
\]
for the homogeneous solutions and their signed incoming components.
For a boundedly invertible reference operator with inverse norm at most
$C$, a fixed-domain perturbation of norm at most $\delta$ has inverse
norm at most $C/(1-C\delta)$ if $C\delta<1$.
With fixed incoming and outgoing trace maps and a vanishing same-end
reference block, the resolvent identity bounds the same-end reflection
by the perturbation norm times the two inverse and trace factors.
\end{proposition}
\begin{proof}
\leanok
Integrate each signed scalar equation from its incoming end. At distance
$L/2$ the squared amplitude decreases by at least the indicated exponential
factor. Sum using the Hilbert norm and integrate the floor on the slice.
The perturbation bound follows by a Neumann inverse on the fixed spaces.
\end{proof}

\begin{proposition}[Matching and regular jump corrections]
\label{prop:spin-core-and-jump}
\lean{FirstOrderTwistor.matchSuppliedCoreTraceGraphs, FirstOrderTwistor.closedJumpCorrectionOfRightInverse, FirstOrderTwistor.criticalJointFloorOfJumpBounds}
\leanok
For two supplied bounded core graphs $G_L,G_R$ and propagation norm
bounds $e^{-3L/2}$ in both directions, the affine graph equations have
a unique solution when $\|G_L\|\|G_R\|e^{-3L}<1$. For an original-source
jump problem supply four measured face maps, a right lift of both jumps,
a compatible closed inverse, and zero-jump gluing. There is a unique
piecewise homogeneous correction with the specified jumps. If its regular
fields satisfy $\|\psi_a\|_{8/3}\le C J_a$ and their joint floor on $U$ is
$c^2$, then $c^2\operatorname{vol}(U)^{3/4}\le C^2\sum_a J_a^2$ for
nonnegative jumps and bound constants as in the critical norm estimate.
\end{proposition}
\begin{proof}
\leanok
Substitute one affine trace equation into the other and apply contraction.
For the jump correction subtract from the piecewise lift the restriction
of the global inverse applied to its Dirac error. Use the $L^{4/3}$
triangle inequality on the sum of squared field norms.
\end{proof}

\begin{proposition}[Attainment in a supplied compact gauge]
\label{prop:spin-compact-kernel}
\lean{FirstOrderTwistor.existsUnitKernelOfCompactEmbedding, FirstOrderTwistor.weylBlock_eq_zero_of_dense_spinor}
\leanok
Fix a real Hilbert domain with compact inclusion into the reference
$L^2$ space. A bounded sequence whose image norms tend to one, whose
operator defects tend to zero, and whose operators converge in norm to
$P$ yields a kernel element with unit image norm. Separately, a continuous
Weyl block vanishes throughout a chart if its projected defect vanishes
and the spinor is nonzero on a dense set satisfying the pointwise identity.
\end{proposition}
\begin{proof}
\leanok
Take a strongly converging image subsequence and a weak cluster point in
the bounded Hilbert domain. Continuous linear maps identify the kernel
and the image norm of this point. For the Weyl block, apply the pointwise
identity on the dense set and continuity to all coefficients.
\end{proof}

\begin{openproblem}[Differentiated control over an entire physical residual]
\label{problem:physical-spin-residual}
\notready
Construct measured fields or local gauges on the actual residual of the
geometric sequence with summably small independent differentiated bounds
and negligible uncovered energy. A Green pole, compatible endpoint germs,
a fixed-domain Sobolev inverse, and a normalized limit on the same
smooth manifold each require their own geometric realization.
\end{openproblem}

\section{Fixed-type minimization in the master manuscript}
\label{sec:master-variational}

Fix a closed connected oriented smooth homotopy four-sphere \(M\), smooth
Riemannian metrics, and \(T=4\pi^2\). Write \(E_\pm(g;U)\) for the
Hilbert--Schmidt energy of the corresponding Weyl operator on \(U\).
Signature gives equality of the two totals on \(M\). Regional energies
retain their chirality. A positive class with \(E_+\le T\) recognizes
the smooth type by the sphere and equality theorems of Chang, Gursky
and Yang. This implication and its topological hypotheses are external
dependencies in Lean. The dependencies below distinguish checked conditional comparisons from
geometric realization on a candidate smooth type.

\begin{definition}[Admissible geometric energies]
\label{def:master-admissible}
\lean{FirstOrderTwistor.geometricPositiveEnergies}
\leanok
For an oriented closed smooth Riemannian four-manifold in the geometric
library, let \(A_M\subset\mathbb R\) consist of the energies witnessed by
a smooth metric with positive class Yamabe constant and a pointwise
oriented orthonormal frame in which the squared chiral block of its
Levi--Civita curvature is integrable. The energy uses the volume of
that metric. The witness is a choice of basis at each point.
\end{definition}

\begin{lemma}[Lower bound and a marked competitor]
\label{lem:master-competitor}
\uses{def:master-admissible}
\lean{FirstOrderTwistor.geometricPositiveEnergies_bddBelow, FirstOrderTwistor.geometricEnergy_mem_of_pullback}
\leanok
\(A_M\) is bounded below. Given an orientation-preserving
diffeomorphism from \(M\) to a closed smooth oriented four-manifold \(X\)
and an admissible framed metric on \(X\), its energy belongs to \(A_M\).
\end{lemma}
\begin{proof}
\leanok
The integrated squared norm is nonnegative. Pull back the metric and
transport its frame along the given diffeomorphism. The Yamabe constant,
integrability and the curvature integral transfer along that same map.
\end{proof}

\begin{proposition}[Independent replacements on the original type]
\label{prop:master-replacement}
\uses{lem:master-competitor,prop:ed03-measured-comparisons}
After a positive-class witness is supplied, write \(I=\inf A_M\) in
\(\mathbb R\). Let \(m_j\) be the raw chiral masses of finitely many
disjoint marked standard balls in a source of energy \(e\). If each
individual repair gives an admissible metric on \(M\) of energy at
most \(e-m_j+\gamma_j\), then
\[
  \sum_j m_j\le N(e-I)+\sum_j\gamma_j .
\]
Apply the infimum lower bound to each competitor before summing.
For the manuscript's universal entry question the convention is
\(\inf\varnothing=+\infty\). The geometric identification of \(A_M\)
with the manuscript functional also needs independence from unrelated
oriented frames.
\end{proposition}

\begin{corollary}[Independent repairs and the uncovered mass]
\label{cor:master-repair-residual}
\uses{prop:master-replacement,lem:ed03-measured-residual,lem:ed03-repair-threshold}
On a nonstandard type in the band, fix \(\varepsilon>0\).
Suppose a subsequence admits the marked repairs above, with bounded
cardinality at this tolerance, total surcharge of limsup at most
\(b(\varepsilon)\) and full uncovered chiral mass of limsup at most
\(\omega(\varepsilon)\). Then \(I\le b(\varepsilon)+\omega(\varepsilon)\).
If such families exist for arbitrarily small tolerances and the sum
tends to zero, this contradicts \(I\ge T\). A physical volume estimate
supplies the complement bound under a separate uniform absolute
continuity hypothesis for chiral mass.
\end{corollary}

\begin{proposition}[Recognition at an attained threshold]
\label{prop:master-recognition}
On a closed oriented smooth homotopy four-sphere, a positive conformal
class with \(E_+(g)\le T\) gives the standard smooth type. At equality
apply the equality classification to the actual class and exclude
the other topologies by simple connectivity and homology.
An unattained value \(I=T\) remains a possibility in the band.
\end{proposition}

\begin{proposition}[Nonsplitting for two marked caps]
\label{prop:master-nonsplitting}
\uses{lem:master-competitor,prop:cut-marked-reassembly,lem:source-remainder-split,prop:ed03-measured-comparisons}
Let the two full raw masses be \(m_1+m_2=e\). Assume two positive
marked caps with costs \(E_j\le m_j+\gamma_j\), \(\gamma_j\ge0\), and
\(e+\gamma_1+\gamma_2<2T\). On a nonstandard candidate at least one
cap is recognized as standard and the marked reassembly makes the other
cap \(s\) a competitor on \(M\). Hence \(m_{3-s}\le e-I+\gamma_s\).
Along a minimizing sequence two uniformly positive full sides are
impossible if both costs tend to zero. Either label may survive.
The geometric theorem needs recognition of the capped homotopy spheres
and the actual attaching maps.
\end{proposition}

\begin{lemma}[Positive paths below the scalar-flat wall]
\label{lem:master-wall}
\uses{lem:scalar-component-budgets,lem:ed03-path}
On the original closed smooth type, a zero Yamabe class has a
scalar-flat representative \(q\), for which
\(E_+(q)=2T+\frac14\int_M|\mathrm{Ric}_q|^2\,dV_q\).
Along a smooth path starting in a positive class with \(E_+<2T\)
throughout, the class stays positive. Its Yamabe constant is at least
\(\sqrt{48(2T-E_+)}\). The geometric inputs include the signature and
Euler formulas, conformal invariance, the Yamabe representative and
continuity of the first conformal-Laplacian eigenvalue along this path.
\end{lemma}

\begin{lemma}[First-order interpolation in the splice coordinates]
\label{lem:master-splice}
\uses{lem:cut-cap-form}
Fix the marked annulus and its smooth cutoff \(\psi\) in advance. In the
oriented comparison metric \(m=\delta+v\) put \(a=\|v\|_\infty\),
\(d=\|Dv\|_4\) and \(K=\|\mathrm{Rm}_m\|_{2}\) on the splices for the last
norm. Under the first-order comparison assumptions, the interpolants
\(\delta+(1-s\psi)v\), \(0\le s\le1\), incur an \(L^2\) chiral density
error at most \(C(aK+d^2+d+a)\), uniformly in \(s\). For incoming
splice energy \(q\) and error bound \(\Delta\), the net surcharge is at
most \(2\sqrt q\,\Delta+\Delta^2\). The checked norm estimate consumes
the density error. The curvature commutator and the comparison of chiral
projections in this chart are the proposed geometric proof.
\end{lemma}

\begin{proposition}[The short marked cut]
\label{prop:master-short-cut}
\uses{prop:master-nonsplitting,lem:master-splice,lem:master-wall}
On a nonstandard type take positive volume-one Yamabe representatives
with \(e_i\to I<2T\). Supply one fixed short oriented spherical chart
on each source, including its separating middle slice. In the conformal
comparison gauge assume uniform and strong \(W^{1,4}\) convergence to
the Euclidean metric and bounded full \(L^2\) curvature on the splices.
Straighten in this chart on the joined smooth type, stretch the exact
product for each fixed source, and transfer compactly supported tests
to each completed cap. Zero point capacity and the path bound give both
positive cap signs. The smaller full raw chiral mass tends to zero.
\end{proposition}

\begin{theorem}[Local replacement on a marked comparison]
\label{thm:master-local-replacement}
\uses{lem:master-competitor,lem:master-wall,lem:master-splice,lem:ed03-weighted-defect}
On an oriented open domain supply embeddings of neighborhoods of a
compact exhaustion with \(\Phi_i^*g_i=u_i^2m_i\), \(u_i>0\).
Assume \(m_i\to h\) locally uniformly and strongly in \(W^{1,4}\),
\(h\) smooth, and local full \(L^2\) curvature bounds in \(m_i\).
For the above minimizing sequence the two chiral matrix densities
converge strongly in \(L^2_{\rm loc}\). The limit minimizes against
sufficiently small smooth compactly supported variations and is
Bach-flat. The proposed proof tests an actual positive competitor on
the original type. For \(r_i=F_\pm(m_i)-F_\pm(h)\) it gives
\(\int(2\chi-\chi^2)|r_i|^2\le e_i-I+o(1)\) for each fixed cutoff.
The existing two-jet theorem has stronger source regularity.
\end{theorem}

\begin{lemma}[Physical chiral quantiles]
\label{lem:master-quantiles}
\uses{lem:source-remainder-split,lem:scalar-component-budgets,lem:ed03-yamabe-quantiles}
In physical distance set \(r_i(\theta)\) equal to the least ball radius
capturing \(\theta e_i\) over all centers, and \(r_i=r_i(1/2)\).
Assume positive total chiral mass bounded away from zero, uniformly
bounded total full curvature energy, and continuous ball masses with
attainment as in the smooth compact case. If
\(r_i(\theta)/r_i\to0\) on a subsequence for some \(0<\theta<1/2\),
use its attaining inner center. The entire exterior at the median
radius about that same center carries at least \(e_i/2\). Among the
disjoint intervening annuli of any fixed log width one has full
curvature integral tending to zero. The other case gives
\(\lim_{b\downarrow0}\limsup_i\sup_x E_+(B(x,b r_i))/e_i=0\).
\end{lemma}

\begin{proposition}[A fixed-distortion harmonic shell]
\label{prop:master-harmonic-annulus}
\uses{lem:master-quantiles,prop:master-short-cut}
On a separated subsequence at the positive-class infimum, supply a
whole oriented chart \(1/8<|x|<8\) in the intervening physical shell.
The attaining inner ball and entire median exterior must lie on
opposite sides and outside the closed trim \(1/4\le |x|\le4\).
Assume a constant physical rescaling has vanishing full curvature
energy on the chart and is uniformly within a fixed sufficiently small
distortion of Euclidean. If the coordinates are harmonic, or if
\(\|D(\Delta_{m_i}x)\|_2\to0\), this configuration is impossible on
the nonstandard type. The chart and the distortion are fixed before
the source limit.
\end{proposition}

\begin{theorem}[Complete physical limits and curvature scales]
\label{thm:master-complete-limit}
\uses{thm:master-local-replacement,prop:master-short-cut}
For a positive-class minimizing sequence on a fixed nonstandard type
in the band, let \(a_i\downarrow0\). Suppose the constant rescalings
\(a_i^{-2}g_i\) converge locally uniformly and strongly in \(W^{1,4}\)
on oriented embeddings of neighborhoods of a compact exhaustion to
a smooth connected complete boundaryless limit. The limit is flat.
The argument needs ends, link topology and a marked cut whenever a
nonflat limit is tested. As a further target, for
\(K_i=\max_M|\mathrm{Rm}_{g_i}|\to\infty\), prove that full energy on
physical balls of radius \(A/\sqrt{K_i}\) tends to zero uniformly in
their centers for every fixed \(A\). Bounded \(K_i\) on a subsequence
gives smooth closed attainment on \(M\).
\end{theorem}

\begin{proposition}[Physical Ricci scales]
\label{prop:master-ricci-scale}
\uses{thm:master-complete-limit}
In the same band take a physical scale \(a_i\to0\), points \(x_i\)
and one fixed \(2<p<\infty\). Suppose that for each fixed \(D\) one has
\(a_i^{2-4/p}\|\mathrm{Ric}_{0,g_i}\|_{L^p(B_{g_i}(x_i,Da_i))}\le C_D\)
eventually on the same subsequence. Full curvature energy on
\(B_{g_i}(x_i,Aa_i)\) then tends to zero for every fixed \(A\).
Bounded global \(L^p\) norm of the trace-free Ricci tensor on a
subsequence gives a smooth closed positive Yamabe minimizer on \(M\).
\end{proposition}

\begin{proposition}[Recognition for one Euclidean hypersurface immersion]
\label{prop:master-immersion}
\uses{lem:routes-middle-weyl,lem:routes-shape-deficit,lem:routes-indexed-mass,prop:routes-middle-count,lem:ed03-equality,cor:ed03-balanced-endpoint,lem:ed03-codazzi-jets}
For a smooth immersion of \(M\) in \(\mathbb R^5\), let \(g\) be the
induced metric, \(A\) its shape operator, \(J=\int_M(-\det A)_+\,dV_g\)
and \(B=\|A^\circ\|_{L^4(g)}^2\). Each of \(J<8\pi^2/3\),
\(E_+(g)\le64\pi^2/9\), or \(B\le16\sqrt{2}\,\pi/3\) implies that
\(M\) is standard. The energy criterion also holds in the induced
conformal class, independently of its Yamabe sign. Separate the
pointwise principal-curvature inequalities from the Morse height
count, smooth handle reduction and exclusion of equality by the
Gauss and Codazzi equations. Constructing a sufficiently controlled
immersion on every type is open.
\end{proposition}

\begin{openproblem}[Realization and the remaining alternatives]
\label{problem:master-realization}
\uses{prop:master-harmonic-annulus,cor:master-repair-residual,prop:master-ricci-scale,prop:master-immersion,prop:master-recognition}
Construct entry on an arbitrary smooth type and find a qualifying
marked cut or same-type repair on a minimizing sequence. In the diffuse
case charge the full complement of the repaired balls. For a supplied
whole marked base or complete constant rescaling, audit regularity,
end topology, compactification and all mass outside fixed truncations.
The master derives flatness of smooth complete shrinking limits and
loss of full energy on fixed maximum-curvature-scale balls. A closed
nonround minimizer and a round restart at the same energy still require
rigidity, a strict improvement, or termination. The working immersion,
physical \(L^p\) Ricci (\(p>2\)), and fixed-distortion harmonic-annulus
criteria identify additional hypotheses for these routes.
\end{openproblem}

\section{Descent, immersed curvature and Morse heights}
\label{sec:parallel-investigations}

This section records the third program edition. Checked implications below specify their witnesses and their remaining geometric inputs. For
physical shell selection and the residual of a repair family, the targets
are Lemma~\ref{lem:master-quantiles} and
Corollary~\ref{cor:master-repair-residual}. A supplied shell is tested by
Proposition~\ref{prop:master-harmonic-annulus}. The following results use
the normalization and oriented smooth category of
Section~\ref{sec:master-variational}.

\begin{proposition}[A joint competitor on the original smooth type]
\label{prop:routes-joint-descent}
\lean{SmoothPoincare.jointCompetitorGain_le_excess, SmoothPoincare.excludesStrictJointCompetitor}
\leanok
\uses{lem:master-competitor}
Fix a nonempty set of admissible positive-class energies on \(M\) with
finite infimum \(I\). A source has energy \(e\). Suppose a marked
replacement of a region, including its transitions, supplies a smooth
positive metric on the original type, whose total energy is at most
\(e-m+\gamma\). Here \(m\) is the removed raw energy and \(\gamma\)
includes the whole replacement cost and any net change outside the
region. Then \(m-\gamma\le e-I\). A strict reverse inequality for one
actual competitor gives a contradiction.
\end{proposition}

\begin{proof}
\leanok
The marked pullback supplies membership on the original type. Subtract the
infimum bound from the complete energy comparison.
\end{proof}

\begin{lemma}[Weyl energy on the middle-signature locus]
\label{lem:routes-middle-weyl}
In the induced metric of a smooth Euclidean hypersurface, take principal
curvatures \((-a,-b,c,d)\), all four parameters positive. The three pair
sums are \(ab+cd,-ac-bd,-ad-bc\). For either Weyl endomorphism the squared
Hilbert--Schmidt norm is one twelfth of the sum of the three squared
differences between these sums. It is at least \((8/3)abcd\), with
equality precisely when \(a=b=c=d\). The scalar algebra and the Gauss
equation identifying the actual curvature are separate proof targets.
\end{lemma}

\begin{lemma}[Shape-operator deficit]
\label{lem:routes-shape-deficit}
For arbitrary real principal curvatures in the same block normalization,
\[
  \tfrac16 |A^\circ|^4-|W^\pm|_{\rm HS}^2
     \ \ge\ \tfrac89(-\det A)_+ .
\]
The tuple \((-1,3,3,3)\) realizes equality with negative determinant.
\end{lemma}

\begin{lemma}[Indexed Gauss mass]
\label{lem:routes-indexed-mass}
Let \(m_2(u)\) count the index-two critical points of the height in a
regular direction \(u\) for one smooth immersion of a closed oriented
homotopy four-sphere in \(\mathbb R^5\). With normalized spherical
probability, \(\omega=|S^4|\) and \(U\) the middle-signature locus,
\[
  \mathbb E[m_2]=\frac{2}{\omega}\int_U\det A\,dV_g .
\]
Both normal directions contribute. The area formula counts with
multiplicity on the smooth domain. Identify a measurable count and the
null exceptional directions before integrating.
\end{lemma}

\begin{proposition}[The topological floor for the middle count]
\label{prop:routes-middle-count}
On a nonstandard smooth homotopy four-sphere every Morse function has at
least two index-two critical points. The proposed formal argument uses
the actual extremal handlebody reductions, preserves the framed
attachments, and applies the one-middle-pair recognition theorem and
final disk gluing when the reduced count is at most one.
\end{proposition}

\begin{lemma}[Selecting a two-middle height]
\label{lem:routes-height-selection}
\lean{SmoothPoincare.prob_twoMiddle_gt_threeQuarters, SmoothPoincare.twoMiddle_probability_boundary}
\leanok
On a probability space let \(m\) be an integrable, measurable
natural-valued count. If \(m\ge2\) almost everywhere and
\(\mathbb E[m]<9/4\), then \(\Pr(m=2)>3/4\).
\end{lemma}

\begin{proof}
\leanok
Outside the two-count event the integer count is at least three almost
everywhere. Integrate this extra unit. A probability of one quarter
at three and three quarters at two gives equality at the excluded mean.
\end{proof}

\begin{corollary}[Heights of a subcritical immersion]
\label{cor:routes-subcritical-heights}
\uses{lem:routes-middle-weyl,lem:routes-indexed-mass,prop:routes-middle-count,lem:routes-height-selection}
If the induced class of the specified immersion on a nonstandard
candidate has \(E_+(g)<8\pi^2\), more than three quarters of directions
have exactly two middle critical points. The direction count uses no
Yamabe sign. Recognition still requires control of a selected diagram.
\end{corollary}

\begin{proposition}[A marked mixed cancellation]
\label{prop:routes-mixed-cancellation}
\uses{prop:routes-middle-count}
Suppose a reduced smooth handle presentation of \(M\) has one 0-handle,
one 1-handle, two 2-handles, one 3-handle and one 4-handle. Supply a
framed isotopy or a sequence of justified slides that preserves the
marked presentation and gives a transverse single geometric
intersection of one 2-handle attaching circle with the 1-handle belt
sphere. Smooth handle cancellation leaves at most one middle handle,
so \(M\) is standard. Producing this geometric intersection from a
selected height is open.
\end{proposition}

\begin{openproblem}[Coverage on each candidate smooth type]
\label{problem:routes-coverage}
\uses{prop:master-recognition,prop:master-immersion,cor:routes-subcritical-heights,prop:routes-mixed-cancellation,prop:master-short-cut,cor:master-repair-residual,prop:routes-joint-descent}
For each candidate, produce a recognized class, an immersed subcritical
height with a recognized marked diagram, or a positive-class minimizing
sequence on which a strict contradiction is realized. Quantitative
construction and coverage must refer to the same metric, immersion,
or chosen physical sequence in the implication that uses it.
\end{openproblem}

\subsection{Conditional comparisons on supplied witnesses}

\begin{proposition}[Measured replacements and two full cut sides]
\label{prop:ed03-measured-comparisons}
\uses{lem:master-competitor,prop:cut-marked-reassembly}
\lean{SmoothPoincare.geometricRepairedMass_le_excess_sum_of_pullback, SmoothPoincare.geometricRepairedMass_le_excess_sum, SmoothPoincare.geometricDiscardedMass_le_excess, SmoothPoincare.MarkedSideReplacement.reassemble}
\leanok
Fix the geometric positive-energy set of the original oriented smooth
four-manifold and use its physical chiral density in a specified frame.
For individually marked admissible repairs of known global cost,
\(\sum_j m_j\le N(e-I)+\sum_j\gamma_j\).
For two full measured cut sides, suppose their union is the source,
their intersection is volume-null, both caps are positive, and recognition
and smooth collar reassembly data are supplied for either choice of the
discarded cap. If \(e+\gamma_1+\gamma_2<8\pi^2\), a surviving side
\(s\) gives an oriented competitor with
\(m_{3-s}\le e-I+\gamma_s\). The resulting map carries the entire
opposite raw side onto the marked cap.
\end{proposition}
\begin{proof}
\leanok
Integrate the source curvature on the specified sets. For independent repairs,
pull each metric back and sum its comparison with the same infimum.
For the cut, the strict total bound selects a cap below the recognition
threshold. Use its corrected marking on a neighborhood to identify the
other cap with the source. Apply the infimum bound to that pullback.
\end{proof}

\begin{lemma}[A common Yamabe floor along an actual metric family]
\label{lem:ed03-path}
\uses{def:master-admissible}
\lean{SmoothPoincare.positiveAlongSubcriticalPath, SmoothPoincare.admissibleAlongSubcriticalMetricPath}
\leanok
Suppose a connected parameter family has a continuous sign detector,
positive initially, with the stated positive and zero sign identifications
for the class Yamabe constants of its metrics on the original type.
Supply integrable chiral densities and the normalized inequality
\(Y_t^2\ge48(8\pi^2-E_t)\).
If \(E_t\le8\pi^2-\delta\) throughout for one \(\delta>0\), every energy
is admissible and every \(Y_t\ge\sqrt{48\delta}\).
\end{lemma}
\begin{proof}
\leanok
At a zero of the sign detector the Euler inequality forces
\(E_t\ge8\pi^2\). The intermediate value theorem propagates the positive
sign. Take the square root of the Euler inequality and use each actual
framed metric to witness membership.
\end{proof}

\section{Global Gauss area and the cost of marked collars}
\label{sec:ed03-global-gauss-caps}

\begin{theorem}[Oriented normal and the ordinary area of the direction sphere]
\label{thm:ed03-global-normal}
\uses{lem:ed03-normal-heights}
\lean{SmoothPoincare.orientedGaussVector_volume_pos, SmoothPoincare.orientedGaussNormal_smooth, SmoothPoincare.orientedPrincipalFrameAt_middle_order, SmoothPoincare.gaussRoundArea_eq_gaussConeArea, SmoothPoincare.gaussRoundArea_univ, SmoothPoincare.gaussRoundArea_antipodal}
\leanok
Let $F:(M^4,g)\to\mathbb R^5$ be a smooth oriented isometric hypersurface
immersion. Fix the ambient orientation by an orthonormal basis. The unique
positive unit normal defines a smooth map $N:M\to S^4$. At each point
there is an oriented orthonormal principal frame for this normal with
nondecreasing signed principal curvatures. Ordinary Riemannian area on
the Euclidean unit sphere agrees with radial cone area, has total mass
$\omega=8\pi^2/3$, and is invariant under the antipodal map.
\end{theorem}
\begin{proof}
\leanok
Normalize the oriented cross product of the four tangent vectors of an
oriented chart. A positive change of basis leaves its normalization fixed,
also where the second fundamental form vanishes. Diagonalize the symmetric
shape operator, order its eigenvectors, and adjust the source orientation.
For the area calculation use $(u,t)\mapsto e^tu$. The pulled-back ambient
volume has density $e^{5t}$ relative to round area times $dt$, with integral
$1/5$ over $t<0$. Comparing the volumes of cones gives equality of the
two measures on every Borel set.
\end{proof}

\begin{theorem}[Open middle locus and the indexed area identity]
\label{thm:ed03-ordinary-indexed-area}
\uses{thm:ed03-global-normal,prop:ed03-gauss-sheets}
\lean{SmoothPoincare.isOpen_middlePrincipalLocus, SmoothPoincare.gaussParamDensity_eq_mul_frameJacobian, SmoothPoincare.ImmersionGaussNormal.measuredSheets_of_principal, SmoothPoincare.orientedImmersion_indexedGaussMean}
\leanok
On a nonempty compact Polish source modeled on Euclidean four-space, let $U$
be the locus of signature $(2,2)$ and $J=|\prod_{j=1}^4 k_j|$. Write
$\mu$ for ordinary round area divided by $\omega$. For a unit direction $u$
let $m_2(u)$ count the nondegenerate index-two critical points of
$\langle F,u\rangle$ on the source. The locus $U$ is open, $J$ is
continuous, the point set is finite almost surely, and $m_2$ is measurable
and integrable, with
\[
    \int_{S^4}m_2(u)\,d\mu(u)=\frac{2}{\omega}\int_U J\,dV_g .
\]
\end{theorem}
\begin{proof}
\leanok
In a single smooth normal chart use the determinant of the signed second
form, and two fixed tangent coordinate vectors where the quadratic values
have opposite strict signs. Positive determinant excludes degeneracy and
forces an even number of negative eigenvalues. The two evaluations exclude
both definite signatures. All three strict tests persist near the center.
Changing the unit normal preserves middle signature. In the same source
parametrization the target density of the composed Gauss map equals the
source density times its orthonormal absolute Jacobian. Ratios on charts
give continuity of $J$ even at singular points. Integrate on countably many
regular sheets and use the disjoint fibers over $u$ and $-u$. Multiplicity
counts points in $M$, including distinct points with the same immersed image.
\end{proof}

\begin{theorem}[Morse directions at the immersed threshold]
\label{thm:ed03-morse-threshold}
\uses{thm:ed03-ordinary-indexed-area,cor:ed03-strict-actual-height}
\lean{SmoothPoincare.riemannianVolumeMeasure_criticalValues_eq_zero, SmoothPoincare.orientedHeight_ae_nondegenerate, SmoothPoincare.inducedEnergy_strict_of_orientedHeightFloor, SmoothPoincare.exists_morseHeight_le_one_of_inducedEnergy_le}
\leanok
Almost every round direction gives a smooth height with nondegenerate
Hessian at each critical point. On the compact source above, use the
ordered frames of the constructed normal. With integrable squared
chiral norm and continuous middle deficit, energy at most $64\pi^2/9$
gives a Morse height with a finite set of at most one index-two point.
\end{theorem}
\begin{proof}
\leanok
Apply the equal-dimensional Sard theorem in charts and transfer the null
critical values to Riemannian volume. Remove both this set and its antipode.
The signed Weingarten formula then gives a nondegenerate Hessian at
every critical point of the height. If almost every direction had at
least two middle points, the indexed identity and exclusion of open
balanced principal frames would force strict energy above the threshold.
Intersect the conull Morse and finiteness conclusions before selecting
a direction. Coincident critical values are allowed.
\end{proof}
Smooth recognition from one such direction still requires a marked handle
presentation and its boundary gluings on the original smooth source.

\begin{proposition}[Physical chart tests and moving frames]
\label{prop:ed03-chart-frame-tests}
\uses{prop:ed03-anchored-shell,lem:ed03-splice-curvature}
\lean{SmoothPoincare.physicalRadialVolume_eq_map_scaleMetric, SmoothPoincare.physicalAnchor_and_anchoredBuffer_fullCurvature_tendsto_zero, SmoothPoincare.shiftedPhysicalEuclideanShellVolume_strict_decrease, SmoothPoincare.exists_orientedMovingFrame, SmoothPoincare.movingCurvatureSlots_tensor, SmoothPoincare.tendstoUniformlyOn_movingChiralProjector}
\leanok
On a supplied wide chart, radial approximation on the closed buffered
band gives a physical anchor and transfers the curvature ceiling to
that band. Already in Euclidean space a positive shift of the radial
coordinate gives a strictly decreasing normalized adjacent-shell volume
quotient. On one oriented coordinate basis, positive metric matrices
admit a continuous choice of compatible orthonormal frames. On a fixed
uniform ellipticity window uniform convergence of the metric coefficients
gives uniform convergence of the density-weighted chiral projections.
\end{proposition}
\begin{proof}
\leanok
Push forward the restricted physical volume by the normalized physical
distance; retain the omitted region when comparing whole shells.
For the shifted quotient apply the mean value theorem to differences of
fourth powers. On the compact matrix window apply uniform continuity
to orthonormalization and the four tensor slots in this same basis.
\end{proof}

\begin{proposition}[Positive marked caps and their individual costs]
\label{prop:ed03-cap-tests}
\uses{prop:ed03-inclusive-caps,lem:ed03-splice-curvature}
\lean{SmoothPoincare.HasJoinedCapTests.yamabeFloor, SmoothPoincare.twoCapSigns_of_joinedTests, SmoothPoincare.HasJoinedCapTests.markedSide, SmoothPoincare.cappedEnergy_le_raw_add_halfCollarError}
\leanok
Set $T=4\pi^2$. On one connected joined path, assume the continuous sign
test, positive start, Euler bound and energy ceiling $2T-\gamma$ for
$\gamma>0$. If every smooth log-conformal test on each completed cap
has joined transplants with both numerator and volume converging,
each cap has Yamabe constant at least $\sqrt{48\gamma}$. The actual
cap metric, raw side and complementary marked core then define a positive
marked side. On that cap, strip mass $s$, norm error $\beta$, and a comparison
of the entire output complement give a cap cost at most
$E_{\mathrm{raw}}+2\sqrt{s}\,\beta+\beta^2$.
\end{proposition}
\begin{proof}
\leanok
Pass the joined lower bound to each test quotient before taking the
infimum. The limiting test volume is positive. For the cost, apply the
$L^2$ triangle inequality on the trim and add the complementary integral.
The two caps use their own half-collar errors. Once both signs and markings
are available, the sum test in Proposition~\ref{prop:ed03-inclusive-caps}
includes equality.
\end{proof}

\begin{proposition}[The selected disk widths and finite band charge]
\label{prop:ed03-same-width}
\uses{prop:ed03-endpoint-windows,prop:ed03-cap-tests}
\lean{SmoothPoincare.endpointPhysical_omitted_charges, SmoothPoincare.exists_selectedDiskWidths_small_chiral_sum, SmoothPoincare.endpointAllowance_on_selectedChartDisks, SmoothPoincare.bandPhysical_omitted_charges, SmoothPoincare.bandPhysical_yamabeEuler_sandwich}
\leanok
Use the volume-one physical Yamabe sequence with the ordered spectral,
Euler and signature data and the modified probability gauge with its
constant test and same-region spectral transport. At $I=T$ and
$e_i\to T$, suppose marked disjoint chart disks capture all but vanishing
modified volume. Their full physical exterior has chiral charge $r_i\to0$.
There is one $k_i\ge i$ per finite family for which the sum of physical
trim charges tends to zero. Write $s_i$ for their union charge at width
$1/(k_i+1)$. For any nonnegative allowance $\eta_i$ on these widths,
suppose that eventually the strict joined ceiling
$e_i+2\sqrt{s_i}\eta_i+\eta_i^2<2T$ yields a positive competitor on
the original type costing at most $(\sqrt{r_i}+\eta_i)^2$. Then for each
$\delta>0$, eventually $\eta_i>2\pi-\delta$.

At a single index, if $U$ is a supplied retained set with omitted modified
probability at most $d$, and $\ell^2$ is the total modified spectral second
moment with $\ell\ge0$, then each physical chiral charge satisfies
\[
 E_\pm(g;U^c)\le 2e-\tfrac32\max\{Y/6-\ell\sqrt d,0\}^2 .
\]
\end{proposition}
\begin{proof}
\leanok
The endpoint scalar test and spectral deficit transfer vanishing omitted
volume to vanishing physical charge. The marked coordinate spheres have
zero physical volume. Thin their collars together at each finite index
and apply the repair test at those widths. A fixed saving below $2\pi$
would violate the fixed-type infimum. At a single index Cauchy--Schwarz
bounds the omitted first moment by $\ell\sqrt d$. Transport the captured
second moment on $U$ and use the common total chiral energy of either
sign. The auxiliary positive-Ricci metric supplies the disk selection in
the manuscript; the charges here are measured in $g$.
\end{proof}

\section{The selected middle band and signed disk repairs}
\label{sec:ed03-marked-reduction}

\begin{theorem}[The ordered height and its middle attachment]
\label{thm:ed03-selected-middle-attachment}
\uses{thm:ed03-morse-threshold}
\lean{SmoothPoincare.exists_orderedMorsePresentation, SmoothPoincare.OrderedMorsePresentation.exists_middleIndexCuts, SmoothPoincare.exists_orderedMiddleBand_of_inducedEnergy_le, SmoothPoincare.MiddleIndexCuts.regular_transport, SmoothPoincare.exists_isolatedMiddleAttachment, SmoothPoincare.exists_selectedHandleReduction_of_inducedEnergy_le}
\leanok
For a compact oriented isometric hypersurface as in
Theorem~\ref{thm:ed03-morse-threshold}, suppose $E_+\le 64\pi^2/9$.
There is a direction $u$ and a smooth auxiliary Morse function $h$
on the same domain, with the critical points and indices of the
height in direction $u$. Its values at critical points are distinct
and strictly ordered between index classes. There are regular cuts
$a<b$ such that the critical points in the closed band are exactly
the original height's index-two points. If that set is empty, an
ambient diffeomorphism carries the lower sublevel and level to
their upper counterparts. Otherwise it is a singleton $\{p\}$.
A smooth attachment of $D^2\times D^2$ along its parametrized
$S^1\times D^2$ face identifies the local upper sublevel at
$h(p)+\varepsilon$. Regular ambient diffeomorphisms transport both
levels and sublevels between the two local cuts and the outer cuts.
\end{theorem}
\begin{proof}
\leanok
Enclose the compact range in one open Morse strip. Change the
normed model by a continuous linear equivalence, transport the
Hessian inertia at every critical point, rearrange by index, and
separate the remaining ties. The resulting finite critical-value
sets admit strict separators on either side of index two. In the
single-point case shrink a Morse chart inside the band, leaving room
for both regular cuts and its attaching tube. The smooth adjunction
identifies the lower sublevel and handle by their native inclusions
and fixes the deep lower part. Apply the regular flow on either side.
\end{proof}

\begin{proposition}[End counts and marked reassembly]
\label{prop:ed03-end-counts}
\uses{thm:ed03-selected-middle-attachment}
\lean{SmoothPoincare.oddBalance_of_morse_on_homotopySphere, SmoothPoincare.morseCriticalCount_congr, SmoothPoincare.alternating_morseCriticalCounts_eq_two, SmoothPoincare.morseCriticalCount_reducedOddCases, SmoothPoincare.morseCriticalCount_neg, SmoothPoincare.morseCriticalCount_zero_le_one_add_one_of_homologyZero, SmoothPoincare.morseCriticalCount_four_le_three_add_one_of_dual_homologyZero, SmoothPoincare.transportedMiddleTraceHomeomorph, SmoothPoincare.transportedMorsePresentationHomeomorph, SmoothPoincare.transportedOutgoingDiffeomorph}
\leanok
On a compact smooth four-manifold homotopy equivalent to the unit
four-sphere, let $n_j$ count the actual critical points of index $j$
of a Morse function. Rearrangement preserves those counts even when
critical values coincide. Then
$(n_1-n_0+1)+(n_3-n_4+1)=n_2$ as an integer identity. With
$n_0\le n_1+1$ and $n_4\le n_3+1$, both summands vanish in the
zero case. In the one case exactly one is one. The two inequalities
follow from degree-one boundary maps on the actual critical generators
of $h$ and $-h$ whose cokernels compute singular $H_0(M;\mathbb Q)$.

For supplied end-body homeomorphisms and compatible boundary maps,
changing the lower marking carries every framed tube by the inverse
boundary change. The trace homeomorphism fixes each handle coordinate.
It combines with the actual upper-body change to give a homeomorphism
of the closed adjunctions. Smooth end markings give a smooth boundary
gluing.
\end{proposition}
\begin{proof}
\leanok
The rational Euler characteristic is two by homotopy invariance.
Group its finite Morse sum by the five possible indices. Function
reversal complements every index. Rank-nullity and the rank-one
zeroth homology then give the bounds from the supplied boundary maps. For the
gluing, descend the maps on the two pieces through the attaching
relations, and descend their inverses through the new relations.
A trace change $s$ transports its outgoing gluing to
$(\partial\beta)^{-1}q^{-1}\eta\,\partial s$.
\end{proof}

\begin{proposition}[The standard product seam]
\label{prop:ed03-smooth-product-seam}
\lean{SmoothPoincare.StandardSurgeryModel.pieces_agree_on_face, SmoothPoincare.StandardSurgeryModel.leftRegion_union_rightRegion, SmoothPoincare.StandardSurgeryModel.leftRegion_inter_rightRegion, SmoothPoincare.StandardSurgeryModel.seamChart, SmoothPoincare.StandardSurgeryModel.seamChart_at_zero, SmoothPoincare.StandardSurgeryModel.seamChart_sides, SmoothPoincare.diffeomorphOfSmoothMarkedSeam, SmoothPoincare.diffeomorphToS4OfSmoothMarkedSeam}
\leanok
In $S^4\subset\mathbb R^2\oplus\mathbb R^3$, the two product pieces
$D^2\times S^2$ and $S^1\times D^3$ meet along $S^1\times S^2$.
In signed coordinate $t=|x|^2-1/2$ a smooth chart across their common
face is
\[
 (u,v,t)\longmapsto (\sqrt{1/2+t}\,u,\sqrt{1/2-t}\,v).
\]
A supplied homeomorphism to this model is a diffeomorphism if it
and its inverse are smooth on the two strict interiors and it agrees
on a full signed source collar with this chart composed with a smooth
boundary diffeomorphism times the identity in $t$. An orthogonal
identification of the two ambient models gives the unit sphere in
$\mathbb R^5$ as the target.
\end{proposition}
\begin{proof}
\leanok
Normalize each nonzero block to recover its unit direction. The
inverse signed parameter and the normalized blocks are smooth on the
overlap. These formulas give the inverse on the entire marked strip.
Together with the two interiors it covers the source and target.
\end{proof}

\begin{proposition}[A fixed family of comparison disks]
\label{prop:ed03-comparison-disks}
\uses{prop:ed03-same-width}
\lean{SmoothPoincare.exists_selectedChartDisk_of_compact_subset, SmoothPoincare.eventually_exists_selectedComparisonDisks, SmoothPoincare.eventually_exists_markedBandCharge, SmoothPoincare.selectedDiskTrim_eq_full_image, SmoothPoincare.eventually_fullDiskTrims_small_both_charges}
\leanok
Fix the E58 spectral, Euler and same-set transport data for the physical
metrics $g_i$ and modified probabilities $\mu_i$. If $e_i\to I<5T/4$,
choose the omitted-probability tolerance first. Suppose a finite family of
fixed radii, with late smooth open comparison-ball charts, satisfies the
strict separation, omitted-probability and curvature margins. The comparison
metric measures the ball radii, while $g_i$ measures full curvature on the
outer balls. On late sources there are disjoint smooth closed disk markings
with these ceilings, and both physical chiral energies are strictly below
$T$ on the whole complement of their union.

For any one finite marked family with integrable chiral densities, all
sufficiently thin common trims are full radial annuli. The sum of their
two physical chiral charges is arbitrarily small.
\end{proposition}
\begin{proof}
\leanok
The preimage of each compact inner ball lies inside an open model
ball. Enclose it in a strictly smaller concentric ball and dilate to
unit radius. This preserves the outer comparison and retains a buffer.
Apply the E58 estimate to that chosen disk union. A compact unit disk
has a closed buffer in the open chart. Its boundary sphere has zero
physical volume, so both trim charges vanish by continuity from above.
Take the finite intersection of the resulting eventual width conditions.
\end{proof}

\begin{proposition}[The full exterior of a balanced marked disk]
\label{prop:ed03-balanced-disk}
\lean{SmoothPoincare.interior_selectedChartDisk, SmoothPoincare.signedCharge_fullMarkedExterior, SmoothPoincare.eventually_balancedMarkedExterior_below_threshold, SmoothPoincare.balancedMarkedDisk_saving_gt, SmoothPoincare.exists_dyadicDepth_for_recognitionGap, SmoothPoincare.eventually_balancedMarkedDisk_saving_gt_gap}
\leanok
Let $Q_i\subset\operatorname{int}B_i$ be actual marked disks.
Suppose the exterior of $B_i$ carries $q=2^{-k}$ times the charge
of the exterior of $Q_i$, to error $\tau_i\to0$ in both physical
chiralities. If $k\ge1$ and $e_i\to I<2T$, the full closed exterior
of $B_i$ has both charges strictly below $T$ eventually.
At one finite source $g$, suppose a joined endpoint $h$ has energy
below $2T$, and a positive cap on the original smooth type has
energy at most $E_s(h;M\setminus\operatorname{int}B)+\delta$ in
each sign $s$. Then
\[
 E_s(g;B)-E_s(h;B)>(1-q)e+qE_s(g;Q)+I^+(M)-2T-\tau-\delta .
\]
For a given $\varepsilon_*>0$ choose one $k$ with
$q(T+\varepsilon_*)\le\varepsilon_*$. On a fixed type with
$I^+(M)>T+\varepsilon_*$, if $e_i\to I^+(M)$, the endpoint and cap
conditions hold eventually and $\tau_i+\delta_i\to0$, both required
disk-side savings eventually exceed $\varepsilon_*$.
\end{proposition}
\begin{proof}
\leanok
The marked sphere has zero mass for each physical chiral measure,
so the charge on the closed exterior is the full energy minus the
closed-disk charge. Since $q\le1/2$, the raw residual is bounded
by $e_i/2+\tau_i$. At the finite endpoint, the actual positive cap
has energy at least the same-type infimum. Subtract this bound and
the measured division error from the two exact disk decompositions.
The limiting lower bound $(2-q)I^+(M)-2T$ strictly exceeds the gap
for the chosen depth, independently of the smooth type.
\end{proof}

\begin{proposition}[Signed gain of a marked repair]
\label{prop:ed03-signed-repair}
\uses{prop:ed03-cap-tests,prop:ed03-comparison-disks}
\lean{SmoothPoincare.normSq_sub_le_signedJoinedCollarCost, SmoothPoincare.SignedCollarAssembly.Estimates.ofWeightedDensities, SmoothPoincare.SignedCollarAssembly.joinedEnergy_le, SmoothPoincare.SignedCollarAssembly.outputEnergy_le, SmoothPoincare.SignedCollarCapTestPath.outputYamabeFloor_of_ceiling_lt, SmoothPoincare.excludesStrictSelectedDiskSignedCollarRepair_of_capTests}
\leanok
Fix disjoint marked patches, single-patch and simultaneous metrics on the
original type, and both complete complement comparisons. On the selected
width of disk $B_j$, write $S$ for raw annular charge, $H$ for the squared
norm of the leading cutoff-weighted endpoint density, and $\beta$ for
its measured error. For either chirality put
\[
 U_j=\max\{0,(\sqrt H+\beta)^2-S\},\qquad
 V_j=(\sqrt{H_{\rm out}}+\beta_{\rm out})^2-S_{\rm out}.
\]
With the affine-in-parameter error bound and signature equality on the
single-patch metrics, the joined energy is at most
$C=e+\sum_j\min_{\pm}U_j^{\pm}$. The capped energy is at most $e-G$,
where $G=\sum_j\max_{\pm}(\lambda^{\pm}(B_j)-V_j^{\pm})$.
Suppose $C<2T$. On a connected test family starting at a positive
source, with a continuous Yamabe sign detector, the Euler-square
inequality, energy bounded by the joined path, and scalar and volume
test transplants to the capped output, the latter has Yamabe constant
at least $\sqrt{48(2T-C)}>0$. Hence $G\le e-I^+(M)$.
\end{proposition}
\begin{proof}
\leanok
The triangle inequality bounds the joined density by interpolation
between the two endpoint bounds. A single patch has the same total
increment in each sign. Equality on its complete complement localizes
this identity, permitting a sign choice separately on each patch.
The continuous sign level cannot cross zero below the wall. Pass
the resulting uniform lower bound through the transplanted Yamabe
quotients and compare that finite output with the same-type infimum.
\end{proof}

\begin{proposition}[A measured choice of exact collar widths]
\label{prop:ed03-measured-choices}
\uses{prop:ed03-signed-repair}
\lean{SmoothPoincare.exists_measuredCollarChoice, SmoothPoincare.signedCollarCosts_le_measuredScore, SmoothPoincare.excludesStrictRepair_of_measuredCollarChoice, SmoothPoincare.excludesStrictRepair_of_measuredSelectedDiskCollarChoice}
\leanok
Suppose each protected core $C_j$ has $N_j>0$ available marked
charts and widths with row budget $D_j\ge\sum_k(w_{jk}^2+\beta_{jk}^2)$.
Choose a minimizer in each row. Suppose the choices admit the
simultaneous marked assembly and cap tests of
Proposition~\ref{prop:ed03-signed-repair}. On the chosen annuli
assume the measured weighted densities and allowances satisfy the
local curvature and error bounds for both signs. Set
$\Gamma=2\sum_j D_j/N_j$. A positive source obeying
$e+\Gamma<2T$ must then satisfy
\[
 \sum_j\max_{s=\pm}\lambda^s(C_j)\le \Gamma+e-I^+(M).
\]
\end{proposition}
\begin{proof}
\leanok
At least one score in each finite row is at most its average.
On the full annulus and its exterior half, the surcharge is bounded
by $2\sqrt{S}\,\beta+\beta^2\le S+2\beta^2\le2(w^2+\beta^2)$.
Protected charge is retained on each chosen disk. Summing on the
same family gives the joined ceiling and the output comparison.
\end{proof}

\begin{proposition}[Physical defects at the isotropic-gap entry]
\label{prop:ed03-isotropic-entry}
\lean{SmoothPoincare.orderedEndpointSpectrum_hilbertSchmidt_bounds, SmoothPoincare.endpointPhysicalIsotropicGapBounds, SmoothPoincare.IsotropicGapBounds.quantity_sandwich, SmoothPoincare.isotropicGap_smoothingEntry, SmoothPoincare.endpointPhysicalIsotropicGap_smoothingEntry}
\leanok
For each volume-one positive Yamabe metric with the endpoint Euler,
signature, ordered eigendata and strict spectral inequality, let $e$
be its common chiral energy, $Y$
its class Yamabe constant, $D$ its integrated Ricci defect, and
$\sigma$ the integrated square of the larger top Weyl eigenvalue of
the two signs. Set $P=4e+3\sigma$ and $T=4\pi^2$. Then
$9e/2\le P\le8e$. Along a sequence of these data, even on varying
smooth types, $P\to8T$ implies $D\to0$,
$36\sigma-Y^2\to0$, and eventually $e<2T-T/9$.
\end{proposition}
\begin{proof}
\leanok
The ordered trace-free spectra bound the top eigenvalue square on
each sign. Integrate and apply signature. The physical scalar-square
and Euler identities give $P-8T=D+(36\sigma-Y^2)/12$, with both
terms nonnegative. Squeezing them separately gives the two limits.
The lower bound for $P$ supplies the uniform energy margin.
\end{proof}

\section{The marked end filtrations and four-measure division}
\label{sec:ed03-end-bands-division}

\begin{theorem}[Both ends of the selected height]
\label{thm:ed03-end-filtrations}
\uses{thm:ed03-selected-middle-attachment}
\lean{SmoothPoincare.EndHandleFiltration, SmoothPoincare.exists_orderedEndHandleFiltrations, SmoothPoincare.orderedCriticalPartition, SmoothPoincare.exists_selectedEndHandlePresentation_of_inducedEnergy_le, SmoothPoincare.IsolatedZeroHandle.birthHomeomorph, SmoothPoincare.IsolatedOneHandle.attaching_radius_two_lt, SmoothPoincare.IsolatedOneHandle.attachingFeet_smooth, SmoothPoincare.IsolatedOneHandle.attachingFeet_isClosedEmbedding}
\leanok
At $E_+\le64\pi^2/9$ the same selected direction, ordered function $h$
and regular cuts $a<b$ admit finite zero- and one-handle filtrations from
the empty set up to $\{h\le a\}$ and $\{h\ge b\}$. The latter uses $-h$.
The two lists and the original middle set partition the actual critical
points of the height. The middle alternative retains its regular transport
or its rounded two-handle with both adjacent transports.
\end{theorem}
\begin{proof}
\leanok
List the critical points below the cut by strictly increasing value and
insert regular separators. At each point the local Morse trace supplies
a birth with empty attaching face or a smooth attachment on its two
parametrized feet. The chosen chart retains its normal form up to twice
its recorded radius $R$, and $\sqrt{2\varepsilon+8r^2}<2R$. This
leaves room for the entire radius-two foot graph. The disjoint-sum description at a birth is topological.
Regular bands connect the local cuts to the separators. Apply the same
construction to the reversed function and use regularity of the selected
levels to separate the three point sets.
\end{proof}

\begin{proposition}[Regular markings and a normalized tree edge]
\label{prop:ed03-end-markings}
\uses{thm:ed03-end-filtrations,prop:ed03-end-counts}
\lean{SmoothPoincare.regularEndMarked_of_image_eq, SmoothPoincare.RegularEndMarked.exists_tubeTransport, SmoothPoincare.exists_regularBandEndMarked, SmoothPoincare.HasRegularLevelTransport.exists_nativeMarked, SmoothPoincare.nativeOppositeLevelDiffeomorph, SmoothPoincare.exists_leafBallHomeomorph, SmoothPoincare.exists_roundLeafAdjunction, SmoothPoincare.exists_leafBallTransport}
\leanok
A regular ambient flow taking both the sublevel and level to their
targets restricts to smooth diffeomorphisms of the native end body and
its boundary. Their forward and inverse maps commute with the two
inclusions and transport any supplied boundary parametrization.
This applies to the actual regular comparisons on the selected ends by
conjugating the ambient diffeomorphism through the identity change of model.
The native level maps for $h=b$ and $-h=-b$ also agree on their
underlying points in both directions. For two normalized product zero-handles joined along opposite standard feet
of one product one-handle, a homeomorphism to one ball identifies the
exposed boundary. It transports the remaining attaching and outgoing maps in
the adjunction.
\end{proposition}
\begin{proof}
\leanok
Restrict the ambient map and its inverse using the sublevel and level
embedding charts. In the normalized leaf use an affine axial coordinate on
each of the three cylinders. Their maps agree on the two entire feet, and
the inverse descends across the joins. Radial projection takes the boundary of
the resulting convex cylinder to the round sphere. Functoriality of the
quotient transports the other markings.
\end{proof}

\begin{theorem}[Four measures on smooth marked disks]
\label{thm:ed03-atlas-division}
\lean{SmoothPoincare.exists_bisecting_measured_subfamily, SmoothPoincare.exists_four_measure_ball_division, SmoothPoincare.exists_fine_atlas_disk_packing, SmoothPoincare.exists_four_measure_atlas_division, SmoothPoincare.exists_bisecting_marked_disks}
\leanok
On a compact Hausdorff smooth boundaryless four-manifold, let $U$ be
open and let $\mu_1,\ldots,\mu_4$ be finite Borel measures vanishing on
singletons. For every $\varepsilon>0$ there is a finite disjoint family
of smooth marked closed disks in $U$ whose union $V$ satisfies
$|\mu_j(V)-\mu_j(U)/2|<\varepsilon$ for all four measures.
\end{theorem}
\begin{proof}
\leanok
Pack the open remainder of each successive chart with disjoint closed
balls of small mass for $\sum_j\mu_j$, losing arbitrarily little mass.
A finite chart cover suffices. For each packed disk form its vector of
four masses and choose its sign so the cross term with the previous
signed sum is nonpositive. The squared final discrepancy is bounded by
the sum of the squared lengths. The small individual masses bound that
sum, while the uncovered set accounts for the error in replacing the
packing mass by the mass of $U$.
\end{proof}

\begin{proposition}[A single exterior and its protected collars]
\label{prop:ed03-one-disk-separator}
\uses{thm:ed03-atlas-division}
\lean{SmoothPoincare.markedDisk_seam_null, SmoothPoincare.exists_marked_four_measure_half_of_separation, SmoothPoincare.exists_finitePointRouting_into_opens, SmoothPoincare.exists_diffeomorph_contracting_finiteMarkedDisks, SmoothPoincare.exists_finite_markedDisk_separator, SmoothPoincare.exists_nestedMarkedDisk_fourMeasure_dyadic, SmoothPoincare.exists_bufferedNestedMarkedDisk_fourMeasure, SmoothPoincare.exists_balancedPhysicalExteriorDivision, SmoothPoincare.exists_balancedPhysicalSpectralExteriorDivision, SmoothPoincare.exists_eventually_balancedPhysicalExterior_below_threshold, SmoothPoincare.halfDivision_of_packing_trace, SmoothPoincare.exists_iteratedMarkedDiskDivision, SmoothPoincare.NestedBalancedDisks.eventually_trim_disjoint_core, SmoothPoincare.eventually_protectedFourMeasureTrim}
\leanok
On a connected source, take four finite measures absolutely
continuous with respect to physical volume. There is a single smooth
marked outer disk enclosing the core strictly whose full exterior
simultaneously approximates half of every original exterior mass to
arbitrary accuracy. With error $\tau$ at each step, after $k>0$ steps the error
relative to the fraction $2^{-k}$ is bounded by $2\tau(1-2^{-k})$.
For a nested pair, every sufficiently thin outer trim lies in a full
chart annulus, avoids the entire original core, and has arbitrarily
small total incoming mass for the four measures. Apply this on each
physical source to its two chiral measures, a finite spectral density
and the full curvature measure. At any fixed positive depth, if
$e_i\to I<2T$ and the errors tend to zero, both raw exterior chiral
charges eventually fall strictly below $T$. Here $T=4\pi^2$.
\end{proposition}
\begin{proof}
\leanok
Include the original core in the finite list. Use supported smooth
point motions in complements of finitely many points to route their
centers into the desired open sides of a seed sphere. Before this
motion contract each full disk into the inverse image of its assigned
open side, fixing every other disk. Pulling the chart back produces a
single outer disk on the original physical source. The core seam has
volume zero. On the packing the exterior trace is exact. Off the packing its discrepancy costs at most half of the
uncovered mass in addition to the packed discrepancy. Iterate the
half-mass recurrence. The inverse radial coordinate on the compact
core has maximum strictly below one. Intersect this uniform gap with
the open chart buffer and the four null-seam estimates.
\end{proof}

\begin{proposition}[Halfspaces in a supplied physical chart]
\label{prop:ed03-marked-chart-division}
\uses{thm:ed03-atlas-division}
\lean{SmoothPoincare.fourMeasureImbalance_continuous, SmoothPoincare.fourMeasureImbalance_odd, SmoothPoincare.exists_fourMeasureBall_bisection, SmoothPoincare.markedDiskPullback_absolutelyContinuous, SmoothPoincare.exists_markedDisk_affine_bisection, SmoothPoincare.exists_markedDisk_finiteBall_division}
\leanok
In a supplied smooth marked disk, pullback takes its four finite
volume-absolutely-continuous measures to Lebesgue-absolutely-continuous
measures. Given the antipodal-zero principle for continuous odd maps
from the four-sphere to $\mathbb R^4$, there is an affine hyperplane
bisecting all four restricted masses exactly. The finite-ball
approximation in the chart holds unconditionally.
\end{proposition}
\begin{proof}
\leanok
Transport the measures by the inverse chart, keeping their full
restrictions to the marked disk. Hyperplanes for a nonzero normal have
zero mass. Dominated convergence gives continuity of the signed
halfspace imbalance, including at the poles of the parameter sphere.
Opposite parameters reverse the imbalance. A pole can have zero
imbalance only when all four masses vanish, in which case a central
plane works.
\end{proof}

\begin{proposition}[Cost on the selected widths and the matched endpoint]
\label{prop:ed03-measured-endpoint-wall}
\uses{prop:ed03-one-disk-separator}
\lean{SmoothPoincare.measuredCollarAllowance_sq_le, SmoothPoincare.measuredCollarRowBudget_of_localAllowances, SmoothPoincare.SignedCollarCandidateControls.ofCutoffCharges, SmoothPoincare.excludesStrictRepair_of_localMeasuredSelectedDiskCollars, SmoothPoincare.balancedEndpointEnergy_lower_of_localChange, SmoothPoincare.eventually_balancedEndpointEnergy_gt_wall_on_selectedTrim}
\leanok
On the actual candidate annuli write $b=w+\tau+d$.
For nonnegative local data with $a\le\min(A,1)$ and
$0\le\beta\le C_0(ab+\sqrt{db}+d)$, the squared error is at most
$12 C_0^2[A^2(w^2+\tau^2)+d(w+\tau)+d^2]$.
The row sum controls the finite measured width selection. For the
balanced disk $B$ around a protected core $Q$, put $q=2^{-k}$.
Fix a chirality $s$. Suppose the initial exterior division has upper
error $\tau$, the matched same-type cap is positive and has signed
energy at most the endpoint exterior charge plus $\delta$, and the
initial and endpoint charges agree on $B\setminus W$. Then
\[
 E_+(h)\ \ge\ I+(1-q)E_+(g)+q\lambda_g^s(Q)
          -\tau-\delta-\lambda_g^s(B\cap W).
\]
Along a sequence tending to $I>T+\varepsilon$, if
$q(T+\varepsilon)\le\varepsilon$ and the three losses tend to zero,
every such endpoint eventually exceeds $2T+\varepsilon$.
\end{proposition}
\begin{proof}
\leanok
Square the local allowance and apply finite Cauchy--Schwarz to its two
mixed sums. For the endpoint, the positive cap has energy at least the
fixed-type infimum. Split each full signed measure on $B$ and its
complement, using the null seam and the agreement off $W$. Drop the
nonnegative endpoint charge on $B\cap W$. The limiting lower bound
is at least $(2-q)I>2T+\varepsilon$.
\end{proof}

The selected filtrations still require smooth normalization of actual feet,
with the other tubes and the outgoing signed collar fixed. The one-disk
division now supplies the outer marking and full raw complements. For a successful
repair, the physical jet bounds, matched positive cap and simultaneous signed
saving must hold for one path on its selected widths. Positive caps used in
the last proposition are finite attained comparisons on the original type;
$I$ denotes its limiting infimum.

\section{Native handle feet and variations on the selected physical disk}
\label{sec:ed03-native-feet-jets}

\begin{theorem}[Parametrizations on the actual incoming level]
\label{thm:ed03-native-foot-charts}
\uses{thm:ed03-end-filtrations}
\lean{SmoothPoincare.IsolatedOneHandle.footPartialDiffeomorph, SmoothPoincare.IsolatedOneHandle.footPartialDiffeomorph_apply_nativeAttaching, SmoothPoincare.IsolatedOneHandle.footPartialDiffeomorph_disjoint, SmoothPoincare.IsolatedOneHandle.footComponentChart_closedBall_two_subset_source, SmoothPoincare.IsolatedOneHandle.footComponentPoint_val, SmoothPoincare.IsolatedOneHandle.footComponentChart_symm_point}
\leanok
For each of the two feet of an actual native one-handle in either selected
end filtration, its full unit disk extends to a smooth chart containing the
closed radius-two ball. The two images lie on disjoint branches of the
incoming level. Each chart restricts to the open connected component of
its own foot center. Its inverse recovers the transverse parameter.
\end{theorem}
\begin{proof}
\leanok
In Morse coordinates the transverse coordinate is $rz$ and the signed
axis coordinate has magnitude $\sqrt{2\varepsilon+r^2|z|^2}$. The squared
model radius is $2\varepsilon+2r^2|z|^2$. The retained double-radius
normal form therefore covers the graphs for $|z|\le2$. Division of the
transverse coordinate by $r$ gives the smooth inverse. The image of the
connected closed ball lies in the component of its center.
\end{proof}

\begin{proposition}[Relative oriented feet in a signed collar]
\label{prop:ed03-relative-foot-collar}
\uses{thm:ed03-native-foot-charts}
\lean{SmoothPoincare.IsolatedOneHandle.exists_orientedNativeFootComponentChart, SmoothPoincare.protectedFootBalls_compact_and_compl_preconnected, SmoothPoincare.exists_bufferedFootSignedCollarComparison, SmoothPoincare.SupportedLevelIsotopy.signedCollarComparison_protected, SmoothPoincare.SupportedLevelIsotopy.signedCollarComparison_end}
\leanok
Give the closed native component an orientation. Its complete foot chart
has an oriented representative after either the identity or the central
reflection of the three parameters. A finite pairwise disjoint family of
protected closed chart balls has compact union and preconnected complement.
For a supplied oriented model foot in the same connected level, disjoint
from these protected balls, the supported isotopy matches every foot
parameter. On supplied full signed collar charts it gives a smooth
comparison in both directions. On the central plateau its foot parameter
is the chosen linear isometry; on the protected tubes and both outer
regions it preserves every coordinate under the two markings.
\end{proposition}
\begin{proof}
\leanok
Restrict the actual chart to a connected open source and use the
orientation dichotomy there. Reflection corrects the negative case on
the entire marking. Apply relative ball isotopy off the union of the
protected balls. In signed coordinates the comparison and inverse are
$(x,t)\mapsto(J_{\rho(t)}x,t)$ and $(x,t)\mapsto(J_{\rho(t)}^{-1}x,t)$,
conjugated by the given charts. Choose $\rho=1$ on the central plateau
and $\rho=0$ near both outer ends.
\end{proof}

\begin{theorem}[Calibration by the actual height]
\label{thm:ed03-native-height-collar}
\uses{thm:ed03-end-filtrations}
\lean{SmoothPoincare.exists_nativeSignedHeightCollar, SmoothPoincare.exists_nativeOppositeSignedCollars}
\leanok
On a nonempty regular level $L=h^{-1}(c)$ of a smooth height on a compact
boundaryless source, a smooth partial diffeomorphism with source
$L\times(-r,r)$ has target exactly $h^{-1}((c-r,c+r))$. Its zero slice is
the native inclusion. The inverse signed coordinate at $y$ is $h(y)-c$.
The opposite height has a collar with the same physical image; the
identification is the identity on level points and negates the transverse
coordinate in both directions.
\end{theorem}
\begin{proof}
\leanok
Compactness separates the selected level from the critical values.
Use the smooth gradient-like field with unit negative height speed
throughout a smaller closed regular band. The maps are its flow for time
$-t$ from $x\in L$ and its return flow for time $h(y)-c$ from $y$. The
flow law supplies the two inverse identities. Smooth factorization through
the native level gives smoothness of the inverse projection. Conjugate
by the native opposite-level map and $t\mapsto-t$ to obtain the second chart.
\end{proof}

\begin{proposition}[Physical cutoff jets on the balanced disk]
\label{prop:ed03-physical-marked-jets}
\uses{thm:ed03-atlas-division}
\lean{SmoothPoincare.exists_balancedPhysicalSpectralExteriorDivision_of_endpointData, SmoothPoincare.exists_selectedMarkedRadialCutoff, SmoothPoincare.exists_selectedRadialCutoff_radial_deriv, SmoothPoincare.fderiv_fderiv_selectedRadialCutoff_apply, SmoothPoincare.exists_selectedMarkedMetricLift, SmoothPoincare.selectedMarkedMetricLift_chartGram, SmoothPoincare.exists_selectedRadialSpliceSecondPartialError_on_physicalTrim}
\leanok
Choose one outer marked disk for the two physical chiral measures, the
tested top-eigenvalue measure in the same physical frame, and the full
curvature measure. On a buffered width $w=1/(n+1)$ the physical metric
has a smooth model lift agreeing with its marked pullback on a neighborhood
of the entire closed buffer. The squared-radius step equals one in the
inner region and zero in the outer region. Along every unit ray its
derivative attains $-1/(2w)$ somewhere in the full trim.
For two smooth coefficients $G,B$, write $H=B-G$. In the pullback of the
actual incoming trim volume, the second-partial cutoff error, at path
parameter $0\le t\le1$, has $L^2$ norm bounded by
\[
 t\left[\left(4C_2/w^2+C_1/w\right)\|H\|_2
   +(2C_1/w)\left(\|\partial_m H\|_2+\|\partial_i H\|_2\right)\right].
\]
\end{proposition}
\begin{proof}
\leanok
Transport the restricted volume by the inverse marking. Differentiate
the smooth profile of squared radius twice. The product rule gives one
Hessian term multiplying $H$ and the two displayed mixed terms. The
compact buffer admits a single smooth extension of the pullback metric;
its fixed chart Gram coefficients agree throughout an open neighborhood.
\end{proof}

\begin{theorem}[Attained signed variation on the disk side]
\label{thm:ed03-selected-side-variation}
\uses{prop:ed03-physical-marked-jets}
\lean{SmoothPoincare.coreVariation_taylor_upper, SmoothPoincare.signedCoreSaving_eq_energySaving, SmoothPoincare.positiveCoreVariation_bounds, SmoothPoincare.exists_selectedDiskSideVariation}
\leanok
Put $U=B\cap W_n$ for the chosen outer disk and its marked trim. Test
a two-sided metric path with energy $f(s)=E_+(g_s)$, positive endpoints
at $\pm r$, and the same total charge on the entire complement $U^c$
as the source, separately in each sign. Suppose $f$ is $C^4$ and its
second and fourth derivatives are bounded in absolute value by $D$ on
$[-r,r]$. One of these same two positive endpoints realizes, in each
chirality, a saving on $U$ of at least
\[
 \max\{ |f'(0)|r-Dr^2/2,\;-f''(0)r^2/2-Dr^4/24\}.
\]
The saving is at most the smaller of $f(0)-I$ and the raw charge of $U$
in that sign. For a positive source this yields
$|f'(0)|\le (f(0)-I)/r+Dr/2$ and
$[-f''(0)]_+\le2(f(0)-I)/r^2+Dr^2/12$.
\end{theorem}
\begin{proof}
\leanok
Use Taylor's theorem on both signed intervals and choose the endpoint of
smaller energy. Averaging the cubic Taylor polynomials cancels both odd
terms. Signature and equality on the full complement identify each local
saving with the total drop. Positivity puts that attained endpoint above
the fixed-type infimum. The coordinate width $w$ and variation step $r$
are independent parameters.
\end{proof}

\begin{theorem}[Localized Ricci charge in the modified flow gauge]
\label{thm:ed03-modified-euler-transport}
\lean{SmoothPoincare.metricFrameTensorNormSq_eq_chiral_ricci_scalar, SmoothPoincare.modifiedCutoffMoments_full_eq, SmoothPoincare.modifiedEuler_timeWeight, SmoothPoincare.modifiedEuler_weightedTimeBound, SmoothPoincare.modifiedEuler_transport_of_flux, SmoothPoincare.modifiedTerminalRicci_lower_of_eulerTransport_of_measurableCharges}
\leanok
Let a common cutoff $0\le\chi\le1$ equal one on the initial charge set
and vanish off $B$ at terminal time $\tau>0$. Suppose the initial full
curvature charge is at least $\epsilon^2$. Write $m_+,m_-,m_S$ for
the two cutoff Weyl-square moments and the cutoff scalar-square moment at
time zero. Assume the integrability and measurability of the tested
densities. Given the spatial Euler transgression estimate with coefficient
$K$, a nonnegative remaining time factor of $L^2$ norm at most $L$, and
$b<1/2$, the actual terminal metric satisfies
\[
 \int_B |Z_\tau|^2\,dV_\tau\ \ge \epsilon^2/2-4(m_++m_-)-m_S/6-2KL\sqrt{\tau/(1-2b)}.
\]
\end{theorem}
\begin{proof}
\leanok
The Levi--Civita curvature symmetries and the first Bianchi identity
give the ordered-tensor norm $4(|W^+|^2+|W^-|^2)+2|Z|^2+R^2/6$.
Integrate with the same test on the initial slice. Cauchy--Schwarz in
time uses the exact integral $\tau/(1-2b)$ for the squared time
weight. At the terminal slice drop the nonnegative Weyl and scalar
moments and bound the cutoff Ricci integral by the whole supporting set.
\end{proof}

\section{Rounded edges and the cost of a flat disk reference}
\label{sec:ed03-rounded-edges}

\begin{theorem}[Birth bodies and edge incidence]
\label{thm:ed03-birth-edge-incidence}
\uses{thm:ed03-native-foot-charts}
\lean{SmoothPoincare.IsolatedZeroHandle.birthBodyChart_component, SmoothPoincare.IsolatedZeroHandle.exists_transportedBirthBody, SmoothPoincare.IsolatedOneHandle.traceComponentMap_fibers, SmoothPoincare.IsolatedOneHandle.levelFeet_disjoint_of_bodyFoot_ne, SmoothPoincare.IsolatedOneHandle.not_simplyConnected_of_same_bodyFoot, SmoothPoincare.IsolatedMiddleAttachment.traceComponentMap_bijective}
\leanok
At a selected birth at $p$, the radius-one closed four-ball in its rescaled
Morse chart is the whole sublevel component of $p$ at $f(p)+\varepsilon$.
The chart extends across its boundary and the ball and its marked sphere
transport together across the following regular band.
For a selected one-handle, the map on body components induced by the lower
insertion and the trace diffeomorphism is onto. Its only identifications are
the component labels of the two native feet. Distinct feet in the same body
preclude a simply connected outgoing body. The corresponding component map
for a selected two-handle is bijective.
\end{theorem}
\begin{proof}
\leanok
The birth chart has height $f(p)+\varepsilon |x|^2$. Its closed unit ball
is connected, compact and open relative to the selected sublevel. To compute
the component map, collapse the component labels of the attaching region
in the adjunction. The handle is connected. For a loop edge the joined-feet
obstruction applies to the actual trace. For an index-two attachment the circle
of normal-disk centers is connected, so it carries a single incoming label.
\end{proof}

\begin{theorem}[Rounded comparison at both native rims]
\label{thm:ed03-rounded-native-rims}
\uses{thm:ed03-birth-edge-incidence,thm:ed03-marked-edge-feet,thm:ed03-native-height-collar}
\lean{SmoothPoincare.IsolatedOneHandle.edgeBlockChart_nativeFoot, SmoothPoincare.IsolatedOneHandle.edgeBlockChart_afterFootNormalization, SmoothPoincare.IsolatedOneHandle.edgeBlockChart_signedOutgoing, SmoothPoincare.IsolatedOneHandle.roundedRimToOutgoing_signed_inverse, SmoothPoincare.IsolatedOneHandle.exists_roundedOutgoingCollar_nearNativeRims, SmoothPoincare.IsolatedOneHandle.exists_outgoingSignedBandTransport, SmoothPoincare.IsolatedOneHandle.exists_roundedRimTransport_toNextCut}
\leanok
The four-dimensional block chart restricts to both entire native attaching
disks with their parameters. Supplied foot normalizations retain their
individual parities. For a rounding on the Morse core with radius exceeding
$\sqrt{2\varepsilon+2r^2}$, its radial comparison $T$ has a smooth inverse
near both full rims. If $\Phi$ is the actual outgoing product chart obtained
in the construction, there are an open set $U$ containing both rim images and
$0<w\le\varepsilon/2$ on which the rounded collar satisfies
\[
 \Psi=T^{-1}\circ\Phi,\qquad
 \Psi^{-1}=\Phi^{-1}\circ T,\qquad
 f_{\rm rounded}(\Psi(z,t))=f(p)+t
 \quad(z\in U,\ |t|<w).
\]
Let $B_1$ be the collar at the next regular cut, $d$ the regular level
comparison, and $H$ the band comparison based on this same $\Phi$. On
a common positive substrip the maps satisfy
\[
 H(T(\Psi(z,t)))=B_1(dz,t),\qquad
 T^{-1}(H^{-1}(B_1(dz,t)))=\Psi(z,t).
\]
These equations retain every tube parameter supplied in the outgoing level,
and on the rim strip they retain the original label and both normal directions.
\end{theorem}
\begin{proof}
\leanok
At fixed axial square $s$, add $2\varepsilon+s-\operatorname{smoothCap}
(\varepsilon,r,\delta,s)$ to the squared transverse radius. Subtraction
recovers the starting point on the inverse domain. Both native rims lie
in the common open core. Compactness of $S^0\times S^2$ gives a uniform
two-sided interval for the pullback of the outgoing chart. At the next cut
conjugate the level transport times the identity by the two signed collars.
\end{proof}
A global edge contraction still needs the two incoming body markings, the
remaining tubes and compatible maps on their open overlaps with the rounded
block. A product extension chosen from two collars may shear tangentially
relative to another extension that agrees on the zero section.

\begin{theorem}[Realized reference on a selected width]
\label{thm:ed03-realized-disk-reference}
\uses{prop:ed03-physical-marked-jets,prop:ed03-selected-quadratic-jets}
\lean{SmoothPoincare.exists_selectedEuclideanReferenceSplice, SmoothPoincare.exists_selectedEuclideanSecondPartialError_on_physicalTrim, SmoothPoincare.selectedEuclideanSplice_flat_near_innerDisk, SmoothPoincare.eventually_balancedFlatReference, SmoothPoincare.selectedReference_endpoint_saving}
\leanok
Let $Q\subset\operatorname{int} B$ be the nested marked disks on one physical
source, and $W_n$ its selected trim. For every positive tolerance all sufficiently
large $n$ give a single width for which the four original trim masses have
sum below tolerance and $Q\cap W_n$ is empty. On that width one can splice
a smooth positive definite Euclidean reference in the original coordinates.
The time-one metric is flat on $B\setminus W_n$, including $Q$. All exterior
metric germs on $(B\cup W_n)^c$ agree throughout the path.
\end{theorem}
\begin{proof}
\leanok
Extend the reference from the buffered chart by a smooth cutoff. Convexity of
positive definite forms gives the path. The cutoff is constant on the inner
and outer open regions. The physical jet estimate uses the original measure
and the same width. The unit seam can cross the transition. A smaller closed disk
has an open flat neighborhood at positive $n$.
\end{proof}

\begin{theorem}[Signed price of a positive flat replacement]
\label{thm:ed03-mandatory-flat-cost}
\uses{thm:ed03-realized-disk-reference}
\lean{SmoothPoincare.balancedFlatReference_cost, SmoothPoincare.balancedFlatReference_trim_cost, SmoothPoincare.eventually_balancedFlatReference_cost, SmoothPoincare.eventually_balancedFlatReference_trim_charge_pos, SmoothPoincare.selectedReference_charge_eq_of_germ}
\leanok
Put $q=2^{-k}$, $I=\inf E_+$ over positive Yamabe classes on the fixed
oriented smooth type, $e=E_+(g)$ and $R=(B\cup W)^c$.
If the marked exterior division at depth $k$ has error $\epsilon$ and $h$
is an attained Yamabe-positive metric flat on $B\setminus W$, then for
each $s\in\{+,-\}$,
\[
 \lambda_h^s(W)+\lambda_h^s(R)-\lambda_g^s(R)
 > I-q\bigl(e-\lambda_g^s(Q)\bigr)-\epsilon .
\]
Matching exterior germs and frames put the full lower bound on $W$.
When $e_i\to I$ and $\epsilon_i\to0$, for every $\eta>0$ the left side
exceeds $(1-q)I-\eta$ for both signs at all sufficiently late indices,
for every such positive endpoint. If $I>0$ and $k>0$, an endpoint preserving
the full exterior charge must eventually have strictly positive trim charge
in each sign.
\end{theorem}
\begin{proof}
\leanok
The flat portion contributes zero, so the two new charges sum to
$E_s(h)\ge I$. Signature zero gives this inequality in both chiralities.
The full incoming exterior controls $\lambda_g^s(R)$ by the marked division.
Subtract and pass to the stated eventual inequality, keeping the finite core
charge nonnegative. This argument requires no minimizing metric at $I$.
\end{proof}

\begin{proposition}[Exact saving on a matched flat trim]
\label{prop:ed03-native-matched-flat}
\uses{thm:ed03-realized-disk-reference,thm:ed03-mandatory-flat-cost}
\lean{SmoothPoincare.exists_nativeConstantMetric_at, SmoothPoincare.exists_selectedMatchedFlatDisk, SmoothPoincare.selectedMatchedFlatDisk_trim_charge, SmoothPoincare.selectedMatchedFlatDisk_exact_saving, SmoothPoincare.selectedSplice_exterior_chiralMeasure_restrict_eq}
\leanok
Suppose the incoming metric in the given marking agrees with a constant
positive Gram matrix $L^*L$ on an open neighborhood of the complete closed
trim. The path to that flat reference has zero trim curvature at every time
and matches the original unit seam. Its endpoint is flat on the whole disk.
For an attained positive endpoint $h$ with the stated exterior frame agreement,
\[
 E_s(h)=\lambda_g^s((\operatorname{int}B)^c),\qquad
 E_+(g)-E_+(h)=\lambda_g^s(B)\le e-I\qquad(s=+,-).
\]
Equality of restricted measures for independent exterior frames is also
proved when their orthogonal action on two-forms is supplied with Hodge
projection and curvature-conjugacy identities.
\end{proposition}
\begin{proof}
\leanok
Pull back the Euclidean metric by a constant linear isomorphism to define the
reference tensor in the native parameters. The interpolation is constant on the
transition. Flatness of the endpoint and equality on the complete exterior
give the two charge identities. Membership of the positive endpoint in the
fixed-type energy set gives the upper bound on its saving.
\end{proof}

\begin{proposition}[Euler flux on the supplying flow]
\label{prop:ed03-spatial-euler-flux}
\uses{thm:ed03-modified-euler-transport}
\lean{SmoothPoincare.metricFrameEulerPfaffian_eq_chargeDensity, SmoothPoincare.metricFrameEulerCovariantDerivative_bianchi, SmoothPoincare.frameEulerPfaffian_hasDerivAt_of_bianchi, SmoothPoincare.ricciFlow_frameEulerTransgression_norm, SmoothPoincare.eulerCurvatureCoefficient_hasDerivAt_skewJet, SmoothPoincare.ricciFlowEulerPfaffian_fixedBase_hasDerivAt, SmoothPoincare.eulerSpatialFlux_integral_bound, SmoothPoincare.modifiedEuler_transport_of_spatialFlux}
\leanok
In an oriented orthonormal frame the normalized scalar Euler density is
twice the Pfaffian coefficient in the frame convention used here. The
transgression coefficients obey $|T|\le144|\mathrm{Rm}|\,|\nabla\mathrm{Ric}|$
in the proved frame norm. On an actual smooth Ricci solution the coefficient
variation is computed with the four base vectors fixed and endomorphism-frame
velocities $\mathrm{Ric}^{\sharp}e_a$ at the tested time. Suppose a smooth slice current $X$ realizes
$m'(s)=\int\chi\operatorname{div}_{g(s)}X_s\,dV_s$ in the open time interval.
With the corresponding current bound, slice curvature norm at most $A$,
integrated squared derivative norm at most $L^2$, and fixed smooth cutoff
$|d\chi|_s\le\delta(\tau/s)^b$, one has
\[
 |m(\tau)-m(0)|\le144\,\delta A L\sqrt{\tau/(1-2b)}
 \qquad (b<1/2,\ \tau>0)
\]
under the stated endpoint continuity and flux measurability.
\end{proposition}
\begin{proof}
\leanok
Apply the differential Bianchi identity to the frame product rule and the
Ricci equation to the orthogonal connection variation. Spatial integration
by parts and Cauchy--Schwarz on each supplying slice, followed by weighted
Cauchy--Schwarz in time, give the estimate. The current, the two endpoint
moments and every slice volume belong to the same evolution.
\end{proof}

\section{Marked birth bodies, curved disk cost and the tested Euler current}
\label{sec:ed03-marked-glue}

\begin{theorem}[Two actual incoming birth bodies]
\label{thm:ed03-birth-survivors}
\uses{thm:ed03-birth-edge-incidence}
\lean{SmoothPoincare.IsolatedZeroHandle.exists_survivingBirthTransport, SmoothPoincare.IsolatedZeroHandle.exists_twoBirthCharts_after_emptyCut, SmoothPoincare.IsolatedOneHandle.nativeFeet_on_oppositeIncomingSpheres, SmoothPoincare.IsolatedOneHandle.exists_twoBirthCharts_withNativeEdgeFeet}
\leanok
Suppose an empty incoming sublevel meets two successive isolated births with
regular intervening bands, followed by a selected index-one cut. Two ambient
smooth ball charts cover its entire incoming sublevel by disjoint closed
unit balls. Their unit spheres give the complete incoming level. If the two
body labels of the edge differ, each whole native attaching disk lands on a
different one of these spheres, up to interchange.
\end{theorem}
\begin{proof}
\leanok
Raise the height by a smooth bump on the second birth chart. The modified
regular band carries the old sublevel and its whole boundary past the new
ball. Transport both charts to the incoming level and separate their compact
images by open neighborhoods. Each native disk is connected. Its component
label determines which sphere contains all of its parameters.
\end{proof}

\begin{theorem}[Gluing at a selected trace]
\label{thm:ed03-marked-glue}
\uses{thm:ed03-birth-survivors,thm:ed03-rounded-native-rims}
\lean{SmoothPoincare.markedEdgeInverseOverlap, SmoothPoincare.markedEdgeSignedCollar_height, SmoothPoincare.IsolatedOneHandle.nativeDeepTraceDiffeomorph, SmoothPoincare.IsolatedOneHandle.markedEdgeTraceMap_transportTube, SmoothPoincare.IsolatedOneHandle.markedEdgeTraceMap_transportOutgoingTube, SmoothPoincare.IsolatedOneHandle.exists_markedEdgeTraceDiffeomorph_of_cover, SmoothPoincare.IsolatedOneHandle.exists_outgoingTraceAdjustment}
\leanok
On the strict deep lower region the stored trace is the actual lower
inclusion, with a smooth inverse. On a cover of the upper body and its
model by the two incoming patches, the rounded patch and the outgoing
patch, agreement of both maps on their open overlaps glues them to a
diffeomorphism $F$. Its map from the stored trace is $F\circ e$, for the
actual trace identification $e$. It acts on the lower and handle
representatives and every furnished framed tube by the corresponding patch
maps. A collar comparison preserves the signed height difference in each
direction. A compactly supported level isotopy can change the outgoing
trace marking while fixing the deep lower region and protected tubes.
\end{theorem}
\begin{proof}
\leanok
Glue the forward local maps on the source and their inverses on the target.
The two composites restrict to identities. On a common product strip the
transition is $B\circ(d\times\mathrm{id})\circ A^{-1}$. The trace
equations follow by composition with $e$ and $e^{-1}$. The two deep
identities in the isolated-handle data prove smooth restriction there.
A prescribed rounded model elsewhere requires its open comparison with
the actual trace.
\end{proof}

\begin{theorem}[A curved replacement on the selected width]
\label{thm:ed03-matched-curved}
\uses{thm:ed03-realized-disk-reference,prop:ed03-native-matched-flat}
\lean{SmoothPoincare.exists_selectedMatchedCurvedDisk, SmoothPoincare.exists_selectedClampedCorePath, SmoothPoincare.selectedMatchedCurvedDisk_trim_charge, SmoothPoincare.selectedMatchedCurvedDisk_energy, SmoothPoincare.selectedMatchedCurvedDisk_cost_window, SmoothPoincare.selectedMatchedCurvedDisk_core_charge_of_division}
\leanok
Let $Q\subset B$ be the selected marked disks, $W$ the actual native trim,
$e$ the incoming chiral energy, $q=2^{-k}$ and $I$ the fixed smooth-type
positive infimum. A curved native reference matching the source near the
closed trim extends to an affine path of positive-definite metrics, with
the same full trim and closed exterior germs at all times. The endpoint
realizes the reference on an open neighborhood of $B$. If this endpoint
$h$ lies in a positive class and exterior frames match, then for either
sign $s$ the attained saving is the decrease of disk charge. The retained
charge inside the disk away from its trim obeys
\[
 \lambda_h^s(B\setminus W)>
 I-q\bigl(e-\lambda_g^s(Q)\bigr)-\epsilon-\lambda_g^s(B\cap W).
\]
\end{theorem}
\begin{proof}
\leanok
Use the marked metric lift and the cutoff of width $1/(n+1)$.
On the full transition the two metrics agree, so interpolation fixes its
germ. Decompose the actual signed energy over the disk and its complete
exterior and apply the four-measure division and the positive-class floor.
A matched curved trim can carry positive charge throughout the path.
\end{proof}

\begin{proposition}[Price of a separate positive cap]
\label{prop:ed03-curved-cap-price}
\uses{thm:ed03-matched-curved,thm:ed03-mandatory-flat-cost}
\lean{SmoothPoincare.curvedCap_exteriorPrice, SmoothPoincare.curvedCap_outerTrimPrice, SmoothPoincare.exists_depth_eventually_curvedCap_price_gt}
\leanok
Write $V=(\operatorname{int} B)^c$ and
$\Delta_s(A)=\lambda_h^s(A)-\lambda_g^s(A)$. A positive cap with
$E_s(\mathrm{cap})\le\lambda_h^s(V)+\delta$ requires
\[
 I-q\bigl(e-\lambda_g^s(Q)\bigr)-\epsilon
 < \Delta_s(V\cap W)+\Delta_s(V\setminus W)+\delta .
\]
Along a minimizing sequence, for each $a<I$ one can fix a positive depth
before choosing the divisions. Eventually every such cap in either sign
has $\Delta_s(V)+\delta>a$ as the division errors tend to zero.
\end{proposition}
\begin{proof}
\leanok
Compare the attained cap energy with $I$ and use the upper estimate
on the complete original exterior. Split this exterior into two disjoint
measurable pieces at the actual trim. Choose $k$ with
$(1-2^{-k})I>a$ before passing along the sequence.
\end{proof}

\begin{theorem}[The Euler current on the modified flow]
\label{thm:ed03-tested-euler}
\uses{prop:ed03-spatial-euler-flux}
\lean{SmoothPoincare.globalEulerCurrent, SmoothPoincare.globalEulerCurrent_norm_le, SmoothPoincare.globalEulerCurrent_divergence_eq_primitiveJet, SmoothPoincare.exists_localEulerGaugeFrame, SmoothPoincare.riemannianVolumeDensity_eq_orientedBasisDet, SmoothPoincare.intrinsicEulerDensity_eq_modifiedChargeDensity, SmoothPoincare.ricciFlow_eulerVolume_hasDerivAt, SmoothPoincare.ricciFlow_testedEuler_hasDerivAt, SmoothPoincare.ricciFlow_testedEuler_hasDerivAt_gradient}
\leanok
Let $g(t)$ be the supplied smooth Ricci flow on a closed oriented
four-manifold, in the gauge of the modified solution. For every fixed
continuous test $\chi$ and every regular time,
\[
 \frac{d}{dt}\int\chi e(g)\,dV_g=\int\chi\,\operatorname{div}_g X\,dV_g,
 \qquad e(g)=\tfrac14|\mathrm{Rm}|^2-|\mathrm{Ric}|^2+\tfrac14R^2.
\]
In an oriented orthonormal frame $X_i=2\sum_{jab}
({*}\mathrm{Rm}{*})_{ijab}\nabla_a\mathrm{Ric}_{jb}$. The smooth global
field satisfies $|X|\le144|\mathrm{Rm}|\,|\nabla\mathrm{Ric}|$ in the
ordered tensor norms. For a smooth test the derivative equals
$-\int d\chi(X)\,dV_g$.
\end{theorem}
\begin{proof}
\leanok
Raise the one-form obtained by contracting the double dual of curvature
against the first covariant derivative of Ricci. Trace its covariant
derivative and use differential Bianchi to identify it with the variation
of twice the fixed-base Pfaffian. Transport the two bundle slots by the
Ricci gauge through a prescribed oriented frame. The determinant changing
the four base vectors is the evolving volume ratio. Relative to one fixed
volume the weighted density is smooth in spacetime. Its time derivative
may therefore be integrated against a fixed continuous cutoff. Spatial
integration by parts gives the smooth-test formula.
\end{proof}

\begin{theorem}[The retained mark on the whole original trace]
\label{thm:ed03-original-rounded-trace}
\uses{thm:ed03-birth-survivors,thm:ed03-marked-glue}
\lean{SmoothPoincare.IsolatedOneHandle.exists_bufferedNativeFootBirthOverlap, SmoothPoincare.IsolatedOneHandle.exists_bufferedNativeBlockBirthOverlap, SmoothPoincare.IsolatedOneHandle.exists_retainedRoundedTrace, SmoothPoincare.IsolatedOneHandle.RetainedRoundedTrace.inverse, SmoothPoincare.IsolatedOneHandle.RetainedRoundedTrace.exists_originalTraceComparison, SmoothPoincare.IsolatedOneHandle.RetainedRoundedTrace.outgoingLevel_inclusions, SmoothPoincare.IsolatedOneHandle.RetainedRoundedTrace.exists_outgoingLevelCollars_toNextCut, SmoothPoincare.IsolatedOneHandle.RetainedRoundedTrace.exists_twoSidedCollar_extending_toUpper}
\leanok
For each full incoming foot of the selected merger there is an open
four-dimensional comparison on $B(b,r)\times B(0,\sigma)$ with $r>0$ and
$1<\sigma<2$. The actual block coordinates and their pulled-back ball
coordinates agree under the birth chart $P$, and $P^{-1}$ recovers them
on the whole image. In the same Morse chart refine the two regular cuts.
The new trace factors as $U\circ C$, with explicit lower and handle-cell
representatives in the rounded sublevel. If $H$ is the upper regular
comparison from the new cut to the old and $e$ is the original stored
trace, the whole original trace maps to this rounded sublevel by
\[
 F=U^{-1}H^{-1}e,\qquad F^{-1}=e^{-1}HU .
\]
On each outgoing level the restrictions respect the marked-band square
in both directions, including every parameter of a supplied tube. Restricting $U$ to its entire boundary gives a smooth level map. On a common two-sided height strip its product transport goes from the rounded level through the actual outgoing level to the next cut. It agrees with $U$ at height zero. There is also an open two-sided ambient collar $D$ that agrees with $U$ on the entire inward half of the supplied rounded collar $A$. For every $0\le s<r$, its forward equation is $D(x,s)=U(A(x,-s))$ under sublevel inclusion, and its inverse recovers $A(x,-s)$.
\end{theorem}
\begin{proof}
\leanok
The compact native feet lie in the open target of their assigned birth
chart. Shrink the product neighborhood of each complete disk inside the
open domain of the block followed by $P^{-1}$. Retain both smooth factors
in the isolated-handle construction, including their literal point maps
and deep-lower identities. Restrict the upper regular band to sublevels
and levels, compose these diffeomorphisms and apply the inclusion square.
The inward restriction has bijective derivative through both sublevel inclusions at the boundary. Extend this smooth half collar across its zero section. The extension and the height-calibrated product collar have the same boundary marking. Comparing them on a whole strip requires further control.
This constructs the comparison for the particular rounded sublevel of
the height. A comparison carrying independently prescribed standard
ball charts still requires agreement on their model overlaps.
\end{proof}

\begin{theorem}[Positive curved references and signed charge]
\label{thm:ed03-positive-conformal-core}
\uses{thm:ed03-matched-curved,prop:ed03-curved-cap-price}
\lean{SmoothPoincare.exists_selectedNativeRoundCorePath, SmoothPoincare.exists_selectedPositiveConformalCoreReference, SmoothPoincare.conformalCore_scalar_quadratic, SmoothPoincare.signedCollarWeylBlock_conformal, SmoothPoincare.SignedCollarMetric.chiralMeasure_conformal, SmoothPoincare.SignedCollarMetric.orderedSpectrum_conformal, SmoothPoincare.balancedPhysicalSpectralMeasure_conformal}
\leanok
On a connected closed positive-class source, every buffered native disk
and every $\eta>0$ admit a smooth $u$ supported strictly inside the native
inner ball and an attained $0<|t|\le\eta$ such that $g_t=e^{2tu}g$ is
distinct from $g$ and has nonzero scalar curvature there. All $g_s$ lie
in the original positive class. Both chiral measures and the spectral-square
measure agree on every region in their respective metrics. The entire
trim and closed exterior remain fixed as metric germs. A second construction
inserts an actual unit-round neighborhood on the same disk by an affine
splice. Its Yamabe sign needs a separate estimate.
\end{theorem}
\begin{proof}
\leanok
Take a nonconstant smooth bump supported strictly inside the protected
width. At a point with $du\ne0$, the scalar-curvature numerator
$R-6t\Delta u-6t^2|du|^2$ cannot vanish at all three parameters
$\eta/2,-\eta/2,\eta$. Translation of the logarithmic conformal factor
gives equality of the class infima. In rescaled oriented frames the
Kulkarni--Nomizu correction has pure trace in each chiral block. After
the trace is removed both operators scale by $e^{-2tu}$ and their
squared coefficients cancel the factor $e^{4tu}$ in the volume.
\end{proof}

\begin{theorem}[Terminal ball charge in the modified metric]
\label{thm:ed03-terminal-ball-euler}
\uses{thm:ed03-tested-euler,thm:ed03-modified-euler-transport}
\lean{SmoothPoincare.lipschitzEulerSpatialFlux_integral_bound, SmoothPoincare.ricciFlow_terminalDistance_power_bound, SmoothPoincare.ricciFlow_terminalLipschitzEuler_transport, SmoothPoincare.terminalBallCutoff_properties, SmoothPoincare.ricciFlow_terminalBall_RicciCharge, SmoothPoincare.intrinsicTerminalCharge_eq_modifiedChargeDensity, SmoothPoincare.ricciFlow_terminalBall_intrinsicRicciCharge}
\leanok
Assume $[0,\tau]$ belongs to the carrier of the compact modified solution,
with regular interior and $\mathrm{Ric}_s\ge-(b/s)g_s$, $b<1/2$.
Suppose $\|\mathrm{Rm}\|_{L^2(g_s)}\le A$ and the time integral of the
squared $L^2(g_s)$ norm of $\nabla\mathrm{Ric}$ is finite and at most $L^2$.
For one fixed terminal $\delta$-Lipschitz cutoff the absolute change
of its Euler integral is at most $144\delta A L\sqrt{\tau/(1-2b)}$.
On a connected source take the terminal distance cutoff with plateau
$\{d_\tau\le R\}$ and vanishing on $\{d_\tau\ge R+1\}$. If a measurable
initial set in that plateau has full curvature charge at least
$\epsilon^2$, the terminal open ball satisfies
\[
 \int_{B_\tau(p,R+1)} |Z_\tau|^2\,dV_\tau\ \ge \tfrac{1}{2}\epsilon^2-4(m_++m_-)-\tfrac{1}{6}m_{R^2}
 -288 A L\sqrt{\tau/(1-2b)}.
\]
Here $Z_\tau$ denotes the trace-free part at the terminal time.
The $m$'s are the initial cutoff moments. The frame-independent version
combines their contribution as one integral of $\chi(|\mathrm{Rm}|^2/2+2e)$
in the initial volume. Endpoint integrability then follows from continuity
on the compact slices.
\end{theorem}
\begin{proof}
\leanok
The Ricci inequality makes $s^{-2b}g_s(v,v)$ nonincreasing for fixed
$x,v$. Thus $d_\tau\le(\tau/s)^b d_s$. Smooth approximation on each
slice gives the spatial flux bound with the exact Lipschitz constant
of the original cutoff. Integrate the derivative of its Euler moment,
use the carrier continuity at the endpoints and apply time
Cauchy--Schwarz. At the initial set the full tensor lower bound and
the two Weyl charges control the sign of the Euler moment. At the
terminal ball the two nonnegative Weyl terms and scalar-square term
give the displayed lower bound.
\end{proof}

\begin{proposition}[Four unchanged exterior tests and the cost inside the core]
\label{prop:ed03-protected-conformal-cost}
\uses{thm:ed03-positive-conformal-core,prop:ed03-curved-cap-price}
\lean{SmoothPoincare.metricFullCurvatureMeasure_conformal_restrict_closedExteriors, SmoothPoincare.BalancedPhysicalExteriorDivision.conformal_of_support, SmoothPoincare.exists_protectedPositiveConformalDivision, SmoothPoincare.conformalCore_scalarCost_eq, SmoothPoincare.conformalCore_fullCost_two_direction, SmoothPoincare.exists_selectedConformalCoreCosts_unbounded}
\leanok
One can make the positive curved change inside the protected disk $Q$
itself. At every amplitude the original four exterior divisions hold on
the same marked pair. The restriction of the full-curvature measure to
either complete closed exterior and to the supplied disjoint outer trim
is unchanged. On any measurable $A$, let $C(t,A)$ be the integral of
$|\mathrm{Rm}_{g_t}|^2$ in its own volume. Then
\[
 \tfrac12\bigl(C(t,A)+C(-t,A)\bigr)
 \ge 3t^4\int_A |du|_g^4\,dV_g-\tfrac16\int_A R_g^2\,dV_g .
\]
The coefficient of $t^4$ is positive on the inner ball for the constructed
bump. Full costs there exceed arbitrary ceilings at one of the two signs
for sufficiently large amplitudes.
\end{proposition}
\begin{proof}
\leanok
Off the core support the conformal metric, its adapted frames and both
volume forms agree as germs. Restrict this identity to each closed
exterior and use the signed and spectral measure identities globally.
In dimension four the transformed volume cancels the exponential in
the square of the scalar-curvature formula. Average the two resulting
squares at $t$ and $-t$, and use $|\mathrm{Rm}|^2\ge R^2/6$.
A nonzero differential on an open subset gives positive integral
against the actual source volume.
\end{proof}

\section{The outgoing collar, the protected disk and terminal Ricci charge}
\label{sec:ed03-collar-charge}

\begin{theorem}[Comparing the retained and calibrated collars]
\label{thm:ed03-actual-collar-transition}
\uses{thm:ed03-original-rounded-trace}
\lean{SmoothPoincare.IsolatedOneHandle.RetainedRoundedTrace.exists_outgoingCollarComparison, SmoothPoincare.IsolatedOneHandle.RetainedRoundedTrace.exists_outgoingCollarComparison_toNextCut, SmoothPoincare.IsolatedOneHandle.outgoingCollarTransition_fixedOn_iff, SmoothPoincare.IsolatedOneHandle.outgoingCollarTransition_normalDerivative_pos}
\leanok
Let $A$ be the calibrated source chart, $F$ the two-sided extension of
the retained inward map, and $P$ the target product parametrized by the
actual boundary marking. There is a smooth coordinate change
$J=P^{-1}F$ fixing the entire zero section. Its domain and inverse domain
contain the same full strip of positive width. Both $J$ and $J^{-1}$
preserve the inward and outward sides. Their tangential derivatives at the zero section are the identity and their normal derivatives there are strictly positive. At calibrated parameter $q=(x,t)$
the source chart becomes $A(J^{-1}q)$. Its rounded height is
$c_{\rm round}+(J^{-1}q)_2$. The actual target height is $c+t$
and the corresponding next cut has height $h_{\rm next}+t$.
On the inward side the retained body map and inverse give these actual
points. A prescribed parameter set is fixed by the transition precisely
where the two parametrized collars agree.
\end{theorem}
\begin{proof}
\leanok
Compose the partial diffeomorphisms on their open overlap and use
compactness of the full boundary level to put a two-sided strip in both
coordinate domains. The inward map takes strict interior to strict
interior. The calibration and injectivity separate the opposite side.
Apply the two inverse equations to each parameter of a framed tube.
\end{proof}

\begin{theorem}[Covering an independently prescribed ball]
\label{thm:ed03-target-ball-cover}
\uses{thm:ed03-marked-glue,thm:ed03-birth-survivors}
\lean{SmoothPoincare.exists_markedPatchDiffeomorph_of_unitBallCover, SmoothPoincare.IsolatedOneHandle.connectedTrace_of_twoBirthCharts, SmoothPoincare.IsolatedOneHandle.RetainedRoundedTrace.exists_markedBall_of_nativeRepCover, SmoothPoincare.markedEdgeModelFlatten_images, SmoothPoincare.markedEdgeModelCollar_coordinates, SmoothPoincare.markedEdgeRoundBallCollar_coordinates}
\leanok
On a binary merger the two connected incoming body images are joined by
the handle at its two feet. Suppose four local diffeomorphisms to the
independently specified closed four-ball agree forward on source overlaps,
and cover the images of every lower and handle representative in the
rounded body. They cover the entire target, glue to a diffeomorphism and
give its inverse on every patch target. One explicit Euclidean model is
$q(u,z)=(u^2-1)^2+|z|^2$. For $a>1$ the change
\[
 (u,z)\longmapsto
 \left(\frac{u}{\sqrt{a+1}},
      \frac{z}{\sqrt{(a+1)(a-1+u^2)}}\right)
\]
carries $q\le a^2$ onto the unit ball and its full level onto the unit sphere.
Writing this map as $d_a$, a collar of the round target is
$d_a d_{\sqrt{a^2+t}}^{-1}(v)$, for $v$ on the unit sphere and
$|t|<w<a^2-1$. It fixes every $v$ at zero. It lies in the closed ball
precisely for $t\le0$, with a smooth inverse recovering both parameters.
\end{theorem}
\begin{proof}
\leanok
The native lower and handle representatives exhaust the adjunction.
The glued map is a local diffeomorphism from a compact connected body.
The resulting covering of the simply connected ball has one sheet.
Its inverse agrees with each patch inverse by injectivity. The explicit
model follows by factoring the height defect into a positive function
times the difference of the squared norm from one.
\end{proof}

\begin{proposition}[The closed split and slid controls]
\label{prop:ed03-closed-slide-control}
\uses{thm:ed03-marked-glue}
\lean{SmoothPoincare.ClosedMarkedControl.splitExterior_torus, SmoothPoincare.ClosedMarkedControl.splitTorusCollar_transport, SmoothPoincare.ClosedMarkedControl.exists_splitClosingHomeomorph, SmoothPoincare.ClosedMarkedControl.slid_dualTube_and_inverse, SmoothPoincare.ClosedMarkedControl.slidClosedAssembly_representatives}
\leanok
In the split control the spectator in $S^3\subset\mathbb C^2$ has tube
$(z,v)\mapsto (\sqrt{1-\epsilon^2|v|^2}\,z,\epsilon v)$,
$0<\epsilon<1$. The explicit complementary solid torus agrees with
this tube on both boundary parameters. Its smooth signed collar recovers $t=|y_2|^2-\epsilon^2$ and both circle parameters throughout $-\epsilon^2<t<1-\epsilon^2$. The two marked disk pieces glue to $S^1\times S^2$ by the lower and upper hemisphere maps, with an inverse on each piece. Given the slid trace comparison
$F$, its boundary restriction and its agreement on the outgoing
spectator chart, the upper closing map is $\gamma'=\gamma\circ\partial F$.
The dual framed spectator and its inverse agree with those of the split
control. The full closed adjunction comparison is a homeomorphism,
acting by $F$ below and the identity above, with the inverse formulas
on every representative.
\end{proposition}
\begin{proof}
\leanok
Evaluate the complementary torus chart at its two boundary circle
parameters. For the upper dual chart, swap the two factors and apply
the closing map. Its inverse applies the inverse closing map before
unswapping. The two representative equations descend through the
closed attachment because the closing square commutes.
\end{proof}

\begin{theorem}[A finite nonconformal disk variation]
\label{thm:ed03-signed-first-variation}
\uses{thm:ed03-positive-conformal-core,prop:ed03-curved-cap-price}
\lean{SmoothPoincare.exists_clampedAffineMetricFamily, SmoothPoincare.exists_selectedNonconformalDiskSaving_of_weightedSquare, SmoothPoincare.protectedDisk_positiveCap_allowance_lower}
\leanok
Fix one finite positive source below $2T$, its disk $B$, protected core
$Q$, and native trim $W$. A localized variation supported in the interior
of $B$ away from $Q\cup W$ preserves the prescribed metric germs.
Suppose its signed energy derivative at zero is $-D/2$,
where $D=\int\chi b^2\,dV_g>0$, its tangent at a point of the indicated
affine path is nonscalar, and a continuous sign detector is positive there
at parameter zero. With the stated frame agreement off the support, some
$t>0$ gives a positive nonconformal endpoint. Its saving in each
chirality exceeds $tD/4$ and is at most $e-I$. The whole short joined path
stays strictly below $2T$. If a same-type positive cap costs at most
$\lambda_h^s(V)+\delta_s$, $V=(\mathrm{int}\,B)^c$, retention of disjoint
$Q$ and $B\cap W$ on a subcritical joined endpoint forces
\[
 \delta_s>I+\lambda_g^s(Q)+\lambda_g^s(B\cap W)-2T .
\]
\end{theorem}
\begin{proof}
\leanok
Clamp the real parameter to $[-1/2,1/2]$ for a symmetric direction
with $|D_x(v,v)|\le g_x(v,v)$. Every metric is spatially smooth and
positive definite, and the path is affine near zero. The derivative
and continuity at zero select the finite interval. The signature identity
and the unchanged complete exterior transfer the saving to each signed
disk charge. For the cap bound, add the retained terms inside $B$ and
use $I\le E_s(\mathrm{cap})$ in the cap's own positive class.
\end{proof}

\begin{proposition}[Paying a cap from the net saving]
\label{prop:ed03-net-saving}
\uses{thm:ed03-signed-first-variation}
\lean{SmoothPoincare.diskCharge_add_joinedCost_le_excess_of_cap_paid_from_saving, SmoothPoincare.eventually_no_paid_cap_of_protectedDisk_charge, SmoothPoincare.eventually_no_paid_cap_of_original_charged_markedDisk, SmoothPoincare.original_markedDisk_charge_eq_and_le_excess_of_arbitrarily_small_cap_drift}
\leanok
On an actual finite source let $e$ be the old energy, $h$ the whole
replacement, and $j_s\ge0$ a measured additional joined cost. If the cap
above has allowance $\delta_s$ and $\delta_s+j_s\le e-E_s(h)$,
then $\lambda_h^s(B)+j_s\le e-I$. Along a sequence $e_i\to I$, a
uniformly charged retained core in even one sign, allowed to depend on $i$,
eventually rules out this payment in both signs simultaneously. This
applies when the original marked disk carries a fixed positive charge in
one sign on a selected subsequence and lies inside the protected core.
At one fixed source, suppose for every $\eta>0$ there is a single
positive same-type cap whose energy in each sign differs by less than
$\eta$ from the corresponding original exterior charge. Then the two
original disk charges agree, and each is at most $e-I$.
\end{proposition}
\begin{proof}
\leanok
Use the full identity $E_s(h)=\lambda_h^s(V)+\lambda_h^s(B)$ and the
cap bound. The replacement exterior cancels even when its charge has
changed. Monotonicity on the protected subset and convergence of $e_i-I$
give the sequence conclusion. At one fixed source apply the same-type cap
floor and signature identity at each tolerance before letting it decrease to zero.
\end{proof}

\begin{theorem}[Trace-free Ricci charge in a pointed terminal ball]
\label{thm:ed03-pointed-terminal-charge}
\uses{thm:ed03-terminal-ball-euler}
\lean{SmoothPoincare.tendsto_terminalBallCutoff_weylScalarCharge_of_initial_region, SmoothPoincare.eventually_ricciFlow_terminalBall_intrinsicRicciCharge, SmoothPoincare.intrinsicTraceFreeRicciSq_eq_of_partialDiffeomorph_inner, SmoothPoincare.terminalTraceFreeRicciSq_of_twoJets, SmoothPoincare.pointedTerminalBuffer_tendsto_traceFreeRicci, SmoothPoincare.pointedTerminalBuffer_traceFreeRicci_charge, SmoothPoincare.ricciFlow_terminalBall_intrinsicRicciCharge_pointedLimit}
\leanok
Consider a sequence of smooth modified Ricci solutions with the same fixed
terminal time $\tau>0$, curvature norm bound $A$, integrated derivative
budget $L^2$, and lower Ricci bound $-(b/s)g_s$, $b<1/2$. Suppose each
initial attaining set has full curvature mass at least $\epsilon^2$ and
lies in the radius-$R$ terminal plateau. At time zero the combined
Weyl and scalar-square charge tends to zero on measurable regions
containing the entire radius-$(R+1)$ terminal ball. Fix
\[
 0<c<\frac{\epsilon^2}{2}-288 A L\sqrt{\frac{\tau}{1-2b}} .
\]
Then the terminal ball eventually carries trace-free Ricci mass at least
$c$. Suppose these same terminal metrics and centers have complete pointed
smooth convergence on a cofinal subsequence in the canonical pullback
domains. For every $r>R+1>0$ the limit satisfies
\[
  \int_{B(p,r)} |\mathrm{Ric}^{\circ}|^2\,dV\ \ge c .
\]
\end{theorem}
\begin{proof}
\leanok
Bound the moment of the fixed cutoff by the nonnegative initial charge on
its entire support. Apply the terminal Euler estimate separately on each
solution. On a compact buffer the pulled-back metric two-jets converge.
Coercivity and the jet bounds give pointwise convergence of
$|\mathrm{Ric}|^2-R_{\rm scal}^2/4$ and a uniform spatial bound.
Use the Riemannian volume change of variables and dominated convergence
on that compact set. Choose $R+1<a<r$. The lower metric bound on
$\overline B(p,a)$ prevents a path of source length less than $R+1$ from
exiting its image for all sufficiently large indices. Thus the comparison
of integrals includes the entire source charge ball.
\end{proof}

\begin{proposition}[Volume escape and a positive scalar point]
\label{prop:ed03-volume-and-scalar}
\uses{thm:ed03-pointed-terminal-charge}
\lean{SmoothPoincare.ricciFlow_terminal_sqrtVolume_change_le, SmoothPoincare.ricciFlow_terminal_volume_tendsto_atTop, SmoothPoincare.pointedTerminal_noncompact_of_volume_escape, SmoothPoincare.ricciFlow_terminalPointed_noncompact_of_initialVolume, SmoothPoincare.ricciFlow_exists_scalar_pos_in_charged_ball}
\leanok
For the same connected compact source slices, write $V_i(s)$ for their
total volumes. The spatial curvature budget gives
\[
 |\sqrt{V_i(\tau)}-\sqrt{V_i(0)}|
 \le (\sqrt 6/2)A\tau .
\]
If their initial volumes diverge, their smooth pointed terminal limit
is noncompact. Further suppose the charged limit slice belongs to a
regular limiting Ricci solution. If scalar curvature is nonnegative
on the fixed charged open ball for all times in a neighborhood of the
slice, at a point of that ball both scalar curvature and
$|\mathrm{Ric}^{\circ}|^2$ are strictly positive.
\end{proposition}
\begin{proof}
\leanok
Differentiate $\sqrt{V}$ in the regular interior, use
$|R_{\rm scal}|^2\le6|\mathrm{Rm}|^2$ and Cauchy--Schwarz, then use
continuity at both endpoints. A compact pointed limit would globalize
the comparison and have finite total volume equal to the limit of the
source volumes. Positive integrated charge gives a point of positive
trace-free Ricci square. If its scalar curvature vanished, the time
derivative there would be zero and the spatial Laplacian nonnegative.
The scalar evolution equation has a strictly positive Ricci term there.
\end{proof}

The target-ball theorem still requires native forward charts with the
prescribed overlap and tube equations. The finite disk variation still
requires its energy derivative on the selected physical metric. The
charged terminal limit uses the supplied pointed comparison at positive
time. A comparison at the separated quantiles in the original physical
metric has its own distance, curvature and marking hypotheses.

\section{Edition 04: geometric inputs for smooth recognition}
\label{sec:ed04}
The following are proposed formalization targets and a coverage question.
The smooth category, the markings, and the quantifiers are part of each
target. Write $T=4\pi^2$. For an oriented metric, $E_\sigma(g;U)$
is the integral of the squared Hilbert--Schmidt norm of the signed Weyl
block over $U$. Write $I^+(X)$ for the infimum of $E_+$ over the
positive Yamabe classes on the indicated smooth type.

\begin{openproblem}[Round product at a fixed physical face]
\label{ed04:exact-face}
On a closed oriented smooth homotopy four-sphere in a positive class
of energy $e<2T$, fix one smoothly marked separating three-sphere with
conformally round induced metric and umbilic second fundamental form.
Formalize its conformal Gaussian normalization and cutoff on two
independently chosen positive subwidths $w_L,w_R$. For each tolerance,
construct a smooth path fixed on the outer germs with product endpoint,
peak joined energy below $2T$, and signed drift on each complete side
smaller in absolute value than the tolerance. All choices are made on
this fixed source.
\end{openproblem}

\begin{openproblem}[Positivity on the two actual caps]
\label{ed04:caps}
\uses{ed04:exact-face,lem:master-wall,prop:ed03-inclusive-caps}
From the joined path, insert a round product cylinder of arbitrary
finite length. Transplant each compactly supported test to the joined
metric, retain a uniform positive Yamabe bound, and pass through the
added center on each cap by density. Identify each capped smooth type
in its prescribed angular marking and recover the original oriented
manifold by its actual boundary map. Each cap energy in either sign is
the corresponding endpoint side energy.
\end{openproblem}

\begin{openproblem}[Limiting raw masses on fixed capped types]
\label{ed04:side-masses}
\uses{ed04:caps,prop:ed03-inclusive-caps}
Under the exact-face hypotheses, prove that on each original side the
two chiral masses agree at a value $m_j$ and $I^+(\widehat P_j)\le m_j$.
For a nonstandard source precisely one capped side has its smooth type
and the other has raw mass at most $e-I^+(M)$. First construct actual
positive caps for each tolerance on this fixed source, then take the
infimum on each fixed capped type. For a parametrized disk the exterior
cap carries the original smooth type. Its excess bound is the disk
specialization recorded in Proposition~\ref{prop:ed03-net-saving}.
\end{openproblem}

\begin{openproblem}[Recognition of the intermediate three-manifold]
\label{ed04:prefix}
For a closed smooth homotopy sphere in a reduced mixed presentation
with one lower, two middle and one upper handle of indices one, two and
three respectively, realize the winding row $(1,0)$ by marked slides.
For the winding-one first boundary, a map of nonzero degree from a closed
oriented three-manifold of positive scalar curvature implies that the
boundary is diffeomorphic to $S^3$. For the winding-zero first boundary,
the corresponding conclusion is the connected sum of two copies of
$S^1\times S^2$. In the latter case the smooth filling and its
compressions are required to exclude residual homology-sphere summands.
In either case, use the original framed trace and top attachment to
identify the closed four-manifold.
\end{openproblem}

\begin{openproblem}[Reduction followed by a transported second test]
\label{ed04:successor}
For a reduced presentation with handle counts $(1,1,3,2,1)$ in index
order and normalized winding row $(1,0,0)$, a positive-scalar-curvature
domination certificate on the first winding-one boundary gives a
presentation on the same closed smooth type with no lower one-handle
and two middle handles. Carry both surviving attaching maps and the
upper identification through the reduction. Domination on the second
actual prefix then recognizes the closed type. For counts $(1,2,3,1,1)$
use the reversed presentation with its own cuts and markings.
\end{openproblem}

\begin{openproblem}[Positive entry on a marked double]
\label{ed04:entry}
Fix a smooth compact oriented contractible half $W$, a boundary metric
$\gamma$ in its marking and an orientation-preserving $\gamma$-isometry
$f$. Seek a positive-scalar-curvature extension of the exact tensor
$\gamma$ with strictly positive Robin trace form on every nonzero
$H^{1/2}$ trace. The sign criterion supplies a totally geodesic positive
filling in the marked conformal class. For an actual such filling $h$,
write $B_h=E_+(h;W)+E_-(h;W)$. Construct positive metrics on the fixed
double $W\cup_f(-W)$ with both chiral energies tending to $B_h$.
Writing $\Lambda$ for the greater of the two top Weyl eigenvalues and
$G_h=\int_W\Lambda_h^2\,dV_h$, the closed spectral quantity
$4E_++3\int\Lambda^2\,dV$ should tend to $4B_h+6G_h$.
If $B_h<2T$, one finite member enters the subcritical band. A witness
on a candidate $M$ requires its marked smooth identification with this
double and the budget of the resulting filling.
\end{openproblem}

\begin{openproblem}[Selection on every untreated smooth type]
\label{ed04:coverage}
For each hypothetical nonstandard smooth homotopy four-sphere, take a
finite handle presentation minimizing the number of middle handles.
Produce a marked strict net reduction of this minimum, direct smooth
recognition, or positive entry together with a finite positive
competitor of energy strictly below $I^+(M)$ on this very type.
The alternative used may depend on $M$. A single entry metric, an
excluded shape of a charged disk, and a boundary test on a known
presentation each address only part of this selection problem.
\end{openproblem}

\section{Edition 05: entry, fixed-type attainment and rigidity}
\label{sec:ed05}
The open problems describe the geometric targets. The checked propositions below give their current formal prerequisites. Fix a closed connected oriented
smooth homotopy four-sphere $M$. Put $T=4\pi^2$,
$e=E_+=E_-$ on closed metrics, and $I=I^+(M)$.
The equality of the two closed integrals does not identify their restrictions
to a subdomain. All energy and volume measurements use their supplying metric.

\begin{openproblem}[Positive entry on a specified smooth type]
\label{prop:ed05-entry}
On a specified $M$ construct a positive Yamabe class with $e<2T$.
One target is a smooth representative with $R>0$ and
$|\operatorname{Ric}_0|^2<R^2/12$ everywhere. For a marked handle
boundary establish the smooth identification and the estimates on the entire
closed representative. Record which smooth types the construction covers.
\end{openproblem}

\begin{openproblem}[Budget on the original regularized trajectory]
\label{prop:ed05-bour-energy}
\uses{lem:ed03-positive-sequence,prop:ed47-budget,prop:ed47-motion}
Given entry, choose $0<\alpha_j<1$ decreasing to zero and unit-volume
positive Yamabe metrics with $0\le e_{j,0}-I\le\alpha_j^2$. Write
$D=\int|\operatorname{Ric}_0|^2$ and $S=\int R^2$. Along the original
negative twice-gradient flow of
$\mathcal F^\alpha=2(1-\alpha)e+\alpha D/2$,
prove that for all sufficiently large $j$ and every smooth time,
$Y\ge y>0$, $I\le e(t)\le 2T-\eta$ and
$D(t)\le4I+4(e_{j,0}-I)/\alpha_j\le d_*<32\pi^2$,
with $y,\eta,d_*$ independent of $j$ and time.
\end{openproblem}

\begin{openproblem}[The two outputs of one fixed-parameter flow]
\label{prop:ed05-bour-alternative}
\uses{prop:ed05-bour-energy}
At a fixed sufficiently small positive $\alpha_j$, instantiate Bour's
strictly pinched alternative on that original maximal solution. Return either
a closed $\mathcal F^{\alpha_j}$-critical metric via whole smooth comparison
or a complete nonflat scalar-flat Bach-flat pointed limit with the
physical bounds and coherent embeddings on compact domains. Account for
the short-time gauge and derivative estimates at fixed parameter.
\end{openproblem}

\begin{openproblem}[Smooth identification of a pointed limit]
\label{prop:ed05-bour-pointed}
\uses{prop:ed47-mark,prop:ed47-ends}
Take a smooth pointed limit of rescaled original metrics with the uniform
positive Yamabe, subcritical chiral and strict Ricci bounds of the preceding
budget, finite full curvature, and $R=B=0$ on the complete nonflat limit.
Use one oriented embedding on a neighborhood of the connected truncated
core and every end buffer. Exclude nontrivial end groups before using
$D_X=2(E_+(X)+E_-(X))+32\pi^2(k-1)$ to deduce $k=1$.
Construct the smooth conformal compactification $N$ and its positive
class. Identify $N$ with $M$ by forward and inverse smooth maps
agreeing on the actual cap collar. Prove $I\le e_X\le\liminf_i e_i$.
\end{openproblem}

\begin{openproblem}[Uniform estimates for the mixed critical sequence]
\label{prop:ed05-bour-stationary}
\uses{prop:ed05-bour-alternative,prop:ed05-bour-pointed,prop:ed47-balls,prop:ed47-mixed}
Return one unit-volume constant-scalar metric $q_j$ on $M$ at each late
index. Its parameter is $p_j=\alpha_j$ for a closed regularized return
and $p_j=0$ for an exact pointed return in Yamabe gauge.
Recover the scalar, physical Ricci and energy bounds in the gauges used.
From the constant-scalar trace-free Ricci equation with principal term
$\tfrac12\nabla^*\nabla\operatorname{Ric}_0$, prove noncollapse,
harmonic-coordinate regularity and estimates of every order uniform
as $p_j\to0$, also on curvature-normalized balls in the unbounded branch.
\end{openproblem}

\begin{openproblem}[Closed return at bounded curvature]
\label{prop:ed05-bour-closed}
\uses{prop:ed05-bour-stationary,prop:ed47-return}
If the unscaled curvature of the selected critical sequence is bounded
on a subsequence, obtain smooth convergence with whole oriented
comparison diffeomorphisms. Pass the positive Yamabe floor, the
vanishing Bach limit and $e(q_j)\to I$ to the limit and pull back to $M$.
\end{openproblem}

\begin{openproblem}[An attained positive Bach-flat class on the entered type]
\label{thm:ed05-bour-attainment}
\uses{prop:ed05-bour-closed,prop:ed05-bour-pointed,prop:ed05-bour-stationary}
From one positive subcritical entry on $M$ construct a smooth positive
Bach-flat class on $M$ of energy exactly $I$. In the unbounded-curvature
branch of the mixed critical sequence, use its own nonflat normalized
pointed limit, the marked return, and $I\le e_X\le\liminf e(q_j)=I$.
\end{openproblem}

\begin{openproblem}[Uniform infimum gap on an entered type]
\label{prop:ed05-uniform-gap}
\uses{prop:master-recognition}
Realize the master's uniform recognition gap as a separate formal
dependency. There is an $\varepsilon_*>0$, independent of the
oriented smooth homotopy four-sphere $M$, such that a nonempty
positive-class energy set with $I^+(M)\le T+\varepsilon_*$ implies
that $M$ is diffeomorphic to the standard sphere. Trace the
positive-time smoothing, compactness and smooth type hypotheses.
\end{openproblem}

\begin{openproblem}[Rigidity of an attained minimizer]
\label{prop:ed05-rigidity}
\uses{prop:ed47-variation}
Assume a smooth positive class attains $I$ in the band above the
recognition gap and below $2T$. For a specified candidate prove smooth
recognition of $M$ or produce a strictly lower positive competitor on
$M$. In a Hessian test exhibit an actual admissible negative direction
modulo gauge. Record the surviving Einstein and non-Einstein cases.
\end{openproblem}

\begin{openproblem}[Coverage of the smooth types]
\label{thm:ed05-coverage}
\uses{prop:ed05-entry,thm:ed05-bour-attainment,prop:ed05-rigidity,prop:ed05-uniform-gap,prop:master-recognition,prop:master-immersion}
For every oriented smooth homotopy four-sphere give either an applicable
direct recognition certificate or an entry witness and a rigidity
argument covering the minimizing class returned on that same type.
For a handle argument preserve the actual markings and establish
exhaustiveness of its allowed presentations and successors.
\end{openproblem}

\subsection{Checked results for the original and returned metrics}

\begin{proposition}[Yamabe starts at prescribed excess]
\label{prop:ed47-start}
\lean{SmoothPoincare.BourYamabeNormalization, SmoothPoincare.exists_bourYamabeStarts}
\leanok
Given one measured positive entry below $8\pi^2$ and a conformal Yamabe
normalization preserving the measured energy, every positive sequence
$\alpha_j$ admits unit-volume positive Yamabe starts on this carrier.
Their excess lies in $[0,\alpha_j^2]$, their scalar-square integral is
$Y_{j,0}^2$, and they satisfy a common strict chiral margin.
\end{proposition}
\begin{proof}
\leanok
Use the entry to inhabit the positive energy set. Select a member below the
smaller of $I+\alpha_j^2$ and the midpoint between the entry energy and
$8\pi^2$. Normalize the selected metric in its own class.
\end{proof}

\begin{proposition}[Physical budget at every original time]
\label{prop:ed47-budget}
\uses{prop:ed47-start}
\lean{SmoothPoincare.Bour.bourChargeOn_full_eq, SmoothPoincare.Bour.bourEuler_eq_charges, SmoothPoincare.Bour.bourEnergy_eq_chiral_scalar, SmoothPoincare.Bour.bourEnergy_conformal_minimum, SmoothPoincare.Bour.bourRicci_eq_integral, SmoothPoincare.Bour.bourFullCurvature_eq_euler_add_ricci, SmoothPoincare.Bour.bourOriginalFlow_budget, SmoothPoincare.Bour.eventually_bourOriginalFlow_budget}
\leanok
Supply the closed Euler value $16\pi^2$, the signed metric measurements,
the unit-volume Yamabe representatives, the round initial upper bound,
a continuous sign detector and the original functional sublevel.
Put $\delta=8\pi^2-I>0$ and $\epsilon_j=e_{j,0}-I\le\alpha_j^2$.
When $0<\alpha_j<1$ tends to zero, one cutoff in $j$ gives, at every
original time,
\[
 Y_t\ge\sqrt{24\delta},\qquad I\le E_\pm(g_t)\le8\pi^2-\delta/2,
 \qquad S_t\le Y_{j,0}^2+48\alpha_j .
\]
The physical two-tensor norm satisfies
\[
 D_t\le4I+4\epsilon_j/\alpha_j+(Y_{j,0}^2-384\pi^2)/12
 \le32\pi^2-2\delta.
\]
The curvature-operator integral is $16\pi^2+D_t$.
Both total chiral energies converge uniformly over the original time domains to $I$.
\end{proposition}
\begin{proof}
\leanok
Integrate the local four-slot curvature decomposition before using closed
signature. The forms $2(1-\alpha)e+\alpha D/2$ and
$2e+\alpha S/24-16\alpha\pi^2$ agree. A strictly positive pinching
margin bounds the Yamabe square away from zero. Connectedness propagates
the initial sign. Every slice is now admissible in the same infimum, which
gives the displayed inequalities by functional decrease.
\end{proof}

\begin{proposition}[Dissipation and the squared parabolic time scale]
\label{prop:ed47-motion}
\lean{SmoothPoincare.Bour.original_gradient_sublevel, SmoothPoincare.Bour.metricVelocity_scale_time, SmoothPoincare.Bour.gradientSq_scale, SmoothPoincare.Bour.parabolic_dissipation_eq, SmoothPoincare.Bour.bourEnergy_scale, SmoothPoincare.Bour.volume_constant_of_scale_variation}
\leanok
For a supplied geometric first variation and actual velocity $-2G$ on
a regular slab, $\int_a^b\|G\|_2^2=(F_a-F_b)/2$ and $F_b\le F_a$.
Under $c g(t_0+s/c^2)$ the covariant velocity has weight $c^{-1}$ and
the spatial gradient integral has weight $c^{-2}$. The time-integrated
dissipation on corresponding slabs agrees. The closed energy is
constant under scaling at fixed Euler value. Volume is constant given
the regular volume variation, its velocity trace and scale-direction
first variation.
\end{proposition}
\begin{proof}
\leanok
Insert the metric derivative in the physical tensor pairing and integrate
the derivative of $F$. Four-volume scales by $c^2$. The time substitution
cancels the remaining power. The derivative of the functional in the
scaling direction vanishes, giving zero integrated velocity trace.
\end{proof}

\begin{proposition}[One comparison of the full truncation]
\label{prop:ed47-mark}
\lean{SmoothPoincare.bourWholeCoreComparison, SmoothPoincare.bourComparison_orientationAt, SmoothPoincare.bourPointedTraceFreeRicci_real_le_liminf, SmoothPoincare.bourPointedSignedWeyl_le_liminf, SmoothPoincare.bourGlueMarkedInverse}
\leanok
Given pointed exhaustion maps with the specified orientation, a compact
truncation containing the core and all buffers is eventually covered by
one open comparison and its inverse. Canonical two-jet metric convergence,
actual constant rescalings of compact sources and an eventual Ricci ceiling
give $D_X\le\liminf D_i$ in the original physical norm. Each Weyl sign
has its separate comparison under the signed frame and density hypotheses.
Two supplied partial diffeomorphisms on the core and its exterior, with
the prescribed agreement, image overlap, coverage and orientation at a
core point, glue to an oriented diffeomorphism with both inverse germs.
\end{proposition}
\begin{proof}
\leanok
Restrict the single late comparison. Exhaust the limit with compact sets
before bounding the full lower limit. The Ricci density has scale weight
$-2$ and the four-volume has weight $2$. Glue the two compact neighborhoods,
globalize using source and image coverage, and restrict the inverse.
\end{proof}

\begin{proposition}[End count and both capped exterior inequalities]
\label{prop:ed47-ends}
\lean{SmoothPoincare.bourOneEuclideanEnd, SmoothPoincare.bourSignedEndEnergy, SmoothPoincare.bourCappedExterior_both_lt}
\leanok
After triviality of all end groups and the spherical boundary Euler
identity have been supplied, $k\ge1$ and $D_X<32\pi^2$ imply $k=1$.
Equal signed totals on the compactified carrier then give $D_X=4e_X$.
For a specified cap the two signed exterior tests give strict capped
energies below $T$ under their cut, ball-mass and splice margins.
\end{proposition}
\begin{proof}
\leanok
The boundary identity is $D_X=2(E_++E_-)+32\pi^2(k-1)$. Nonnegative
signed charges exclude $k\ge2$. Subtract the signed interior masses and
add the measured splice on each complete exterior.
\end{proof}

\begin{proposition}[Uniform physical ball volumes for the CSC sequence]
\label{prop:ed47-balls}
\lean{SmoothPoincare.bourYamabe_smooth_test_scale, SmoothPoincare.bourBallVolume_lower_of_cutoff, SmoothPoincare.bourBallVolume_lower_of_yamabe_scale, SmoothPoincare.bourTraceFreeRicci_scale, SmoothPoincare.bourVolume_scale}
\leanok
On the connected closed four-manifold let $Y\ge y>0$, with constant
scalar $0\le R\le C$. For every $a\ge1$, center $x$ and $0<r\le1$,
\[
 \operatorname{Vol}_{a g} B_{a g}(x,r)\ge
 \left(\frac{y}{4(6L^2+C)}\right)^2 r^4,
\]
where $L$ is the fixed smooth cutoff slope. The constant works on the
original and curvature-normalized metrics at every index.
\end{proposition}
\begin{proof}
\leanok
Use the conformal-class inequality on nonnegative smooth tests. In the
scaled metric smooth a Lipschitz cutoff, bound its gradient by $L/r$,
then iterate the ball-volume inequality from the small-radius lower bound.
\end{proof}

\begin{proposition}[Initial Yamabe saturation in a second pointed return]
\label{prop:ed47-mixed}
\uses{prop:ed47-budget,prop:ed47-mark,prop:ed47-ends}
\lean{SmoothPoincare.bourFirstPointed_excess_and_deficit, SmoothPoincare.Bour.fixedReturn_budget, SmoothPoincare.Bour.mixedReturn_energy_and_parameter, SmoothPoincare.Bour.mixedReturn_pointed_saturation}
\leanok
Suppose the closed first returns give convergence of their physical
energies, and the pointed first returns give the marked Weyl and Ricci
comparisons and the scalar bound in their new Yamabe gauge. Positivity,
the closed Euler evaluations and original budgets supply $p_j\to0$ and
$E_\pm(q_j)\to I$. On a second marked pointed subsequence, let the
positive returned metric have $e_X\le\liminf E_+(q_j)$ and
$4e_X\le\liminf D(q_j)$. Then $e_X=I$ and, on that subsequence,
$Y_{j,0}^2\to384\pi^2$.
\end{proposition}
\begin{proof}
\leanok
Write $b_j=384\pi^2-Y_{j,0}^2$. A first pointed return gives
$b_j\le48\alpha_j$. A closed one gives
$b_j\le48\alpha_j+12(4I-D(q_j))$. The oriented positive return gives
$I\le e_X$. Use $4I\le\liminf D(q_j)$ to control the closed terms
even when the two first alternatives alternate.
\end{proof}

\begin{proposition}[Positive return on the original smooth carrier]
\label{prop:ed47-return}
\lean{SmoothPoincare.SignedCollarMetric.pullbackCross, SmoothPoincare.SignedCollarMetric.energy_tendsto_of_metricCPConvergence, SmoothPoincare.bourAttainsOfSignedPullback, SmoothPoincare.bourBachFlat_of_criticalLimit, SmoothPoincare.bourClosedReturn, SmoothPoincare.bourRoundReturn}
\leanok
A supplied oriented diffeomorphism pulls a positive measured metric of
energy at most the source infimum back to an attaining metric in both
chiralities. Bach-flat transport uses the specified covariant operator
and naturality input. On a supplied closed mixed return the energy,
conformal test and covariant Euler limits give the same conclusion. On
a fixed carrier, the two-jet energy comparison uses compatible frames,
weighted measurability and domination. The actual round pullback has both
energies zero.
\end{proposition}
\begin{proof}
\leanok
Transport the signed geometric witness and use membership in the positive
energy set for the opposite inequality. On the closed limit pass each
conformal test separately and take the limit of the supplied Euler identity.
\end{proof}

\begin{proposition}[Second variation at an attained geometric metric]
\label{prop:ed47-variation}
\uses{prop:ed47-return}
\lean{SmoothPoincare.bourAttainedVariation, SmoothPoincare.bourNegativeSecondVariationExcludesAttainment, SmoothPoincare.bourRoundAttainedVariation}
\leanok
Along a two-sided, jointly smooth positive metric variation with twice
continuously differentiable measured Weyl energy, an attaining metric
has zero first and nonnegative second derivative. Negative second
derivative in a supplied direction produces arbitrarily close positive
competitors of strictly lower energy. The signed round pullback satisfies
the necessary variation conditions.
\end{proposition}
\begin{proof}
\leanok
Positive-class membership gives a local minimum along the path.
Apply the second-derivative test to the real energy function.
\end{proof}

\section{Edition 06 geometric suppliers and attained rigidity}
\label{sec:ed06}
These open formal milestones follow the master entry (E), fixed-type
attainment (X) and rigidity (R) plan. Use the oriented smooth carrier and
physical measures of Section~\ref{sec:ed05}. Thus $T=4\pi^2$,
$e=E_+=E_-$ on the closed smooth homotopy four-sphere $M$, and
$I=I^+(M)$. Write $D=\int|\operatorname{Ric}_0|^2\,dV$.
The assessed attainment argument in Investigation 06
supplies the scientific starting point for R. The formal attainment
and recognition-gap targets remain open in
Theorem~\ref{thm:ed05-bour-attainment} and
Proposition~\ref{prop:ed05-uniform-gap}.

\subsection{Entry and independent marked reduction}

\begin{openproblem}[Positive entry on a marked handlebody boundary]
\label{prob:ed06-entry-construction}
\uses{prop:ed04-positive-class,prop:ed04-round-face-control}
Route 11. Fix one actual balanced-presentation five-dimensional
handlebody with handles of indices zero, one and two whose oriented
boundary is a smooth homotopy four-sphere. Specify the attaching circles,
framings, retained exterior and a compatible new cap. The tested product
circle cap already costs at least $64\pi^2/3$ in the sum of the signs.
On the resulting closed metric prove $E_++E_-<16\pi^2$ by separately
accounting for both signs on every piece. Prove that for some $\lambda>0$
the conformal-Laplacian quadratic form dominates
$\lambda\int u^2\,dV_g$ for every smooth test $u$. Record the marked
smooth identification for this construction. An isometric
double comparison requires the actual boundary isometry and estimates
on the same physical half metric as in Problem~\ref{ed04:entry}.
On the marked Miller--Schupp boundary, the regular and long-cuff full-width trials still need a positive closed sign and the two signed costs. Coverage of additional smooth types is the separate target
in Theorem~\ref{thm:ed05-coverage}.
\end{openproblem}

\begin{openproblem}[Native reduction and its first uncovered cases]
\label{prob:ed06-marked-reduction}
\uses{prop:ed45-relative-pair,prop:ed46-reverse-core}
Route 09. Formalize the assessed Investigation 05 criterion on a smooth
four-ball with a split zero-framed lower unknot $A$, an integral pair of
native attachments from $Y_1$ to the immediate $Y_3$, and the actual
rank-three terminal upper body and closing, with
$Y_r=\#_r(S^1\times S^2)$. If either single-attachment cut of the pair is
reducible, the corresponding stationary knot localizes in a ball.
Use the stationary filling and framed slides to localize both tubes,
obtain the local surgery $Y_2$, put $A$ last and cancel with the actual
upper body. The assessed conclusion is a strict same-type
$(0,3,3)\to(0,2,2)$ reduction in indices one through three, with no
stabilization or hypothesis on the initial intersection count. Carry
the forward and inverse maps of the actual closing. With a residual
continuation, account for its transport and inverse before asserting
any decrease. In a least-three presentation of this shape the two
native second cuts are irreducible, have first Betti number two and
admit no nonzero-degree map from a closed oriented positive-scalar
three-manifold. Seek a marked slide, new order or level for one such
input, or recognize one actual localized rank-two closing. Route 10
owns the physical deformation on a supplied qualifying cut. The product-doubled tangle class now has a scientific recognition with three cancellations and no births, including a control with both cuts irreducible. On the native GST outlet, up to three slides do not yield a first unstabilized cancellation in the assessed family.
\end{openproblem}

\subsection{Physical suppliers for fixed-type attainment}

\begin{proposition}[First variation on the original metric]
\label{prob:ed06-original-variation}
\uses{prop:ed50-gradient,prop:ed50-commuted}
\lean{SmoothPoincare.Bour.originalGradient_pairing, SmoothPoincare.Bour.bourCovariantGradient_apply_commuted, SmoothPoincare.Bour.bourOriginalVelocity_eq_negTwoGradient_iff_components, SmoothPoincare.Bour.originalGradient_regularSlab}
\leanok
On a smooth closed oriented four-dimensional regular metric family with
constant Euler integral, the derivative of $\mathcal F^p$ pairs the
symmetric tensor $G^p$ of Proposition~\ref{prop:ed50-commuted} with
the actual metric velocity. Its explicit component evolution on every tangent
pair is equivalent to the physical equation $K=-2G^p$. A regular
solution of the negative twice-gradient equation uses this same tensor
for dissipation and volume preservation. Local conformal Euler constancy
also gives $\operatorname{tr}G^p=-p\Delta R/4$.
\end{proposition}
\begin{proof}
\leanok
Apply the two adjoints to the physical Ricci velocity, include both
inverse-metric and volume derivatives, then commute the crossed Ricci
derivatives. Test on the actual scale, conformal and solution paths.
\end{proof}

\begin{openproblem}[Harmonic radius and uniform stationary interior estimates]
\label{prob:ed06-stationary-interior}
\uses{prop:ed51-morrey,prop:ed47-mixed,prop:ed48-trace}
Route 13. Use the genuine constant-scalar stationary equation for each
mixed return $q_j$, with $p_j=\alpha_j$ for a closed regularized return
and $p_j=0$ for an exact Yamabe return. On the original and
curvature-normalized scales $h_j=a_jq_j$, $a_j\ge1$, start with a common
ordered four-slot curvature bound and the measured volume and injectivity
estimates. Supply the whole-space scalar $L^6$ Laplacian estimate in the pinned library and construct a common interior harmonic radius with controlled
$C^{1,\gamma}$ coefficients. Use the scalar diagonal Ricci principal
part in these coordinates and the $\tfrac12\nabla^*\nabla$ principal
part of the trace-free Ricci equation to obtain coupled interior
$W^{2,b}$ estimates for $b>4$ and then estimates of every fixed order
on smaller balls. The coupled physical estimate at exponent six is Proposition~\ref{prop:ed51-coupled}. Control $s_j=R(q_j)/a_j$ and the coefficient
$(2-p_j)s_j/12$ uniformly, with no division by $p_j$.
\end{openproblem}

\begin{openproblem}[Nonflat ALE end geometry and marked return]
\label{prob:ed06-ale-return}
\uses{prop:ed51-pfaffian,prop:ed49-exterior,prop:ed47-mark,prop:ed47-ends,prop:ed48-transfer,prop:ed48-cap}
Route 13. Treat the fixed-parameter original extraction and the mixed
stationary extraction at their respective scales. Given a complete
nonflat scalar-flat Bach-flat pointed limit with finite full curvature,
positive Yamabe test bounds and the source chiral and Ricci budgets,
prove ALE decay and the limiting Gauss--Bonnet boundary terms including
ambient curvature. Exclude nontrivial link groups before applying
$D_X=2(E_+(X)+E_-(X))+32\pi^2(k-1)$ to exclude $k\ge2$ by the
strict Ricci ceiling. On one connected truncated core with all end
buffers, obtain a coherent oriented embedding, actual collar-compatible
cap data and forward and inverse smooth comparisons. Construct a
positive compactification on the same smooth type with
$I\le e_X\le\liminf_i e_i$ and $D_X=4e_X$. The separate
start-deficit test of Proposition~\ref{prop:ed48-deficit} applies to
a second marked return with the specified tail hypothesis.
\end{openproblem}

\subsection{Rigidity on the attained exact stratum}

\begin{openproblem}[Least Ricci selector at the exact attained value]
\label{prob:ed06-least-ricci}
\uses{prop:ed49-selector,thm:ed05-bour-attainment,prop:ed48-transfer}
Route 07. Fix one entered nonstandard smooth type with attained
$I=a_M$ and $T+\varepsilon_*<I\le 2T-\delta$ for some fixed $\delta>0$,
where $a_M$ is the first positive critical value and $\varepsilon_*$ is
the recognition gap. Formalize
the compactness supplier for unit-volume positive Yamabe representatives
at the exact value $a_M$, modulo oriented diffeomorphisms, with the
fixed-band marked-neck and same-type strict-drop hypotheses. Pass the
Yamabe inequality and equality for constant tests to the limit. Then
$D=4I-(Y_{\rm rd}^2-Y^2)/12$, where $Y_{\rm rd}$ is the spherical
Yamabe constant. Select the minimum of $Y^2$ and hence of $D$.
Strict spherical comparison at the maximum of $Y^2$ bounds the deficit
away from zero on this exact stratum. Relate that bound to near-minimizing
starts only under a separate recovery hypothesis. The post-attainment value derivative is computed in Proposition~\ref{prop:ed51-slope}.
\end{openproblem}

\begin{openproblem}[Polarized Weyl Hessian for a non-Einstein minimizer]
\label{prob:ed06-polarized-hessian}
\uses{prob:ed06-least-ricci,prob:ed06-original-variation,prop:ed47-variation}
Route 07. On the least-$D$ candidate with $D>0$ let
$Z=\operatorname{Ric}_0$ and
$P=(Z^\sharp\circ Z^\sharp)^\flat-\tfrac14|Z|^2g$ in its
constant-scalar Yamabe gauge. Compute the polarized physical Hessian
$H=D^2E_+$ on $Z,P$ modulo conformal and diffeomorphism directions.
Identify a geometric subclass where the span is independent modulo
gauge and the Hessian matrix has negative determinant, or where
$H(Z,Z)=0$ and $H(Z,P)\ne0$. Realize the resulting negative direction
by an admissible metric path and a strictly lower positive same-type
competitor. Record the dependent and nonnegative cases that remain.
\end{openproblem}

\begin{openproblem}[A transverse-traceless spectral window in the Einstein case]
\label{prob:ed06-tt-window}
\uses{prop:ed49-round,prob:ed06-least-ricci,prop:ed47-variation}
Route 07. For $D=0$ the selected positive Yamabe metric is Einstein.
At scalar $s>0$, the physical Weyl Hessian on transverse-traceless
tensors $h$ is $\tfrac18\langle h,(-\Delta_L-s/2)(-\Delta_L-s/3)h\rangle_{L^2}$,
where $\Delta_L$ has leading term $\nabla^i\nabla_i$ in the master
convention. For a specified Einstein subclass, force an actual smooth
transverse-traceless eigentensor of $-\Delta_L$ with eigenvalue in
$(s/3,s/2)$ and realize its admissible negative variation, or give
smooth recognition on that subclass. Endpoint eigentensors have zero
Hessian.
\end{openproblem}

\subsection{Checked results on the physical metric and the marked exterior}

\begin{proposition}[The two spatial adjoints of the connection variation]
\label{prop:ed49-adjoint}
\uses{prop:ed48-slab,prop:ed48-trace}
\lean{SmoothPoincare.Bour.bourRicciDivergence_adjoint, SmoothPoincare.Bour.bourWeightedRicciKoszul_firstAdjoint, SmoothPoincare.Bour.bourWeightedRicciKoszul_secondAdjoint, SmoothPoincare.Bour.bourKoszulFlux_native, SmoothPoincare.Bour.bourClosedTensorPairing_native}
\leanok
At a smooth regular time on a closed physical metric path, write $K=\dot g$,
$A=\operatorname{Ric}+(p-4)Rg/12$, and
$C_{kij}=\tfrac12(\nabla_i K_{jk}+\nabla_j K_{ik}-\nabla_k K_{ij})$.
Then $C$ is the lowered variation of the Levi--Civita connection.
The physical $L^2$ pairing of $\nabla^k C_{kij}$ with $A$ is
$-\langle C,\nabla A\rangle_{L^2}$, with the derivative in the first slot.
Moving the remaining derivative gives a pairing with $K$, using the inverse
slot permutations and coefficient $-1/2$ on the rough Laplacian of $A$.
\end{proposition}
\begin{proof}
\leanok
Apply covariant divergence integration by parts twice, using invariance of the
metric pairing under simultaneous index permutation. The differentiated Koszul
formula and the native smooth velocity identify the connection term on the path.
\end{proof}

\begin{proposition}[The time derivative of the original Ricci tensor in a fixed chart]
\label{prop:ed49-ricci-time}
\uses{prop:ed48-slab}
\lean{SmoothPoincare.Bour.bourChartChristoffel_spatial_hasDerivAt, SmoothPoincare.Bour.bourChartRicci_hasDerivAt, SmoothPoincare.Bour.bourMetricRicciAt_chartBasis_hasDerivAt}
\leanok
On a smooth regular metric family in a fixed chart, time differentiation
commutes with each spatial derivative of the actual Christoffel coefficients.
The derivative of each physical Ricci component on the fixed chart basis
is obtained by taking the two spatial derivatives of the Christoffel speed
in their original order and differentiating both Christoffel products.
\end{proposition}
\begin{proof}
\leanok
Use joint smoothness in time and space on the chart interior to exchange
the derivatives. Differentiate the two contracted products by the product
rule. The chart-basis curvature identity identifies the result with the
corresponding component of the original Ricci tensor.
\end{proof}

\begin{proposition}[Connection estimates in a supplied harmonic chart]
\label{prop:ed49-interior}
\uses{prop:ed48-harmonic,prop:ed48-normal}
\lean{SmoothPoincare.bourChartConnectionDerivative_bound, SmoothPoincare.bourChartConnection_trace_eq_zero, SmoothPoincare.bourChartRoughLaplacian_eq_secondCovariantTrace}
\leanok
In the four-dimensional metric chart with inverse entries bounded by $A$,
first derivatives by $L$ and second derivatives by $D$,
each first derivative of a Christoffel coefficient is bounded by
$96A^2 L^2+6AD$. At a harmonic point the traced coordinate-frame
connection vanishes. The raw tensor connection Laplacian on a smooth
compactly supported tensor reduces there to the metric trace of the
two iterated coordinate covariant derivatives.
\end{proposition}
\begin{proof}
\leanok
Differentiate the inverse-metric identity and the Christoffel formula.
Contract the Laplacians of the coordinate functions, then use the
linearity of the tensor connection in its directional input.
\end{proof}

\begin{proposition}[Euler transgression on a supplied nonflat end]
\label{prop:ed49-end}
\uses{prop:ed48-cap,prop:ed48-transfer}
\lean{SmoothPoincare.bourEulerBoundaryDensity_withArea_error, SmoothPoincare.bourALELevelBoundary_tendsto, SmoothPoincare.bourExhaustedEulerIdentity, SmoothPoincare.bourOneEuclideanEnd_of_boundaryLimits}
\leanok
Use the outward shape and the curvature matrix on the three
complementary oriented tangent two-planes of a normalized link.
If the shape, curvature and area ratio are within $\varepsilon\le1$
entrywise of $\mathrm{Id},0,1$, the area-weighted boundary density
differs from $4$ by at most $412\varepsilon$.
With normalized error tending to zero, integrability and the degree-$m$
round covering measure, its integral tends to $8\pi^2/m$.
For a supplied scalar-flat compact exhaustion with the finite Euler
formula, spherical links, $\chi=2-k$, and $0\le D<32\pi^2$, the
nonempty set of ends has cardinality one.
\end{proposition}
\begin{proof}
\leanok
Expand the cubic determinant and bound the mixed shape--curvature contraction.
Integrate the error on the fixed quotient link and pass all $k$ boundary
terms through the limit. The resulting identity is
$D=2(E_++E_-)+32\pi^2(k-1)$.
\end{proof}

\begin{proposition}[Buffered disk extension on a recognized sphere]
\label{prop:ed49-exterior}
\uses{prop:ed48-cap}
\lean{SmoothPoincare.bourRecognizedCap_commonChart, SmoothPoincare.bourMarkedExterior_of_recognizedSphere}
\leanok
Let the capped source exterior and target double be orientedly recognized
as standard four-spheres. Suppose their cap coordinates extend over a closed
radius-$R$ disk, the seam has radius $0<r<R$, and the two parametrizations
carry the same center orientation. If their composition gives the prescribed
map on the buffered disk, an oriented global diffeomorphism agrees with it
and with its inverse on the two open buffers and identifies the complements
of the open radius-$r$ disks.
\end{proposition}
\begin{proof}
\leanok
Choose an exterior point for each full buffer. On the standard sphere,
move the first point to the second by an oriented map and align the cap
centers by a further motion fixing that exterior point. Stereographic
projection now gives a common chart on the entire buffers. At their common
center the transported frames have equal orientation. Apply the supported
cap extension in this chart and pull it back to the given double.
\end{proof}

\begin{proposition}[Least Ricci energy from first-value compactness]
\label{prop:ed49-selector}
\uses{prop:ed48-transfer}
\lean{SmoothPoincare.attainedYamabe_limit_of_metricConvergence, SmoothPoincare.exists_attainedYamabe_leastRicci}
\leanok
Supply positive subcritical attainment at $0<I\le8\pi^2-\delta$ on the
fixed type, conformal normalization and Bach covariance, exact first-value
compactness modulo oriented pullback, the closed Euler value and the sharp
spherical comparison with equality implying zero chiral energy.
Then $I=a_M$. One unit-volume exact Yamabe representative minimizes
both $Y^2$ and $D$. Another attains the maximum $Y^2$ and its strictly
positive squared spherical deficit bounds the deficit on the exact stratum.
\end{proposition}
\begin{proof}
\leanok
Pass each fixed smooth conformal test through smooth metric convergence and
use the constant test for equality. Apply compactness to a minimizing
and a maximizing sequence separately and use the exact closed Euler identity.
\end{proof}

\begin{proposition}[The round transverse-traceless window and a conditional descent]
\label{prop:ed49-round}
\uses{prop:ed47-variation}
\lean{SmoothPoincare.einsteinWeylTTFactor_neg_iff, SmoothPoincare.standardRoundFour_curvatureAction_traceFree, SmoothPoincare.standardRoundFour_negLichnerowicz_eigenvalue_ge_eight, SmoothPoincare.standardRoundFour_TTMode_weylFactor_ge_one, SmoothPoincare.bourStrictTTWindowExcludesAttainment}
\leanok
On the unit round sphere the trace-free operator $8\mathrm{Id}-\operatorname{tr}\nabla^2$
has eigenvalues $\lambda\ge8$ on nonzero smooth modes. The coefficient
$(s/2-\lambda)(s/3-\lambda)/8$ at $s=12$ is at least one.
On any actual admissible path realizing this signed Hessian on a nonzero
tensor, a strict TT mode $s/3<\lambda<s/2$ excludes attainment.
\end{proposition}
\begin{proof}
\leanok
Compute the round curvature action. Pair the shifted rough operator with
the eigentensor and use the covariant Green identity. Factor the signed
coefficient and use the physical two-sided variation test at a strict sign.
\end{proof}

\begin{proposition}[Regularized Ricci margin from an attained comparator]
\label{prop:ed49-margin}
\uses{prop:ed47-mixed}
\lean{SmoothPoincare.Bour.bourOriginal_attainedComparator_margin, SmoothPoincare.Bour.bourNoMarkedBubble_of_ricciMargin}
\leanok
Fix $0<p<1$ and a comparator of energy $I$ on the original type.
On a supplied regularized trajectory suppose
$F^p(g_0)\le F^p(g_*)+\epsilon<(1-p)16\pi^2$ and
$4\epsilon\le p(4I-D_*)$. Every original smooth time has positive Yamabe
sign, $D(t)\le(4I+D_*)/2$, and $S(t)\le S_*+24\epsilon/p$.
Any sequence of original-type metrics with eventual $D\le d<4I$ excludes
a positive marked return satisfying $4e_X\le\liminf D$.
\end{proposition}
\begin{proof}
\leanok
The strict pinching margin keeps the full trajectory in the positive class.
Insert $e(t)\ge I$ in the two physical formulas for $F^p$.
A marked positive metric on the original type has $e_X\ge I$.
\end{proof}

\subsection{The physical gradient and both extractions}

\begin{proposition}[Ricci velocity on the original path]
\label{prop:ed50-palatini}
\uses{prop:ed49-adjoint,prop:ed49-ricci-time}
\lean{SmoothPoincare.Bour.bourMetricRicciAt_chartKoszul_hasDerivAt, SmoothPoincare.Bour.bourMetricRicciAt_hasDerivAt, SmoothPoincare.Bour.bourKoszulFlux_native_trace, SmoothPoincare.Bour.bourKoszulTraceForm_native}
\leanok
Let $K=\dot g$ on an original regular slice and
$C_{aik}=(\nabla_iK_{ka}+\nabla_kK_{ia}-\nabla_aK_{ik})/2$.
In the fixed chart axes and intrinsically,
$\dot{\mathrm{Ric}}_{ik}=\nabla^a C_{aik}-\nabla_k\tau_i$
with $\tau=\tfrac12 d\operatorname{tr}_gK$.
\end{proposition}
\begin{proof}
\leanok
Differentiate the chart connection and lower its output in the same chart.
Metric compatibility combines the differentiated inverse and connection
terms into the covariant derivative of the lowered flux. Trace output
and last slot to obtain $\tau$, then use a basis to recover the tensor.
\end{proof}

\begin{proposition}[A symmetric gradient and its regular time slab]
\label{prop:ed50-gradient}
\uses{prop:ed50-palatini,prop:ed49-adjoint,prop:ed48-slab,prop:ed48-trace}
\lean{SmoothPoincare.Bour.bourEnergy_hasDerivAt_contractions, SmoothPoincare.Bour.bourWeightedPalatini_secondAdjoint, SmoothPoincare.Bour.originalGradient_pairing, SmoothPoincare.Bour.bourCovariantGradient_trace, SmoothPoincare.Bour.originalGradient_regularSlab}
\leanok
Write $b=(p-4)/12$, $A=\mathrm{Ric}+bRg$ and
$B=(|\mathrm{Ric}|^2+bR^2)/2$.
On a closed regular signed metric family of constant Euler integral,
the covariant $L^2$ gradient is the symmetric part of
\[
 \nabla^a\nabla_i A_{aj}-\tfrac12\nabla^a\nabla_a A_{ij}
 -\tfrac12(\operatorname{div}\operatorname{div}A)g_{ij}
 -(\mathrm{Ric}\circ\mathrm{Ric})_{ij}-bR\,\mathrm{Ric}_{ij}+\tfrac12 B g_{ij}.
\]
Its integrated trace vanishes. Under local constancy of the Euler integral
on conformal tests its trace at each point is $-p\,\Delta R/4$ in the
convention $\Delta=\operatorname{tr}\nabla^2$. On any regular solution
with $\dot g=-2G^p$, volume is constant, and the integrated squared
gradient is half the energy drop.
\end{proposition}
\begin{proof}
\leanok
Differentiate the two inverse metrics in the Ricci square and the
inverse metric in the scalar trace. Include the volume derivative in
the same integral. Transfer both derivatives of each Palatini term.
The symmetric projection pairs with the original velocity. Use actual
scale and conformal paths for the trace tests and integrate along the given solution.
\end{proof}

\begin{proposition}[Physical capture by a common harmonic coordinate ball]
\label{prop:ed50-chart}
\uses{prop:ed48-normal,prop:ed48-harmonic,prop:ed49-interior}
\lean{SmoothPoincare.bourUniformHarmonicBalls_of_yamabe, SmoothPoincare.BourHarmonicChart.exists_buffered, SmoothPoincare.BourHarmonicChart.transport_to_original, SmoothPoincare.BourHarmonicChart.exists_model_metric, SmoothPoincare.bourUniformHarmonicMetricTrace_bound}
\leanok
Under the quantitative Peters harmonic-radius input at a common positive
injectivity radius and ordered curvature ceiling, the normalized metrics
$h=ag$ on the stationary sequence have one radius $r$ and one ellipticity
constant $L$. The chart is defined on $B(0,4r)$ and its closed radius-$r$
image contains the physical $h$-ball of radius $r/L$. The same set
contains the $g$-ball of radius $r/(L\sqrt a)$. An ambient model metric
agrees with the actual pullback on $B(0,3r)$, where the model's fixed
coordinates are harmonic. Its metric entries there satisfy a uniform
pointwise bound on $|h^{ab}\partial_a\partial_b h_{ik}|$ from the actual
ordered curvature ceiling, the first jets and ellipticity.
\end{proposition}
\begin{proof}
\leanok
Apply the physical noncollapse and injectivity estimates before choosing the
Peters chart. Two length comparisons and first exit give the nested balls.
Extend the pullback on the compact inner ball and use Laplacian naturality.
Pull Ricci back on the same ball and insert the curvature and first-jet
bounds in the harmonic metric equation.
\end{proof}

\begin{proposition}[Measured finite Euler correction on an end link]
\label{prop:ed50-link}
\uses{prop:ed49-end}
\lean{SmoothPoincare.bourUnitEuclideanLevel_shape_error, SmoothPoincare.bourImmersedLinkAreaDensity_error, SmoothPoincare.bourImmersedLinkEulerIntegral_error, SmoothPoincare.bourEulerBoundaryDensity_eq_integral_connectionPath}
\leanok
On the normalized unit sphere use $u=|x|^2/2$ and the outward normal.
Under the stated bilinear metric, connection difference, and ordered
ambient curvature bounds of size $\epsilon$, the shape error is at most
$20\epsilon$ and the induced area ratio error at most $3\epsilon$.
On a compact immersed link with the reference unit shape and
outward-adapted orthonormal frames, the physical boundary integral differs
from four times the reference area by at most
$412(20\epsilon)$ times that area when $20\epsilon\le1$.
The scalar connection-path integral has value
$2\sum S_{ab}C_{ab}+4\det S$.
\end{proposition}
\begin{proof}
\leanok
Compare the two Hessians of $u$, bound the length of its gradient and
compare the two induced volume measures. The norm of the ordered curvature
controls each complementary component. Integrate the E49 density estimate
on this fixed link. The tangential quadratic term in the connection
interpolation has weight $1-t^2$.
\end{proof}

\begin{proposition}[The same post-attainment Ricci margin at both extractions]
\label{prop:ed50-second}
\uses{prop:ed49-selector,prop:ed49-margin,prop:ed48-transfer,prop:ed47-return}
\lean{SmoothPoincare.Bour.exists_bourLeastRicci_uniformRegularizedStarts, SmoothPoincare.Bour.bourFirstClosedExtraction_bounds, SmoothPoincare.Bour.bourSecondMarkedExtraction_of_canonicalRescaling_impossible, SmoothPoincare.Bour.bourRegularizedMinimum_of_secondClosed, SmoothPoincare.Bour.bourRegularizedMinima_limit_leastRicci}
\leanok
Given scientific attainment, first-value compactness, the Yamabe and
sharp spherical comparison inputs on a type with $0<I\le8\pi^2-\delta$,
choose one exact least-$D$ metric $g_*$. The same
$d=(4I+D_*)/2<4I$ works for arbitrarily accurate near-minimizers at
every sufficiently small positive $p$. At each fixed $p$, physical closed
endpoints inherit the bound. A second marked positive return on the
original type with the one-end Ricci identity contradicts it.
A supplied whole closed second limit with the stated test and integral
convergences attains the regularized infimum. Along actual global
minima as $p\downarrow0$, a supplied positive whole Yamabe and Bach
limit attains the least $D$ and the chiral excess is $o(p)$.
\end{proposition}
\begin{proof}
\leanok
Compare both physical formulas for the energy with the same exact metric.
Transfer the bound in the original gauges along each first closed extraction.
For a second pointed limit the ordinary Ricci lower bound gives
$4I\le4e_X=D_X\le d$. For whole convergence use both sides of the
variational infimum bound. Finally compare each global minimizer with $g_*$
and pass its Ricci integral to the whole limit.
\end{proof}

\begin{proposition}[Vanishing dissipation on terminal backward windows]
\label{prop:ed50-terminal}
\uses{prop:ed50-gradient,prop:ed47-motion}
\lean{SmoothPoincare.Bour.originalGradient_energy_tendsto_terminal, SmoothPoincare.Bour.originalGradient_dissipation_tendsto_terminal, SmoothPoincare.Bour.originalGradient_parabolicDissipation_tendsto_terminal}
\leanok
For $0\le p\le1$, assume the original equation, Euler constancy, and an
order-connected regular time interval. The energy tends to its infimum
at the right end of this interval. If both endpoints of original slabs
are cofinal there, the dissipation on those slabs tends to zero. The same
holds on the corresponding fixed backward parabolic windows.
\end{proposition}
\begin{proof}
\leanok
The energy is antitone and the chiral--Ricci formula bounds it below by zero.
Take the difference of the two cofinal energy limits and apply the exact
rescaling identity to each window.
\end{proof}


\begin{proposition}[Curvature-commuted gradient on a general slice]
\label{prop:ed50-commuted}
\uses{prop:ed50-gradient}
\lean{SmoothPoincare.Bour.bourRicciWeight_divergence, SmoothPoincare.Bour.bourRicciWeight_doubleDivergence, SmoothPoincare.Bour.bourCrossedRicci_eq_ricciSquare_sub_curvature, SmoothPoincare.Bour.bourRicciFluxAdjoint_apply_commuted, SmoothPoincare.Bour.bourCovariantGradient_apply_commuted, SmoothPoincare.Bour.bourCovariantGradient_apply_csc}
\leanok
For $b=(p-4)/12$, $B=(|\mathrm{Ric}|^2+bR^2)/2$ and
$\Delta=\operatorname{tr}\nabla^2$, the same physical tensor satisfies
\[
\begin{aligned}
G^p={}&-\tfrac12\Delta\mathrm{Ric}-\mathrm{Rm}^{\circ}(\mathrm{Ric})
       +\tfrac{p+2}{12}\nabla^2R\\
     &-\tfrac{p-1}{12}(\Delta R)g-bR\,\mathrm{Ric}+\tfrac12Bg .
\end{aligned}
\]
For constant scalar curvature the two scalar derivative terms vanish.
\end{proposition}
\begin{proof}
\leanok
The contracted Bianchi identity gives
$\operatorname{div}(\mathrm{Ric}+bRg)=(p+2)dR/12$.
In an orthonormal basis the crossed curvature action is
$\mathrm{Ric}^2-\mathrm{Rm}^{\circ}(\mathrm{Ric})$.
Its first term cancels the moving-inverse correction. Use metric
compatibility to differentiate $Rg$ twice and trace its Hessian.
\end{proof}

\subsection{Harmonic interiors and terminal spacetime limits}

\begin{proposition}[Metric estimate on the physical pullback]
\label{prop:ed51-metric}
\uses{prop:ed50-chart}
\lean{SmoothPoincare.Bour.bourFlatHessianEstimate_metricFrozen, SmoothPoincare.bourUniformHarmonicMetricW2_of_flatLaplacian}
\leanok
Assume the scalar whole-space Laplacian estimate
$\|D^2v\|_{L^6}\le C_\Delta\|\Delta v\|_{L^6}$ for compactly supported smooth
$v$ and the model Haar measure. On the realized harmonic ball of
Proposition~\ref{prop:ed50-chart}, the metric entries have a common
$W^{2,6}$ bound on $B(0,5r/2)$. The constant precedes the choice of
metric, base point and realization. The statement applies to $a=1$
and $a=Q$ under the corresponding normalized ordered curvature ceiling.
\end{proposition}
\begin{proof}
\leanok
Conjugate the frozen inverse Gram operator to the Laplacian and cancel the constant
Haar Jacobian. Bound its inverse transformation by the metric upper form.
Use ellipticity and the first metric jets to choose a common freezing radius.
Absorb the oscillation term, differentiate a compact cutoff, and take a fixed finite cover.
\end{proof}

\begin{proposition}[Principal forcing for the trace-free Ricci tensor]
\label{prop:ed51-principal}
\uses{prop:ed50-chart,prop:ed50-commuted}
\lean{SmoothPoincare.bourStationary_roughTraceFreeRicci_eq, SmoothPoincare.bourChartContractedChristoffelDerivative_eq_of_laplacian, SmoothPoincare.bourChartTwo_principal_bound}
\leanok
On a constant-scalar stationary four-metric the rough trace of $Z=\mathrm{Ric}_0$
is algebraic in curvature. In harmonic coordinates on an open neighborhood,
the contracted connection derivative in the ordered Hessian of $Z$ is
controlled by raised curvature, inverse metric first derivatives, and
Christoffel symbols. The scalar principal forcing for each entry of $Z$
is bounded by a zeroth-order constant plus a multiple of the sum of its
first coordinate derivatives in absolute value. For $h=ag$ the scalar
coefficient is $(2-p)R(g)/(12a)$.
\end{proposition}
\begin{proof}
\leanok
Differentiate $h^{ab}\Gamma^k_{ab}=0$, and use the coordinate curvature
identity to switch the two derivative slots. Expand both covariant
derivatives on a two-tensor and insert the resulting identity in each
lower slot. Constant rescaling preserves the covariant tensor $Z$.
\end{proof}

\begin{proposition}[Interior absorption at exponent six]
\label{prop:ed51-cutoff}
\uses{prop:ed51-metric,prop:ed51-principal}
\lean{SmoothPoincare.Bour.sixthPowerDirectional_cutoff_lpNorm, SmoothPoincare.bourCutoffSystem_real_norm_bounds, SmoothPoincare.bourCutoffSystem_sixthInterior}
\leanok
On a single freezing patch suppose the stationary coordinate equations satisfy
$|Pu_j|\le F_0+F_1\sum_{i,k}|\partial_i u_k|$. A compact cutoff $\eta$ gives
bounds on $\|D^2(\eta^2u_j)\|_6$ and $\sum_{i,j}\|\eta\partial_i u_j\|_6$.
Their constants depend on the frozen estimate, coefficient and cutoff
bounds, patch volume and $\sup |u_j|$. The coupled system in four dimensions
has sixteen ordered tensor entries and sixty-four first derivatives.
\end{proposition}
\begin{proof}
\leanok
Integrate the sixth power of the weighted directional derivative by parts.
Young's inequality leaves an arbitrarily small multiple of the matching
$\eta^2\partial_i^2u_j$ norm. The squared-cutoff product rule bounds
these diagonal Hessians by the localized Hessians and weighted first
derivatives. Sum all directions and choose the interpolation length so
the total coefficient is at most one half.
\end{proof}

\begin{proposition}[First hits and stationarity through the final slice]
\label{prop:ed51-terminal}
\uses{prop:ed50-terminal,prop:ed50-commuted}
\lean{SmoothPoincare.Bour.originalGradient_firstHit_nativeWindows, SmoothPoincare.Bour.bourParabolicVelocity_hasDerivAt, SmoothPoincare.Bour.originalGradient_stationary_of_terminalComparison}
\leanok
Fix one parameter and a regular original flow with unbounded ordered curvature
after a regular anchor. There are cofinal first hits $(t_n,x_n)$ for any
diverging levels $Q_n$ above the anchor maximum. Each fixed normalized
backward window eventually has curvature at most one everywhere, value
one at the final base point, and vanishing integrated native gradient square.
Under time-independent compact comparison maps, a metric lower comparison,
and uniform gradient-norm comparison with a continuous limiting field, the
limiting gradient vanishes through the final slice. This transfer works
on an exhausted pointed carrier as well as a whole closed limit.
\end{proposition}
\begin{proof}
\leanok
Minimize the time projection on each compact spacetime crossing set.
The rescaled gradient has weight minus one. A point of positive limiting
gradient norm gives a compact spatial patch of positive measure and a
time interval with uniformly positive dissipation on its actual image.
Continuity also detects positivity at the right endpoint.
\end{proof}

\begin{proposition}[Outward connection forms on a finite end link]
\label{prop:ed51-pfaffian}
\uses{prop:ed50-link,prop:ed49-end}
\lean{SmoothPoincare.bourBoundaryPathTangentialCurvature_of_connectionForms, SmoothPoincare.bourImmersionNormalDefect_eq_neg_shape, SmoothPoincare.bourFrameEulerTransgression_boundary_integral, SmoothPoincare.bourEulerBulk_setIntegral_eq_scalarFlat, SmoothPoincare.bourRoundEulerShell_total}
\leanok
In an outward-adapted oriented boundary frame let $S$ be the shape
matrix and $C$ the ambient tangential curvature on complementary
wedges. The affine connection from the split to the ambient connection
gives the integrated Pfaffian coefficient $2\langle S,C\rangle+4\det S$.
The scalar-flat bulk on a finite truncation is
$E_+(K)+E_-(K)-D(K)/2$ with the physical restricted volume.
For the supplied round polar shell $0<a<b<\pi$, its two outward
boundary integrals and its bulk sum to zero.
\end{proposition}
\begin{proof}
\leanok
Metric compatibility fixes the signs of the two mixed blocks. The curvature
along the path has tangential component $C+(1-t^2)S\wedge S$ in the
complementary-wedge convention. Alternate, contract with $S$, and integrate.
Integrate the scalar-flat density on the same core.
\end{proof}

\begin{proposition}[Right slope of the positive-class value]
\label{prop:ed51-slope}
\uses{prop:ed50-second}
\lean{SmoothPoincare.Bour.bourRegularized_secant_antitone, SmoothPoincare.Bour.bourRegularizedValue_hasRightDerivative}
\leanok
After the fixed-parameter extractions, assume one sequence of actual global
regularized minima with a whole positive unit-volume Yamabe and Bach
limit and convergence of the physical chiral and ordinary Ricci integrals.
For one exact least-$D$ comparator and $J(p)=\inf_g(2(1-p)e(g)+pD(g)/2)$
on the positive metrics, $J(0)=2I$ and $J'_+(0)=D_*/2-2I$.
\end{proposition}
\begin{proof}
\leanok
Since $e(g)\ge I$, the positive right secant of each affine functional
at height $2I$ decreases with the parameter. Its infimum has the same
monotonicity. Test the upper bound on the single least-$D$ metric and
the lower bound along the returned sequence, where $D$ converges to $D_*$.
\end{proof}

\begin{proposition}[Coupled estimates on both stationary scales]
\label{prop:ed51-coupled}
\uses{prop:ed50-chart,prop:ed51-cutoff,prop:ed51-principal}
\lean{SmoothPoincare.bourStationaryRoughTraceFreeRicci_of_partialDiffeomorph_inner, SmoothPoincare.Bour.bourUniformStationaryHarmonicTraceFreeRicci_bounds, SmoothPoincare.bourStationary_harmonic_W26, SmoothPoincare.bourMixedStationary_uniformHarmonicW26}
\leanok
For a common positive Yamabe floor, scalar ceiling and normalized ordered
curvature bound, assume the Peters chart input and the scalar Euclidean
$L^6$ Laplacian estimate. One choice of radius, chart controls and finite bound
works for all stationary constant-scalar metrics with $0\le p\le1$ at both
$a=1$ and the curvature scale $a=Q\ge1$. On the same physical realization,
the sum of the $W^{2,6}$ norms of a metric component and the corresponding
covariant trace-free Ricci component is bounded on $B(0,5r/2)$.
The closed coordinate ball of radius $r$ captures the original metric
ball of radius $r/(L\sqrt a)$.
\end{proposition}
\begin{proof}
\leanok
Transport the ordered curvature and rough trace through the local isometry on
$B(0,3r)$. Bound the stationary reaction on the original carrier in an
orthonormal frame. The squared cutoff and a common oscillation radius give
uniform bounds on each smaller patch. A fixed finite cover and reconstruction
of the full first derivative give all three norms. Metric entries of any two
extensions and their jets agree on the common open physical ball.
\end{proof}

\begin{proposition}[The fixed-parameter pointed scalar vanishes]
\label{prop:ed51-scalar}
\uses{prop:ed51-terminal,prop:ed50-commuted,prop:ed48-normal}
\lean{SmoothPoincare.Bour.terminalMetric_eq_final_of_nativeSpeed, SmoothPoincare.Bour.bourCovariantGradient_trace_eq_scalarLaplacian, SmoothPoincare.bourSmoothPointedBallCutoff, SmoothPoincare.Bour.bourHarmonic_cutoff_energy, SmoothPoincare.Bour.bourPointedStationary_scalarFlat, SmoothPoincare.Bour.bourScalarFlat_stationary_zeroParameter}
\leanok
On the fixed-parameter terminal window, convergence of the metric values
on the same time-independent maps, an upper metric comparison and uniformly
vanishing native speed give a time-independent limit. A connected noncompact
boundaryless stationary four-metric with $p>0$, compact closed physical balls,
a positive unit-ball volume floor and finite $\int R^2\,dV$ has $R=0$.
At $R=0$ the same metric satisfies $G^0=0$ and
$\Delta\mathrm{Ric}+2\,\mathrm{Rm}^{\circ}(\mathrm{Ric})
=|\mathrm{Ric}|^2g/2$. On a connected closed stationary limit with
$p>0$ the scalar is constant.
\end{proposition}
\begin{proof}
\leanok
Compare two times by evaluating the actual velocity on the two transported
tangent vectors. Trace the curvature-commuted gradient to get
$-p\Delta R/4$. On the pointed carrier test harmonicity against $\eta^2 R$.
The compactly supported flux bounds $\int\eta^2|\nabla R|^2$ by
$4\sup|\nabla\eta|^2\int R^2$. Smooth distance cutoffs at increasing
radii and Fatou give $\nabla R=0$. Disjoint unit balls force infinite
volume, so the square-integrable constant vanishes.
\end{proof}


\begin{proposition}[Interior control of the first derivatives]
\label{prop:ed51-morrey}
\uses{prop:ed51-coupled}
\lean{SmoothPoincare.bourW26Norm_volume_eq, SmoothPoincare.bourHarmonicMorrey_of_W26, SmoothPoincare.bourMixedStationary_uniformHarmonicC1Third}
\leanok
In the four-dimensional Euclidean model, on $B(0,2r)$ each metric entry
and each covariant trace-free Ricci entry of the same physical realization
has uniformly bounded value and first derivative. Both first derivatives
are uniformly H\"older continuous with exponent $1/3$. The point, metric,
scale and regularization parameter are selected after the common bound.
\end{proposition}
\begin{proof}
\leanok
Normalize the additive Haar measure and localize the $W^{2,6}$ bound on
$B(0,5r/2)$ using a smooth cutoff equal to one on the smaller ball.
Apply the Euclidean Morrey inequalities to the entry and its first
derivative, with $1-4/6=1/3$.
\end{proof}

